% concatenated from example_paper.tex
%%%%%%%% ICML 2026 EXAMPLE LATEX SUBMISSION FILE %%%%%%%%%%%%%%%%%

\documentclass{article}

% Recommended, but optional, packages for figures and better typesetting:
\usepackage{microtype}
\usepackage{graphicx}
\usepackage{subcaption}
\usepackage{booktabs} % for professional tables
\usepackage{xurl}
% hyperref makes hyperlinks in the resulting PDF.
% If your build breaks (sometimes temporarily if a hyperlink spans a page)
% please comment out the following usepackage line and replace
% \usepackage{icml2026} with \usepackage[nohyperref]{icml2026} above.
\usepackage{hyperref}

\usepackage{fontawesome5}

\hypersetup{
    colorlinks=true,
    urlcolor=blue,
    linkcolor=blue
}

% Attempt to make hyperref and algorithmic work together better:
\newcommand{\theHalgorithm}{\arabic{algorithm}}

% Use the following line for the initial blind version submitted for review:
% \usepackage{icml2026}
% 
% For preprint, use
% \usepackage[preprint]{icml2026}

% If accepted, instead use the following line for the camera-ready submission:
\usepackage[accepted]{icml2026}

\usepackage{hyperref}
\usepackage{url}
\usepackage{graphicx}
\usepackage{enumitem}
\usepackage[table]{xcolor}
\usepackage{wrapfig}
\usepackage{amsmath}
% \usepackage{subfigure}
\usepackage{caption}
\usepackage{subcaption}
\usepackage{titletoc}
\makeatletter
\renewcommand{\addcontentsline}[3]{%
  \addtocontents{#1}{\protect\contentsline{#2}{#3}{\thepage}{}}}
\makeatother
\usepackage{hyperref}
\usepackage{tcolorbox}
\usepackage{tabularx}
\usepackage{pifont}
\usepackage{algorithm}
\usepackage{algorithmic}

% \usepackage{algpseudoscode}
\usepackage{makecell}
\colorlet{ours}{blue!10}
\usepackage[utf8]{inputenc} % allow utf-8 input
\usepackage[T1]{fontenc}    % use 8-bit T1 fonts
\usepackage{hyperref}       % hyperlinks
\usepackage{url}            % simple URL typesetting
\usepackage{booktabs}       % professional-quality tables
\usepackage{amsfonts}       % blackboard math symbols
\usepackage{nicefrac}       % compact symbols for 1/2, etc.
\usepackage{microtype}      % microtypography
\usepackage{xcolor}         % colors
\usepackage{amssymb}
\usepackage{mathtools}
\usepackage{amsthm}
\usepackage{multirow}
\usepackage{amsmath}
\usepackage{appendix}
\newcommand{\fix}{\marginpar{FIX}}
\newcommand{\new}{\marginpar{NEW}}
% \newtheorem{theorem}{Theorem}[section]
% \newtheorem{proposition}[theorem]{Proposition}
% \newtheorem{lemma}[theorem]{Lemma}
% \newtheorem{corollary}[theorem]{Corollary}
% \theoremstyle{definition}
% \newtheorem{definition}[theorem]{Definition}
% \newtheorem{assumption}[theorem]{Assumption}
% \theoremstyle{remark}
% \newtheorem{remark}[theorem]{Remark}
\newcommand{\appendixtoc}{%
    \startcontents[appendix]
    \printcontents[appendix]{l}{1}{\setcounter{tocdepth}{2}}
}
\newcommand{\com}[1]{\tiny$\pm$#1}

% \usepackage{etoc}

% if you use cleveref..
\usepackage[capitalize,noabbrev]{cleveref}

%%%%%%%%%%%%%%%%%%%%%%%%%%%%%%%%
% THEOREMS
%%%%%%%%%%%%%%%%%%%%%%%%%%%%%%%%
\theoremstyle{plain}
\newtheorem{theorem}{Theorem}[section]
\newtheorem{proposition}[theorem]{Proposition}
\newtheorem{lemma}[theorem]{Lemma}
\newtheorem{corollary}[theorem]{Corollary}
\theoremstyle{definition}
\newtheorem{definition}[theorem]{Definition}
\newtheorem{assumption}[theorem]{Assumption}
\theoremstyle{remark}
\newtheorem{remark}[theorem]{Remark}

% Todonotes is useful during development; simply uncomment the next line
%    and comment out the line below the next line to turn off comments
%\usepackage[disable,textsize=tiny]{todonotes}
\usepackage[textsize=tiny]{todonotes}

% The \icmltitle you define below is probably too long as a header.
% Therefore, a short form for the running title is supplied here:
\icmltitlerunning{Black-Box Detection of LLM-Generated Text Using Generalized Jensen-Shannon Divergence}

\begin{document}

\twocolumn[
  \icmltitle{Black-Box Detection of LLM-Generated Text Using \\ Generalized Jensen Shannon Divergence}

  % It is OKAY to include author information, even for blind submissions: the
  % style file will automatically remove it for you unless you've provided
  % the [accepted] option to the icml2026 package.

  % List of affiliations: The first argument should be a (short) identifier you
  % will use later to specify author affiliations Academic affiliations
  % should list Department, University, City, Region, Country Industry
  % affiliations should list Company, City, Region, Country

  % You can specify symbols, otherwise they are numbered in order. Ideally, you
  % should not use this facility. Affiliations will be numbered in order of
  % appearance and this is the preferred way.
  \icmlsetsymbol{equal}{*}

  \begin{icmlauthorlist}
    \icmlauthor{Shuangyi Chen}{yyy}
    \icmlauthor{Ashish Khisti}{yyy}

    %\icmlauthor{}{sch}
    %\icmlauthor{}{sch}
  \end{icmlauthorlist}

  \icmlaffiliation{yyy}{Department of Electrical and Computer Engineering, University of Toronto, Toronto, Canada}
  % \icmlaffiliation{comp}{Company Name, Location, Country}
  % \icmlaffiliation{sch}{School of ZZZ, Institute of WWW, Location, Country}

  \icmlcorrespondingauthor{Shuangyi Chen}{shuangyi.chen@mail.utoronto.ca}
  % \icmlcorrespondingauthor{Firstname2 Lastname2}{first2.last2@www.uk}

  % You may provide any keywords that you find helpful for describing your
  % paper; these are used to populate the "keywords" metadata in the PDF but
  % will not be shown in the document
  \icmlkeywords{Machine Learning, ICML}

  \vskip 0.3in
]

% this must go after the closing bracket ] following \twocolumn[ ...

% This command actually creates the footnote in the first column listing the
% affiliations and the copyright notice. The command takes one argument, which
% is text to display at the start of the footnote. The \icmlEqualContribution
% command is standard text for equal contribution. Remove it (just {}) if you
% do not need this facility.

% Use ONE of the following lines. DO NOT remove the command.
% If you have no special notice, KEEP empty braces:
\printAffiliationsAndNotice{}  % no special notice (required even if empty)
% Or, if applicable, use the standard equal contribution text:
% \printAffiliationsAndNotice{\icmlEqualContribution}

\begin{abstract}
  We study black-box detection of machine-generated text under practical constraints: the scoring model (proxy LM) may mismatch the unknown source model, and per-input contrastive generation is costly. We propose SurpMark, a reference-based detector that summarizes a passage by the dynamics of its token surprisals. SurpMark discretizes surprisals into interpretable states, estimates a state-transition matrix for the test text, and scores it via a generalized Jensen–Shannon (GJS) gap between the test transitions and two fixed references (human vs. machine) built once from existing corpora. Theoretically, we derive design guidance for how the discretization bins should scale with data and provide a principled justification for our test statistic. Empirically, across multiple datasets, source models, and scenarios, SurpMark consistently matches or surpasses baselines, demonstrating strong robustness across domains and generators; our experiments on hyperparameter sensitivity exhibit trends that our theoretical results help to explain.
\par\smallskip
\noindent\faIcon{github}\ 
\href{https://github.com/shuangyichen/SurpMark}{\texttt{github.com/shuangyichen/SurpMark}}
\end{abstract}

\section{Introduction}
% \vspace{-1mm}
Rapid advancements in LLMs have driven their text generation capabilities to near-human levels. 
This has blurred the boundary between human-written and machine-generated text, posing multiple concerns. These include susceptibility to fabrications \cite{Ji_2023} and outdated or misleading information, which can spread misinformation, or facilitate plagiarism \cite{10.1145/3543507.3583199}. LLMs are also vulnerable to malicious use in disinformation dissemination \cite{lin2022truthfulqameasuringmodelsmimic}, fraud \cite{10.1145/3603163.3609064}, social media spam \cite{mirsky2021threatoffensiveaiorganizations}, and academic dishonesty \cite{KASNECI2023102274}.
Moreover, the increasing use of LLM-generated content in training pipelines creates a recursive feedback loop \cite{alemohammad2023selfconsuminggenerativemodelsmad}, potentially degrading data quality and diversity, which poses long-term risks to both society and academia. These concerns motivate the development of detectors that reliably distinguish human-written from machine-generated text and can be deployed at scale across domains. 

Prior work on text detection can be grouped into two categories: classifier-based and statistics-based. Classifier-based detectors require training a task-specific model, which in turn hinges on collecting high-quality, domain-balanced labeled data \cite{guo2023closechatgpthumanexperts,GPTZero,guo2024detective}; this process is costly, time-consuming, and must be repeated when the target domain or generator shifts. Statistics-based methods fall into two categories: global statistics and distributional statistics. The first relies on global statistics such as likelihood or rank \cite{solaiman2019releasestrategiessocialimpacts,gehrmann-etal-2019-gltr}, which can be inaccurate or unstable under calibration mismatch, text-length variability, and domain shift. The second relies on distributional statistics, which are constructed by regenerating a neighborhood around the test passage, via sampling, perturbation, or continuation, thereby tying the detector to that particular input \cite{yang2023dnagptdivergentngramanalysis,su2023detectllmleveraginglogrank,mitchell2023detectgptzeroshotmachinegeneratedtext}. Such per-instance pipelines demand substantial compute and latency and are unrealistic when resources are constrained or throughput is high. Black-box constraints exacerbate calibration drift in global-statistic and regeneration-based detectors due to proxy-model mismatch. This motivates the development of detectors that avoid retraining and per-instance regeneration while remaining reliable under distribution shift in the black-box setting.

\begin{figure*}[t]
\centering
\includegraphics[width=0.88 \textwidth]{framework.pdf}
\caption{SurpMark framework. Offline, we build human/machine reference transition matrices by scoring corpora with a proxy LM, discretizing surprisal via a shared $q_k$, and counting state transitions. Online, a test passage is summarized the same way and assigned a GJS score to measure proximity to human vs. machine references. Details are in Algorithm~\ref{alg:surpmark-offline} and \ref{alg:surpmark-online} in Appendix~\ref{sec:alg}.}
\label{fig:framework}
\end{figure*}

Accordingly, we pursue a design that sidesteps both training-classifier and per-instance regeneration by focusing on stable, dynamics-aware signals, that can be reused across test samples. Viewed through a black-box perspective, the problem naturally invites a likelihood-free hypothesis testing formulation \cite{32134,gerber2024likelihood}:  when the true likelihood is unknown, we compare the empirical summary statistics of a test text against human and machine references. 
% Our choice of summary statistic is driven by a key empirical insight regarding the temporal flow of information. As illustrated in Figure~\ref{fig:transition-heatmap-distribution}(a), distinct patterns emerge in the generation process: while human text sustains high variability, LLMs exhibit a recovery behavior, typically reverting to highly predictable tokens immediately after generating an unexpected one. To capture this signature, which static metrics overlook, we explicitly calculate token surprisal (negative log-probability) and discretize it into interpretable states. This allows us to model the transitions between these states, yielding a representation that is robust to calibration differences and relies on relative structural dynamics rather than absolute likelihood values.
Our summary statistic design is guided by two principles. First, because the references are existing corpora whose contexts differ from the test passage, the summary must be abstract and calibration-robust; second, decisions should exploit token dynamics which expose rich local patterns \cite{xu2025trainingfree}. Empirically, we observe a characteristic “recovery” pattern—LLM text tends to return to highly predictable tokens immediately after an unexpected one (Figure~\ref{fig:transition-heatmap-distribution}(a)). We therefore quantize token surprisal into  interpretable states and summarize texts by their state-transition patterns, allowing decisions to depend on relative structure rather than absolute likelihood levels. This representation captures token dynamics and provides a stable, interpretable basis for likelihood-free comparison to human and machine references.

In this paper, we present SurpMark, a black-box, reference-based detector that frames attribution as a likelihood-free hypothesis test. For each test text, token surprisals from a proxy LM are quantized into $k$ interpretable states. The text is summarized by its state-transition matrix and is then assigned a generalized Jensen-Shannon (GJS) divergence score that measures its proximity to the human or machine reference transitions. 
These design choices motivate the theoretical analysis: Under an idealized Markov model fitted to the discretized surprisal states, we analyze how discretization affects the estimation of GJS and study the properties of our decision statistic.

% \vspace{-3mm}
\subsection{Main Contributions}
% \vspace{-2mm}
\begin{itemize}[noitemsep]
    \item We propose SurpMark, a reference-based detector, as shown in Figure~\ref{fig:framework}. 
    \item A theoretical analysis that guides the discretization bin scaling and justifies the test statistic.
    \item A comprehensive experimental evaluation of SurpMark demonstrates its effectiveness across multiple models and domains.  
\end{itemize}
% \vspace{-3mm}
\subsection{Related Work}
% \vspace{-2mm}
% Prior work on text detection can be broadly categorized into classifier-based and statistics-based methods. Classifier-based detectors train task-specific classifiers to distinguish between human-written and machine-generated text \citep{guo2023closechatgpthumanexperts,GPTZero,guo2024detective}. While effective with sufficient training data, they are costly to build and must be retrained whenever the domain or generator shifts.

% Statistics-based approaches can be divided into two groups based on their design of decision statistics. The first global-statistic methods rely on overall features of the text such as likelihood \citep{solaiman2019releasestrategiessocialimpacts}, LogRank \citep{solaiman2019releasestrategiessocialimpacts} that measures the log of each token’s rank in a model’s predicted distribution , or entropy \citep{gehrmann-etal-2019-gltr} that measures the uncertainty of a model’s next-token distribution. Distributional-statistic methods generate a neighborhood around the test passage via perturbation, continuation, or sampling, and then measure divergence between the test instance and this synthetic distribution. DetectGPT \citep{mitchell2023detectgptzeroshotmachinegeneratedtext} leverages the local curvature of log-probability function, comparing original passages with perturbed variants to enable detection of machine-generated text. Fast-DetectGPT \citep{bao2024fastdetectgpt} introduces conditional probability curvature for faster detection. DNA-GPT \citep{yang2023dnagptdivergentngramanalysis} truncates passages, and analyzes n-gram divergences of the regeneration. DetectLLM-NPR \citep{su2023detectllm} leverages normalized perturbed log-rank statistics, showing that machine-generated texts are more sensitive to small perturbations. Lastde++ \citep{xu2025trainingfree} combines global likelihood with local diversity entropy, where discretization of token probabilities stabilizes the entropy feature. In contrast, our framework discretizes token surprisals to build surprisal-state Markov transitions, enabling likelihood-free hypothesis test. Our method lies between global- and distributional-statistic approaches: it scores each text in a single pass without regeneration, yet makes comparative decisions by measuring alignment with fixed human and machine references.



% Recent work has explored kernel-based statistical tests for machine-generated text detection \citep{zhang2024detecting,song2025deep}. \cite{song2025deep} introduced R-Detect, a relative test framework that reduces false positives by comparing whether a test text is closer to human-written or machine-generated distributions. Our method shares a common foundation with \cite{song2025deep} in that it can also be viewed as a relative test framework. Notably, while the decision rules of these kernel-based approaches are non-parametric and do not rely on supervised classifiers, their optimized variants require training kernel parameters on reference corpora, which increases computational cost. Our approach only requires a lightweight data discretization stage. 

Existing detection methods are broadly categorized into classifier-based and statistics-based approaches. Classifier-based detectors \citep{guo2023closechatgpthumanexperts,GPTZero,guo2024detective} are effective but brittle, incurring high costs to retrain upon domain or generator shifts. Statistics-based approaches offer a more transferable alternative. Global-statistic methods utilize holistic features such as likelihood, LogRank \citep{solaiman2019releasestrategiessocialimpacts}, or entropy \citep{gehrmann-etal-2019-gltr}. Distributional-statistic methods, conversely, estimate local divergence via perturbations or sampling. Prominent examples include DetectGPT \citep{mitchell2023detectgptzeroshotmachinegeneratedtext} and Fast-DetectGPT \citep{bao2024fastdetectgpt} which leverage probability curvature, and DNA-GPT \citep{yang2023dnagptdivergentngramanalysis} which analyzes n-gram divergences. Recently, Lastde++ \citep{xu2025trainingfree} combined global likelihood with local diversity entropy derived from discretized probabilities.

Our framework, SurpMark, bridges these paradigms. Unlike perturbation-based methods that require costly regeneration, SurpMark scores in a single pass. Unlike global-statistic methods, SurpMark incorporates both human and machine reference corpora to perform a comparative analysis. This aligns with recent kernel-based relative tests, specifically R-Detect \citep{song2025deep}, which compares test samples against human/machine anchors. However, while R-Detect requires optimizing kernel parameters on reference corpora, SurpMark only requires a lightweight discretization process. See more discussion in Appendix~\ref{app:related-work}.
% \vspace{-3mm}
\section{SurpMark: Detailed Methodology} \label{sec:metho}
% \vspace{-2mm}
In this section, we introduce the proposed detector SurpMark.

\begin{figure*}[t]
\centering
\includegraphics[width=0.95 \textwidth]{transition-heatmap-distribution-2.pdf}
\caption{(a) The 'Recovery' Phenomenon. Visualization of the key feature driving our detector by comparing the conditional probabilities of transitioning into and out of the "Highly Surprising" state under a 5-bin discretization. This reveals distinct dynamic patterns, including a stronger recovery tendency and a more pronounced spiking tendency from low-surprisal contexts in LLM-generated text. (b) A heatmap illustrating the detector's performance (AUROC) on SQuAD across different hyperparameter settings, using a fixed amount of reference and test data. Higher orders suffer from state-space explosion and sparse transitions, yielding no notable gains beyond the first-order model (see Appendix~\ref{sec:empirical-finding-for-FO} for more discussion). (c) The final score distributions of our detector.}
\label{fig:transition-heatmap-distribution}
\end{figure*}
% \vspace{-3mm}
\paragraph{Surprisal Sequence Estimation via Proxy Model.}
% \vspace{-2mm}
Given a fixed text passage $\mathbf{t}$ and a proxy model $F_{\theta}$, we first
tokenize $\mathbf{t}$ with the tokenizer associated with $F_{\theta}$ to obtain
a token sequence $\mathbf{x} = (x_1, \dots, x_n)$ of length $n$. We then run a
single forward pass of $F_{\theta}$ on this fixed sequence to compute the
token-level surprisal sequence $\{s_t\}_{t=1}^{n}$.
\begin{align}
    \{s_t\}_{t=1}^{n} &= \{s_1,  \dots, s_n\} \nonumber \\
     &= \{ -\log p_{\theta}(x_2|x_1) ,\dots, -\log p_{\theta}(x_n|\{x_t\}_{t=1}^{n-1})  \} \nonumber 
\end{align}
where $p_\theta(\cdot \mid \cdot)$ is the conditional probability estimated by the proxy model $F_\theta$.
% \vspace{-3mm}
\paragraph{Surprisal Discretization by K-means.}
Since surprisal values from the proxy model are continuous, we discretize them into a finite set of surprisal states to enable robust statistical modeling. We employ $k$‑means clustering to partition the surprisal distribution into $k$ levels, denoted as $\mathcal{A}=\{1,\dots,k\}$. For example, when $k = |\mathcal{A}| = 4$, the clusters correspond to interpretable states such as “Predictable,” “Slightly Surprising,” “Significantly Surprising,” and “Highly Surprising.” This abstraction simplifies modeling while preserving the essential structure of predictive uncertainty. Effectively, this step converts the initial sequence of continuous surprisal values, $ \{s_t\}_{t=1}^{n} $, into a discrete state sequence, $ \{a_t\}_{t=1}^{n} $, where $a_t \in \mathcal{A}$.



% \vspace{-3mm}
\paragraph{Modeling State Transitions as Markov Chain.}
After discretizing surprisal values into finite states, we model the resulting sequence as a Markov chain. Notably, LLMs often produce a highly predictable token after a highly surprising one, a recovery effect driven by perplexity minimization, as illustrated in Figure~\ref{fig:transition-heatmap-distribution}(a). To capture this structure, we summarize each text by its empirical first-order transition frequencies. Formally, given a discretized surprisal state sequence $\{a_1, a_2, \dots, a_n\}$, we estimate a transition probability matrix $\hat M$, where each entry $\hat M(j|i)$ represents the empirical probability of transitioning from state $i$ to state $j$, with  $i,j \in \mathcal{A}$.
\begin{align}
   \hat M(j|i)=\frac{\sum_{t=1}^{n-1} \mathbf{1}\{a_t = i,\ a_{t+1} = j\}}{\sum_{t=1}^{n-1} \mathbf{1}\{a_t = i\}}, \quad i, j \in \mathcal{A} \label{eq:counts}
\end{align}
Here, $ \mathbf{1}\{\cdot\}$ is the indicator function.

% Figure~\ref{fig:transition-heatmap-distribution}(b) varies the order of the state-transition summary while keeping the reference and test sets fixed. AUROC deteriorates as the order increases, which we attribute to state-space explosion with limited data: higher-order transition counts on both the reference and test side become extremely sparse, so higher-order summaries bring no notable gains over the first-order one. Table~\ref{tab:aurox-vs-test} in Appendix further quantifies this sample-efficiency issue by sweeping the reference size and test length across different Markov orders. First-order Markov modeling consistently achieve the highest AUROC, while higher orders remain worse \textit{despite more data}. More details are in Appendix~\ref{sec:empirical-finding-for-FO}. These make first-order modeling the most practical and data-efficient choice in our setting.


% \vspace{-2mm}
\paragraph{Reference-based Detection with GJS Divergence.} 
We frame the task of distinguishing between human-written and LLM-generated text as a binary likelihood-free hypothesis testing (LFHT) problem \cite{32134, gerber2024likelihood}. To adapt LFHT for this specific domain, we introduce three key methodological modifications: (1) we utilize token-level surprisals from a fixed proxy LM as observable features; (2) we employ $k$-means quantization to transform continuous values into statistically tractable discrete state sequences; and (3) we propose a new test statistic, $\Delta \text{GJS}_n$. Crucially, unlike standard LFHT which typically evaluates divergence from a single reference distribution \cite{32134}, $\Delta \text{GJS}_n$ leverages a two-sided comparison against both human and machine references to enhance discriminative power. In this framework, the null hypothesis $H_0$ posits that the text is machine-generated, while the alternative $H_1$ suggests it is human-written. Since the true source distributions ($P$ and $Q$) are unknown, our approach remains strictly reference-based, relying on historical corpora to approximate the underlying statistics.

Specifically, given reference texts $\mathbf{t}_P, \mathbf{t}_Q$  from both model source $P$ and human source $Q$, we first compute their empirical surprisal transition probability matrices, denoted by $\hat{M}_P$ and $\hat{M}_Q$, respectively. For a given test text $\mathbf{t}$ coming from either $P$ or $Q$, we similarly compute its surprisal transition probability matrix $\hat{M}_T$ using the surprisal state levels estimated from reference texts.  We then calculate two separate divergence scores using the generalized Jensen-Shannon Divergence (GJS): one measuring the distance between the test text and the machine reference model $ \text{GJS} ( \hat{M}_{P}, \hat{M}_{T}, \alpha )$ and another measuring the distance to the human reference model $ \text{GJS} ( \hat{M}_{Q}, \hat{M}_{T}, \alpha )$, where $\alpha$ denotes the reference–test length ratio. The GJS divergence between $M_A$ and $M_B$ with weight $\alpha$ is defined as 
$
\mathrm{GJS}(M_A,M_B,\alpha) =
\tfrac{\alpha}{1+\alpha}\, D_{\mathrm{KL}}(M_A, M_\alpha)
+ \tfrac{1}{1+\alpha}\, D_{\mathrm{KL}}(M_B , M_\alpha),
M_\alpha = \tfrac{\alpha}{1+\alpha} M_A + \tfrac{1}{1+\alpha} M_B$,
where $D_{\mathrm{KL}}$ denotes the Kullback–Leibler divergence.
We score each test passage with $\Delta\text{GJS}_n = \text{GJS} \left( \hat{M}_{P}, \hat{M}_{T}, \alpha \right) - \text{GJS} \left( \hat{M}_{Q}, \hat{M}_{T},\alpha \right)$. We classify via a tunable threshold $\tau$.
% \vspace{-1mm}
\begin{align}
    \Omega  = 
    \begin{cases} 
    H_0 & \text{if  }  \Delta\text{GJS}_n\leq \tau , \\
    H_1 & \text{if  }  \Delta\text{GJS}_n > \tau 
\end{cases} 
\end{align}
% \vspace{-1mm}
See Algorithm~\ref{alg:surpmark-offline} and \ref{alg:surpmark-online} in Appendix~\ref{sec:alg} for details.

% \vspace{-3mm}
\paragraph{Choice of Markov Order.}
Figure~\ref{fig:transition-heatmap-distribution}(b) varies the order of the state-transition summary while keeping the reference and test sets fixed. AUROC deteriorates as the order increases, which we attribute to state-space explosion with limited data: higher-order transition counts on both the reference and test side become extremely sparse, so higher-order summaries bring no notable gains over the first-order one. Table~\ref{tab:aurox-vs-test} in Appendix further quantifies this sample-efficiency issue by sweeping the reference size and test length across different Markov orders. First-order Markov modeling consistently achieve the highest AUROC, while higher orders remain worse \textit{despite more data}. More details are in Appendix~\ref{sec:empirical-finding-for-FO}. These make first-order modeling the most practical and data-efficient choice in our setting.


% \vspace{-3mm}
\section{Analysis}
% \vspace{-2mm}
% \paragraph{Scope of the Markov approximation.} We do not assume that either human text or LLM text is generated by a first-order Markov process at the token level. Our analysis applies a first-order Markov approximation only to the discretized proxy-surprisal state sequence induced by a fixed proxy LM and a shared quantizer. This stylized modeling choice is used to (i) motivate transition-based statistics and (ii) enable clean finite-sample analysis of our plug-in estimator; the detector itself remains well-defined regardless of whether the underlying text generator is Markov.

% This section theoretically grounds our design choices. First, we justify the first-order Markov modeling of discretized surprisals, supported by Gray’s approximation theory and empirical evidence that second-order gains are negligible. Second, we justify $\Delta \mathrm{GJS}_n$ to be a principled choice and analyze it under this idealized framework to characterize how discretization and sample size influence its behavior.

\paragraph{Scope of the Markov approximation.}
We do not assume human or LLM text is first-order Markov at the token level.
We invoke a first-order Markov \emph{approximation} only for the discretized proxy-surprisal state sequence induced by a fixed proxy LM and a shared quantizer; this is an idealization for analyzing our plug-in transition estimator, not a claim about the true generative mechanism.

This section provides theoretical support for our design. 
First, we argue that first-order modeling on the discretized surprisal state is sufficient in our regime, using Gray’s approximation theory together with the empirical finding that the second-order improvement is negligible.
Second, under this idealized framework, we motivate $\Delta \mathrm{GJS}_n$ and derive a non-asymptotic discretization-estimation tradeoff that guides the scaling choice of the bin count $k$.

% \vspace{-3mm}
\subsection{Approximation Error with First-Order Markov Modeling} \label{sec:theoretical-Justification-fo-Markov}
% \vspace{-2mm}

% Our detector models the discretized surprisal sequence with a first-order Markov chain. Gray’s Markov approximation theory \citep{gray2011entropy}  shows that, \textit{for any stationary discrete-time process, the canonical $K$-th order Markov chain is the best $K$-th order approximation in relative-entropy-rate sense.} Specializing this to our discretized surprisal process and $K=1$, the gain from moving from first order to second order is governed by the conditional mutual information term $I(a_t; a_{t-2} \mid a_{t-1})$. Empirically, we observe this term to be negligible in our data—at most $0.0076$ bits/token (perplexity reduction $\approx 0.5\%$; see Appendix~\ref{sec:theo-justification}), indicating diminishing returns beyond first order under our discretization and budgets. 

Our detector models the discretized surprisal-state sequence $\{a_t\}$ as a first-order Markov chain. Our choice is based on a classical framework of approximating an arbitrary stationary process with a Markov Chain \citep{gray2011entropy} that we present in detail in Appendix~\ref{sec:theo-justification}. In particular, \citep{gray2011entropy} establishes that the natural plug-in estimator provides the best $K$-th order approximation and also provides a quantitative metric to compute the penalty in using the first order Markov chain over higher order choices. For example, using the results in \citep{gray2011entropy} we show in Appendix~\ref{sec:theo-justification} that the approximation error between first and second order approximations is at-most  $I(a_t; a_{t-2} \mid a_{t-1})$. Table~\ref{tab:cond-entropy-perplexity-2} demonstrates experimentally that this value is negligible in our datasets (see Table~\ref{tab:cond-entropy-perplexity} in Appendix for more details). Thus not only is the choice of first-order Markov chain beneficial in that it requires fewer parameters to learn, but it also does not incur a substantial approximation error. 

\begin{table}[htbp]
\centering
{\scriptsize
\setlength{\tabcolsep}{5pt}
\begin{tabular}{l c c}
\toprule
Source & $\hat I$ (bits/token) & Rel.\ PP gain (2nd vs.\ 1st) \\
\midrule
GPT-5-chat & 0.0076 & +0.528\% \\
Human      & 0.0045 & +0.314\% \\
\bottomrule
\end{tabular}}
\caption{Incremental second-order information for discretized surprisal states: $\hat I= I(a_t;a_{t-2}\mid a_{t-1})$. Rel.\ PP gain reports the perplexity reduction from 1st to 2nd order.}
\label{tab:cond-entropy-perplexity-2}
\end{table}



\iffalse
implies that for any stationary finite-alphabet process, the \emph{canonical} $K$-th order Markov chain minimizes the relative-entropy rate among all $K$-th order Markov chains on the same alphabet. 
In our setting, this connects the benefit of increasing the Markov order to a measurable quantity: the per-token reduction in optimal log-loss (entropy rate) from conditioning on one additional lag satisfies
$
H(a_t \mid a_{t-1}) - H(a_t \mid a_{t-1}, a_{t-2})
= I(a_t; a_{t-2} \mid a_{t-1}).
$
Empirically, this conditional mutual information is negligible across our datasets---at most $0.0076$ bits/token (corresponding to $\approx 0.5\%$ perplexity reduction; Appendix~\ref{sec:theo-justification}).

Thus, under our discretization and realistic reference/test budgets, moving beyond first order offers diminishing predictive gains.

\fi 

% \vspace{-2mm}
% \subsection{Analysis Under Idealized First-Order Markov Modeling} \label{sec:setup}
% \vspace{-2mm}
% Let $\{s_t^P\}_{t=1}^N$ and $\{s_t^Q\}_{t=1}^N$ be the surprisal sequences produced by a fixed
% proxy LM on reference corpora from the model source $P$ and the human source $Q$,
% respectively. For the purpose of analysis, we work with a first-order Markov
% \textit{approximation} to these surprisal processes on the real line $\mathbb{R}$, and write
% $\mathcal{S}_P$ and $\mathcal{S}_Q$ for the corresponding transition kernels. This should be viewed as a
% stylized model for the discretized surprisal dynamics, rather than a literal assumption
% that the underlying language model is exactly first-order Markov. For an integer $k \ge 2$, let $q_k : \mathbb{R} \to \mathcal{A} = \{1,\dots,k\}$ be a shared
% quantizer with boundaries $b_1 < \cdots < b_{k-1}$, and define discretized states
% $a_t^P = q_k(s_t^P)$ and $a_t^Q = q_k(s_t^Q)$. The induced $k$-state Markov chains
% have transition kernels
% \vspace{-2mm}
% \begin{align}
%     M_P(j \mid i) = \Pr[a_{t+1}^P = j \mid a_t^P = i],\nonumber \\
% M_Q(j \mid i) = \Pr[a_{t+1}^Q = j \mid a_t^Q = i], \nonumber
% \end{align}
% and their plug-in estimators $\hat M_P, \hat M_Q$ are formed from transition counts as in Eq.~\ref{eq:counts}. We assume the discretized chains are ergodic with unique stationary distributions $\pi_P$ and $\pi_Q$ (see Lemma~\ref{lem:quant-stationary} in Appendix for a formal proof). For our theoretical analysis, we define the minimum state probability as $\pi_{\min} := \min\{\min_{s\in\mathcal{A}} \pi_P(s), \min_{s\in\mathcal{A}} \pi_Q(s)\}$.
% Finally, we observe an independent test surprisal-state sequence
% $a^T_{1:n} = \{a_t^T\}_{t=1}^n$ whose source $M_T$ is either $M_P$ (null $H_0$) or
% $M_Q$ (alternative $H_1$), and all three sequences share the same quantizer $q_k$.


% \vspace{-2mm}
% \subsubsection{Discretization Effect}\label{sec:quant-effect}
% \vspace{-2mm}
% How should we choose the number of bins $k$? Too few bins lose structural information, while too many, given a fixed-length reference, lead to sparse counts, higher estimation noise, and bias from zero-count corrections. Thus, $k$ must balance information preservation and statistical reliability.

% Following \cite{pillutla2023mauvescoresgenerativemodels}, we analyze discretization through a two-term decomposition. Discretization error is a deterministic bias from projecting the continuous object onto $k$ bins, while the statistical error is the finite-sample discrepancy when estimating the discretized object. \cite{pillutla2023mauvescoresgenerativemodels} study IID samples, and control the statistical error by splitting observed vs. unobserved mass and derive non-asymptotic bounds when balanced with their quantization error. Rather than assuming IID samples, we focus on Markov sources and examine empirical transition counts from their sequences.

% For a divergence functional $\mathcal{D}_f$ (we use row-wise GJS), the empirical estimator is $\mathcal{D}_f(\hat M_P,\hat M_Q)$. Our goal is to develop a non-asymptotic bound on the absolute error of the empirical estimator relative to the true target, decomposed as 
% \begin{align}
%     \underbrace{|\mathcal{D}_f(\mathcal{S}_P,\mathcal{S}_Q)-\mathcal{D}_f(M_P,M_Q)|}_{\text{discretization error}}+ \nonumber\\ \underbrace{|\mathcal{D}_f(\hat M_P, \hat M_Q)-\mathcal{D}_f(M_P,M_Q)|}_{\text{statistical error}} \nonumber
% \end{align}
% where $\mathcal{S}_P,\mathcal{S}_Q$ denote the underlying Markov transition kernels. For simplicity we take both references to have the same total lengtht he $N$. $C$ denotes an absolute constant that may change from line to line.
% \vspace{-2mm}
% \paragraph{Discretization Error and Statistical Error.}  At a high level, we decompose the total error into a discretization bias and a finite-sample statistical term. The discretization bias is controlled by adapting \citet{pillutla2023mauvescoresgenerativemodels} to our Markov setting and yields an $O(1/k)$ bound. Theorem 4.2 then bounds the statistical error by tracking three sources—row-wise transition noise, missing-mass from unseen transitions, and an additional stationary-weight estimation error specific to Markov chains—showing how the overall error trades off $k$ and $N$.

% \begin{proposition}\label{prop:quant-error}
%     Let $\mathcal{S}_P,\mathcal{S}_Q$ be the population first-order Markov transition kernels on the continuous surprisal space $\mathbb{R}$. Consider a shared $k$-bin quantizer $q_k:\mathbb{R} \rightarrow \mathcal{A}$ and, from it, form the discretized $k$-state Markov chains $M_P, M_Q$. For any row-aggregated $f$-divergence functional $\mathcal{D}$, there exists such a shared $k$-bin partition satisfying 
%     \begin{align}
%         |\mathcal{D}_f(\mathcal{S}_P,\mathcal{S}_Q)-\mathcal{D}_f(M_P,M_Q)| \le \frac{C}{k}
%     \end{align}
% \end{proposition}

% See Appendix~\ref{sec:proof-quant-error} for the proof. 

% \begin{theorem}\label{thm:sta-error}
%     Suppose we are in the setting described in Section~\ref{sec:setup}. Assume each discretized chain is ergodic with bounded mixing time, $\pi_{\min}\gtrsim 1/k$, and maximum hitting time $\max\{T(M_P),T(M_Q)\}=O(1)$. It holds that 
%     \begin{align}
%         &|\mathcal{D}_f(\hat M_P, \hat M_Q)-\mathcal{D}_f(M_P,M_Q)|\nonumber \\ &\leq  C\Big(\log N\cdot \sqrt{\frac{k^3\log(kN)}{N}}+\frac{k^3}{N}\log\big(1+\frac{N}{k}\big)+\frac{k}{\sqrt{N}} \Big) \nonumber
%     \end{align}
% \end{theorem}
% See Appendix~\ref{sec:proof-sta-error} for the proof.
% \begin{remark}
% In Theorem~\ref{thm:sta-error}, $N$ refers to the number of transitions for brevity; this $O(1)$ difference from the sequence length does not affect the asymptotic bounds.
% \end{remark}
% \vspace{-3mm}
% \paragraph{Balancing Two Errors.} We balance $k$ by trading off the discretization bias against the finite-sample statistical error. The discretization term decays as $O(k^{-1})$, while the leading statistical term from row-wise transition estimation grows like $k^{\frac{3}{2}}/\sqrt{N}$ up to logs, with smaller contributions $O(k/\sqrt{N})$ and $O(k^3/N)$ for $k\ll N^{\frac{1}{3}}$. Neglecting logs and lower-order terms, the dominant balance is between $c_1 k^{\frac{3}{2}}/\sqrt{N}$ and $c_2/k$, yielding $k^* = \Theta(N^{\frac{1}{5}})$. 


\subsection{Analysis Under an Idealized First-Order Markov Approximation}\label{sec:setup}

Let $\{s_t^P\}_{t=1}^N$ and $\{s_t^Q\}_{t=1}^N$ denote proxy-LM surprisal sequences from the model source $P$ and human source $Q$. For analysis, let $\mathcal{S}_P$ and $\mathcal{S}_Q$ denote the (idealized) population first-order transition kernels on the continuous surprisal space $\mathbb{R}$ for the two sources.
Fix an integer $k\ge2$ and a shared quantizer $q_k:\mathbb{R}\to\mathcal{A}=\{1,\dots,k\}$, and define discretized states $a_t^P=q_k(s_t^P)$ and $a_t^Q=q_k(s_t^Q)$. 
The induced $k$-state transition kernels are
% \vspace{-2mm}
\begin{align}
&M_P(j\mid i):=\Pr[a_{t+1}^P=j\mid a_t^P=i],\quad \nonumber\\
& M_Q(j\mid i):=\Pr[a_{t+1}^Q=j\mid a_t^Q=i]. \nonumber
\end{align}
Their plug-in estimators $\hat M_P,\hat M_Q$ are formed from transition counts (Eq.~\ref{eq:counts}). 
We assume ergodicity with stationary distributions $\pi_P,\pi_Q$ and define $\pi_{\min}:=\min\{\min_{s}\pi_P(s),\min_{s}\pi_Q(s)\}$. 
Given an independent test sequence $a^T_{1:n}$, we test whether its source transition kernel $M_T$ equals $M_P$ (null) or $M_Q$ (alternative), using the same $q_k$.

% \vspace{-3mm}
\subsubsection{Discretization--Estimation Tradeoff}\label{sec:quant-effect}
% \vspace{-2mm}
How should we choose the number of bins $k$? Too few bins lose structural information, while too many, given a fixed-length reference, lead to  higher estimation noise. Thus, $k$ must balance information preservation and statistical reliability.
For a row-aggregated $f$-divergence functional $\mathcal{D}_f$ (row-wise GJS in our case), we study
\begin{align}
\underbrace{|\mathcal{D}_f(\mathcal{S}_P,\mathcal{S}_Q)-\mathcal{D}_f(M_P,M_Q)|}_{\text{discretization error}}& + \nonumber\\ \underbrace{|\mathcal{D}_f(\hat M_P, \hat M_Q)-\mathcal{D}_f(M_P,M_Q)|}_{\text{statistical error}} \nonumber
\end{align}
where $\mathcal{S}_P,\mathcal{S}_Q$ denote the population kernels on $\mathbb{R}$ and $M_P,M_Q$ their $k$-bin discretizations.

\begin{proposition}\label{prop:quant-error}
For any row-aggregated $f$-divergence $\mathcal{D}_f$, there exists a shared $k$-bin partition such that
$
|\mathcal{D}_f(\mathcal{S}_P,\mathcal{S}_Q)-\mathcal{D}_f(M_P,M_Q)| \le C/k.
$
\end{proposition}

\begin{theorem}\label{thm:sta-error}
In the setting of Section~\ref{sec:setup}, assume each discretized chain is ergodic with bounded mixing and $\pi_{\min}\ge c/k$ for some absolute constant $c\ge0$.
It holds that
\begin{align}
        &|\mathcal{D}_f(\hat M_P, \hat M_Q)-\mathcal{D}_f(M_P,M_Q)|\nonumber \\ &\leq  C\Big(\log N\cdot \sqrt{\frac{k^3\log(kN)}{N}}+\frac{k^3}{N}\log\big(1+\frac{N}{k}\big)+\frac{k}{\sqrt{N}} \Big) \nonumber
    \end{align}
% where $\tilde O(\cdot)$ hides polylog factors in $k$ and $N$ (see Appendix~\ref{sec:proof-sta-error} for the explicit form and proof).
\end{theorem}
\begin{remark}
In Theorem~\ref{thm:sta-error}, $N$ refers to the number of transitions for brevity; this $O(1)$ difference from the sequence length does not affect the asymptotic bounds.
\end{remark}
% \vspace{-4mm}
\paragraph{Balancing Two Errors.}
Balancing $O(1/k)$ with the dominant estimation term $ O(k^{3/2}/\sqrt N)$ yields
$
k^*=\Theta(N^{1/5})
$
(up to polylog factors).




% \vspace{-3mm}
\subsubsection{Decision Statistic Analysis}
% \vspace{-2mm}
Building on this discretized Markov approximation, we next analyze the decision
statistic $\Delta \mathrm{GJS}_n$. Our detector extends Gutman’s universal hypothesis test \citep{32134} from a single-reference setting to a two-reference setting. In Gutman’s test, the test sequence is compared against one reference source; here we leverage two calibrated references $P$ (LM) and $Q$ (human) and decide by $\Delta \text{GJS}_n$. Our choice of GJS is not ad hoc. Algebraically, $\Delta \text{GJS}_n$ is the log–likelihood ratio (LLR) between the two hypotheses. 
% \vspace{-4mm}
\paragraph{$\Delta \text{GJS}_n$ as Log-Likelihood Ratio.} Proposition~\ref{prop:GJS-n-as-LLR} shows that $\Delta \text{GJS}_n$ exactly equals the normalized log-likelihood ratio $\Lambda_{n,N}$. Here, the log-likelihood ratio represents the maximized data likelihood under the two hypotheses $H_0$ and $H_1$. See Appendix~\ref{sec:proof-llr} for the proof.

\begin{proposition}\label{prop:GJS-n-as-LLR}
     Assume the setting of Section~\ref{sec:setup}. Let $\mathcal{F}_k$ be the family of stationary first-order Markov models on $\mathcal{A}:=[k]$. For sequences $a^P_{1:N}$, $a^Q_{1:N}$, and $a^T_{1:n}$, define the concatenations $(a^P_{1:N},a^T_{1:n})$ and $(a^Q_{1:N},a^T_{1:n})$.  Consider the generalized log-likelihood ratio $ \Lambda_{n,N}$
     \begin{align}
      \Lambda_{n,N} =   \frac{1}{n}\log\!
\frac{\displaystyle \sup_{M,M'\in\mathcal{F}_k}
M\!\bigl((a^P_{1:N},a^T_{1:n})\bigr)\,M'\!\bigl(a^Q_{1:N}\bigr)}
{\displaystyle \sup_{M,M'\in\mathcal{F}_k}
M\!\bigl(a^P_{1:N}\bigr)\,M'\!\bigl((a^Q_{1:N},a^T_{1:n})\bigr)}
     \end{align}
     where the suprema are attained at the empirical Markov models on the respective concatenated sequences. Then,
         $\Delta \rm{GJS}_n =  \Lambda_{n,N}$.

\end{proposition}
In Appendix~\ref{sec:proof-asymp-gaussian}, we further prove the asymptotic normality of our statistics $\Delta \text{GJS}_n$, and empirically verify it through experiments in Appendix~\ref{sec:shapiro-wilk}.

\begin{table*}[h]
\centering
{
\tiny 
\resizebox{\textwidth}{!}{
\begin{tabular}{lcccccccccc}
\toprule
& GPT2-XL & GPT-J-6B & GPT-Neo-2.7B & GPT-NeoX-20B & OPT-2.7B & Llama-2-13B & Llama-3-8B & Llama-3.2-3B & Gemma-7B  & Avg \\
\midrule
Likelihood & 85.02 & 74.82 & 73.32 & 72.03 & 77.22 & 94.39 & 93.93 & 65.22 & 65.8 &77.97 \\

LogRank & 88.2 & 79.25 & 78.29 & 75.37 & 81.99 & 95.9 & 95.05 & 71.04 & 69.18 & 81.59\\

Entropy & 51.1 & 47.15 & 50.94 & 45.94 & 48.88 & 29.03 & 29.31 & 53 & 46.85& 44.69 \\

DetectLRR & 91.07 & 85.81 & 87.12 & 80.27 & 88.48 & 96.43 & 94.85 & 81.54 & 75.5 & 86.79 \\

Lastde & 95.97 & 85.88 & 89.09 & 80.16 & 88.89 & 93.29 & 94.29 & 72.99 & 69.48  & 85.56\\

Lastde++ & \textbf{99.46} & 91.54 & 94.29 & 85.13 & 94.15 & 95.5 & 95.9 & 77.47 & \underline{76.9} & 90.04\\

DNA-GPT & 81.98 & 70.68 & 72.69 & 70.42 & 73.86 & 95.91 & 96.54 & 64.79 & 65.32  &76.91\\

Fast-DetectGPT & 97.94 & 86.83 & 89.15 & 83.17 & 90.55 & \textbf{98.21}  & \textbf{97.98} & 74.32 & 73.95 & 88.01 \\

DetectGPT &94.45 & 79.55 & 84.71 & 75.71 & 82.88 & 86.51 & 86.28 & 64.23 & 69.05  & 80.37\\

DetectNPR & 94.93 & 81.91 & 86.4 & 77.93 & 84.06 & 95.19 & 93.67 & 69.45 & 71.49 & 83.89 \\

R-Detect & 74.38 & 63.58 & 67.7 & 63.35 & 65.83 & 79.98 & 81.64 & 62.22 &55.46 &  68.24\\

Binoculars & \underline{99.19} &85.76  & 87.5 & 82.02 &  86.9 & 96.93 & 96.41 & 68.36 & 73.57 &  86.19\\

FourierGPT & 54.72 & 54.28 & 56.5 & 56.51 &   52.47 & 72.43 & 72.06 & 54.83 & 55.84 & 59.07 \\

\rowcolor{ours}$\text{SurpMark}_{k=6}$ & 98.07 & \underline{92.96} & \underline{95.19} & \textbf{86.78} & \underline{94.49} & 97.41 & 97.06 & \textbf{81.74} & \textbf{77.40} &\textbf{91.23}\\

\rowcolor{ours}$\text{SurpMark}_{k=7}$ & {98.35} & \textbf{93.1} & \textbf{95.42} & \underline{86.40} & \textbf{94.88} & \underline{97.58} & \underline{97.17} & \underline{80.74} & 76.89   &\underline{91.17}\\
\bottomrule
\end{tabular}}}
\caption{Detection results for text generated by 9 open-source models under the black-box setting. The AUROC reported for each model are averaged across three datasets: Xsum, WritingPrompts, and SQuAD. See Table~\ref{tab:opensource-xsum},~\ref{tab:opensource-writing},~\ref{tab:opensource-squad} in Appendix for details.}
\label{tab:opensource}
\end{table*}

\section{Experiments} \label{sec:exp}

\paragraph{Datasets, Configurations and Models.} We evaluate our method on XSum \cite{narayan-etal-2018-dont}, WritingPrompts \cite{fan-etal-2018-hierarchical}, SQuAD \cite{rajpurkar-etal-2016-squad}, WMT19 \cite{barrault-etal-2019-findings}, HC3 \cite{guo2023closechatgpthumanexperts}, DetectRL \cite{wu2024detectrl}, Kaggle DAIGT V2 Train Dataset, and Dolly \cite{DatabricksBlog2023DollyV2}. For Xsum, WritingPrompts, SQuAD, and WMT19, we construct the reference corpora and test set as follows. For each dataset, we randomly sample 300 human-written texts to form the human reference, then generate paired machine outputs by prompting the source model with the first 30 tokens of each human text. For the test set, we sample another 150 human-written texts and create their machine-generated counterparts using the same procedure. For datasets that already provide machine-generated texts, we use the same sampling configuration, except that no additional generation step is required. We select 9 open-source models and 3 closed-source models as our source model. More details are in Appendix~\ref{sec:exp-configuration}. Unless otherwise specified, we use GPT2-Large as our proxy model.

\paragraph{Baselines.}
We benchmark against 13 detectors, including Likelihood \cite{solaiman2019releasestrategiessocialimpacts}, LogRank \cite{solaiman2019releasestrategiessocialimpacts}, Entropy \cite{gehrmann-etal-2019-gltr,ippolito2020automaticdetectiongeneratedtext}, DetectLRR \cite{su2023detectllm}, and Lastde \cite{xu2025trainingfree},  DetectGPT \cite{mitchell2023detectgptzeroshotmachinegeneratedtext}, Fast-DetectGPT \cite{bao2024fastdetectgpt}, DNA-GPT \cite{yang2023dnagptdivergentngramanalysis}, DetectNPR \cite{su2023detectllm}, Lastde++ \cite{xu2025trainingfree}, R-Detect \cite{song2025deep}, Binoculars \cite{hans2024spotting}, and FourierGPT \cite{xu2024detectingsubtledifferenceshuman}. 


\subsection{Main Results}

Table~\ref{tab:opensource} and ~\ref{tab:closesource} present the detection results under black-box scenario. Table~\ref{tab:closesource} shows that SurpMark achieves the best performance on 3 commercial, closed-source LLM. Performance is especially strong on GPT-5-Chat. Table~\ref{tab:opensource} shows that SurpMark ranks first on 6 of 9 open-source models and within the top two on 7 of 9. These results highlight SurpMark’s robustness on proprietary systems and its suitability for real-world commercial deployments. Please note that compared with distribution-based detectors that generate a neighborhood per input at test time, for each dataset, SurpMark builds reference corpora once and reuses them for all test passages. Under a reference-per-test budget $B=\frac{\#\rm{references}}{\# \rm{tests}}$, in Table~\ref{tab:closesource} and \ref{tab:opensource}, SurpMark operates at $B=2$, whereas DNA-GPT uses $B=10$, DetectGPT, DetectNPR, Lastde++ require $B=100$. Thus SurpMark’s reference cost is 5$\times$–50$\times$ lower, while avoiding any per-input contrastive generation at test time, enabling real-time detection as discussed later. 

We attribute SurpMark's superior performance on closed-source models (e.g., GPT-5-chat) to its ability to capture transitional dynamics. As detailed in Appendix~\ref{marginal-transition}, stronger models exhibit a vanishingly small gap in marginal surprisal distributions compared to humans, rendering marginal statistics ineffective. However, the 'transition gap' remains significant, which SurpMark effectively exploits.

% Remarkably, SurpMark exhibits \textit{strong cross-domain robustness}. As shown in Appendix~\ref{app:cross-domain}, using out-of-domain references yields comparable performance to in-domain ones.
\begin{table}
{\scriptsize
\centering
\begin{tabular}{lcccc}
\toprule
& Gemini-1.5-Flash & GPT-4.1-mini & GPT-5-Chat & Avg  \\
\midrule
Likelihood & 56.49 & 66.77 & 49.62 &  57.63 \\

LogRank & 53.87 & 66.8 & 49.83 & 46.53 \\

Entropy & 58.36 & 38.72 & 46.99 & 48.02 \\

DetectLRR & 44.51 & 63.29 & 49.83 & 62.11  \\

Lastde & 48.13 & 57.28& 41.96 & 49.12 \\

Lastde++ & 71.72 & 68.23 & 43.51 & 61.15 \\

DNA-GPT &62.06 & 56.71 & 49.82&  56.2\\

Fast-DetectGPT & 72.49 & 68.32 & 43.39 &  61.4\\

DetectGPT &69.19 & 70.08 & 54.6 & 64.75\\

DetectNPR & 64.96 & 70.83 & 54.99 & 63.59 \\

R-Detect &69.25 & 71.64 & 67.75 & 69.55 \\

Binoculars &74.51 & 71.12 & 49.65 & 65.09 \\

FourierGPT &61.25 & 63.05 & 64.82 &  63.04\\

\rowcolor{ours}$\text{SurpMark}_{k=6}$ & \underline{74.57} & \textbf{80.25} & \underline{78.33} & \underline{77.72} \\

\rowcolor{ours}$\text{SurpMark}_{k=7}$ & \textbf{75.14} & \underline{78.48} & \textbf{81.33} & \textbf{78.32}\\
\bottomrule
\end{tabular}}
\caption{Detection results for text generated by 3 closed-source models under the black-box setting. The AUROC reported for each model are averaged across three datasets: Xsum, WritingPrompts, and SQuAD. See Table~\ref{tab:closesource-xsum},~\ref{tab:closesource-writing},~\ref{tab:closesource-squad} in Appendix for details.}
\label{tab:closesource}
\end{table}

\subsection{Ablation and Sensitivity Analysis} 


\begin{table}[htbp]
\centering
{\tiny
\begin{tabular}{l l c c c}
\toprule
Dataset@Source model      & Method              & $k=7$ & $k=8$ & Best \\
\midrule
         & k-means            & 80.42 & 79.32 & \textbf{80.42} \\
XSum@GPT-4.1-mini          & equal-width  & 80.41 & 76.20  & 80.41 \\
        & equal-mass   & 71.40 & 74.22 & 74.22 \\
\midrule
  & k-means             & 82.05 & 84.86 & \textbf{84.86} \\
WritingPrompts@GPT5-chat  & equal-width & 62.90 & 72.06  & 72.06 \\
  & equal-mass    & 84.42 & 83.59 & 84.42 \\
\bottomrule
\end{tabular}
\caption{AUROC of different discretization schemes under varying number of states $k$.}
\label{tab:discretization-k}}
\end{table}
\vspace{-2mm}
\paragraph{Necessity of $k$-means.}


In Table~\ref{tab:discretization-k}, we evaluate the effect of different discretization schemes, including $k$-means, equal-width, and equal-mass binning. Across all datasets and $k$ values, $k$-means is the most robust quantization scheme: it consistently reaches or matches the best AUROC, while equal-mass can degrade sharply on XSum@GPT-4.1-mini and equal-width is unstable and often much worse.
\begin{figure}
\centering
  \includegraphics[width=0.38\linewidth]{gpt-j-6b-k-impact.png}
  \label{fig:sub1}
  \includegraphics[width=0.38\linewidth]{llama3.2-3b-k-impact.png}
  % \caption{QNLI}
  \label{fig:sub2}
 \caption{Effect of the number of bins 
$k$ on detection performance for source models including GPT-J-6B (left) and Llama-3.2-3B (right).}
 \label{fig:impact-of-k}

 \end{figure}
\begin{figure*}[!t]
\centering
\includegraphics[width=0.9\textwidth]{exp-length.pdf}
\caption{(a) AUROC vs. number of reference samples. The blue curve (“$k$-optimized”) picks the best $k$ at each number of reference. orange/green curves fix $k \in \{7,8\}$. (b) AUROC vs. test length $n$ under different reference lengths. Solid lines are k-optimized for each reference sample truncated to 50/100/200 tokens; shaded bands show the attainable range across $k$ at each $n$. (c) Detection results of 7 detection methods on 6 test lengths.}
\label{fig:k-optmized-k7-k8}
\end{figure*}
\begin{figure*}[!t]
\centering
\includegraphics[width=1 \textwidth]{contour_proxy_model.pdf}
\caption{(a-b) AUROC contour maps (WritingPrompts@Gemma-7B). Left: $k=7$; right: $k=8$. The x-axis is reference length (tokens) and the y-axis is test length (tokens). Colors encode AUROC. In both panels, contours tilt up-right, indicating a trade-off: larger reference length allows smaller test length at similar performance. (c) AUROC vs. proxy model.}
\label{fig:contour-proxy}
\end{figure*}
% \vspace{-3mm}
\vspace{-2mm}
\paragraph{Effect of bins $k$.} Figure~\ref{fig:impact-of-k} shows the effect of the number of bins $k$. Across both models, increasing the number of bins $k$ leads to clear improvements in AUROC up to a moderate range, after which the gains saturate or slightly decline. The best results across datasets are generally observed at $k=6-7$. Our theory yields an optimal bin count of the form $k^* = C N^{1/5}$ for some constant $C$, where $N$ is the total length of the reference samples. We empirically calibrate the constant $C$ in Appendix~\ref{app:choice-of-k}. Next, we further investigate how varying $N$ shift the empirical optimum $k$.

\vspace{-3mm}
\paragraph{Effect of Number of Reference Samples.}
Figure~\ref{fig:k-optmized-k7-k8} (a) shows that AUROC improves sharply as the reference grows from very small number of reference samples to 100 reference samples; beyond 100 reference samples the gains are minor. The $k$-optimized curve picks the best $k \in \{4,\dots,12\}$ at each number of reference. The annotated $k$ values grow mildly with the number of reference samples, and using large $k$ for small number of reference hurts performance. This trend aligns with our theoretical intuition: a larger number of reference samples reduces reference-side estimation error and thus allows for a slightly larger $k$.
\vspace{-3mm}
\paragraph{Effect of Length of Test Sample.} In Figure~\ref{fig:k-optmized-k7-k8} (b), we fix the number of reference samples and study the effect of sample length. AUROC climbs rapidly as test length $n$ grows from 50 to about 150–200. Longer reference lift the curves and make the bands across $k \in \{4, \dots,12\}$ tighter, indicating greater stability. The $k$-optimized curves show that the optimal $k$ is driven more by reference length than by test length. In Figure~\ref{fig:k-optmized-k7-k8} (c), we evaluated detection performance of baselines across varying test length (tokens), focusing on WritingPrompts generated by Gemma-7B. All methods improve with longer texts. SurpMark is competitive at short lengths and becomes the top method for test length larger than 150. Comparison on more source models are presented in Figure~\ref{fig:appendix-test-length} in Appendix.
\vspace{-3mm}
\paragraph{Reference–Test Length Trade-Offs.} Figure~\ref{fig:contour-proxy} (a) and (b) show AUROC contours over reference length and test length $n$ at fixed bins $k$. Performance improves toward the upper-right, and the up-right tilt shows a reference-test length trade-off: larger reference length can compensate for smaller test length at similar accuracy. 

% \begin{figure*}[t]
% \centering
% \includegraphics[width=1 \textwidth]{exp-length.pdf}
% \caption{(a) AUROC vs. number of reference samples. The blue curve (“$k$-optimized”) picks the best $k$ at each number of reference. orange/green curves fix $k \in \{7,8\}$. (b) AUROC vs. test length $n$ under different reference lengths. Solid lines are k-optimized for each reference sample truncated to 50/100/200 tokens; shaded bands show the attainable range across $k$ at each $n$. (c) Detection results of 7 detection methods on 6 test lengths.}
% \label{fig:k-optmized-k7-k8}
% \end{figure*}
\vspace{-3mm}
\paragraph{Effect of Proxy Model.} In Figure~\ref{fig:contour-proxy} (c), x-axis lists the proxy LM used to compute scores. Across both datasets, most baselines improve with stronger proxy models, especially on WritingPrompts with GPT-5-Chat as the source model. SurpMark is consistently top and stable across proxy models. It already performs strongly with the smallest proxy and improves only modestly with larger ones, whereas several baselines are highly sensitive to the proxy choice, some even degrade when the proxy changes. In short, SurpMark achieves strong and reliable performance without expensive proxy models, making it a better default in low-resource deployments.
% \begin{figure*}
% \centering
% \includegraphics[width=1 \textwidth]{contour_proxy_model.pdf}
% \caption{(a-b) AUROC contour maps (WritingPrompts/Gemma-7B). Left: $k=7$; right: $k=8$. The x-axis is reference length (tokens) and the y-axis is test length (tokens). Colors encode AUROC. In both panels, contours tilt up-right, indicating a trade-off: larger reference length allows smaller test length at similar performance. (c) AUROC vs. proxy model.}
% \label{fig:contour-proxy}
% \end{figure*}

\vspace{-3mm}
\paragraph{Throughput.} Figure~\ref{fig:throughput} plots throughput (items/s) against the number of test texts. Baseline methods appear as horizontal lines because their per-item latency is constant. SurpMark improves monotonically as the one-time preprocessing cost is amortized. The curve crosses the Fast-DetectGPT line at roughly $298$, after which SurpMark maintains higher throughput.
% \vspace{-3mm}
\begin{figure}[htbp]
\centering
\includegraphics[width=0.36 \textwidth]{items_per_sec_vs_n_clean.png}
\caption{Throughput (items per second) versus the number of test texts for SurpMark compared to baseline methods (proxy LM: GPT-2 Large; GPU: NVIDIA RTX 4090).}
\label{fig:throughput}
\end{figure}
% \vspace{-3mm}
\begin{figure}[htbp]
\centering
\includegraphics[width=0.36\textwidth]{auroc_grouped_datasets.png}
\caption{AUROC on non-English datasets (HC3-Chi-psy/qa/med and WMT-De). Error bars denote standard deviation. Higher is better.}
\label{fig:non-english}
\end{figure}
% \begin{figure*}[!t]
% \centering
% \includegraphics[width=0.9\textwidth]{exp-length.pdf}
% \caption{(a) AUROC vs. number of reference samples. The blue curve (“$k$-optimized”) picks the best $k$ at each number of reference. orange/green curves fix $k \in \{7,8\}$. (b) AUROC vs. test length $n$ under different reference lengths. Solid lines are k-optimized for each reference sample truncated to 50/100/200 tokens; shaded bands show the attainable range across $k$ at each $n$. (c) Detection results of 7 detection methods on 6 test lengths.}
% \label{fig:k-optmized-k7-k8}
% \end{figure*}
% \begin{figure*}[!t]
% \centering
% \includegraphics[width=1 \textwidth]{contour_proxy_model.pdf}
% \caption{(a-b) AUROC contour maps (WritingPrompts@Gemma-7B). Left: $k=7$; right: $k=8$. The x-axis is reference length (tokens) and the y-axis is test length (tokens). Colors encode AUROC. In both panels, contours tilt up-right, indicating a trade-off: larger reference length allows smaller test length at similar performance. (c) AUROC vs. proxy model.}
% \label{fig:contour-proxy}
% \end{figure*}
\subsection{More Scenarios}
\vspace{-2mm}
\paragraph{Non-English Scenarios.} In Figure~\ref{fig:non-english}, we evaluate on German and Chinese corpora. For German, we use WMT19 with GPT-4o-mini as the source model and Llama-3.2-1B as the proxy model. For Chinese, we use HC3 across multiple domains (psychology, medicine, openqa), which provides paired human and ChatGPT answers to the same questions, and adopt Qwen-2.0-0.5B as the proxy model. SurpMark ranks first on all four datasets, with large margins on HC3-Chi-med.
\vspace{-2mm}
\begin{table}[t]
\centering
\caption{Mixed-domain results on the DetectRL benchmark. Results are grouped by source model. For each source model, the test set contains samples from multiple domains, including arXiv, WritingPrompts, Yelp\_review, and XSum. We report AUROC and TPR@FPR=5\% with standard deviation.}
\label{tab:mix_domain}
{\scriptsize
\begin{tabular}{llcc}
\toprule
Source model & Method & AUROC & TPR@FPR=5\% \\
\midrule
\multirow{3}{*}{Llama-2-70B}
& Fast-DetectGPT & 79.85 \com 1.3 & 55.33 \com 0.32 \\
& Lastde++ & 65.2 \com 0.6 & 44.67 \com 0.83 \\
& \cellcolor{ours} SurpMark & \cellcolor{ours} 94.7 \com 0.6 & \cellcolor{ours} 62.67 \com 0.62 \\
\midrule
\multirow{3}{*}{Claude-instant}
& Fast-DetectGPT & 37.16 \com 0.22 & 4.67 \com 0.2 \\
& Lastde++ & 39.11 \com 0.14 & 2 \com 0.12 \\
& \cellcolor{ours} SurpMark & \cellcolor{ours} 80.92 \com 0.57 & \cellcolor{ours} 34 \com 0.88 \\
\midrule
\multirow{3}{*}{GPT-3.5-turbo}
& Fast-DetectGPT & 74.77 \com 0.68 & 42.67 \com 0.14 \\
& Lastde++ & 63.94 \com 1.06 & 25.33 \com 1.6 \\
& \cellcolor{ours} SurpMark & \cellcolor{ours} 80.57 \com 0.36 & \cellcolor{ours} 47 \com 0.63 \\
\bottomrule
\end{tabular}
}
\end{table}
% \begin{table*}[hbtp]
% \centering
% \caption{Mixed-domain results on the DetectRL benchmark. Results are grouped by source model. For each source model, the test set contains samples from multiple domains, including arXiv, WritingPrompts, Yelp\_review, and XSum. We report AUROC and TPR@FPR=5\% with standard deviation.}
% \label{tab:mix_domain}
% {\scriptsize
% \begin{tabular}{lcc|cc|cc}
% \toprule
% \multirow{2}{*}{DetectRL} 
% & \multicolumn{2}{c}{Llama-2-70B} 
% & \multicolumn{2}{c}{Claude-instant} 
% & \multicolumn{2}{c}{GPT-3.5-turbo} \\
% \cmidrule(lr){2-3} \cmidrule(lr){4-5} \cmidrule(lr){6-7}
% & AUROC & TPR@FPR=5\% 
% & AUROC  & TPR@FPR=5\% 
% & AUROC  & TPR@FPR=5\%  \\
% \midrule
% Fast-DetectGPT 
% & 79.85 \com 1.3 & 55.33 \com 0.32
% & 37.16 \com 0.22 & 4.67 \com 0.2
% & 74.77 \com 0.68 & 42.67 \com 0.14 \\

% Lastde++ 
% & 65.2 \com 0.6 & 44.67 \com 0.83
% & 39.11 \com 0.14 & 2 \com 0.12
% & 63.94 \com 1.06 & 25.33 \com 1.6 \\

% \rowcolor{ours} SurpMark 
% & 94.7 \com 0.6 & 62.67 \com 0.62
% & 80.92 \com 0.57 & 34 \com 0.88
% & 80.57 \com 0.36 & 47 \com 0.63 \\
% \bottomrule
% \end{tabular}
% % \label{tab:mix_domain}
% }
% \end{table*}
\vspace{-2mm}
\begin{table}[t]
\centering
\caption{Mixed-generator results on the DetectRL benchmark. Results are grouped by domain. For each domain, the test set contains samples generated by multiple source models, including Llama-2-70B, Claude-instant, GPT-3.5-turbo, and PaLM-2-Bison. We report AUROC and TPR@FPR=5\% with standard deviation.}
\label{tab:mix_generator}
{\scriptsize
\begin{tabular}{llcc}
\toprule
Domain & Method & AUROC & TPR@FPR=5\% \\
\midrule
\multirow{3}{*}{arXiv}
& Fast-DetectGPT & 71.24 \com 1.03 & 44 \com 1.27 \\
& Lastde++ & 71.29 \com 0.59 & 48.67 \com 1.12 \\

& \cellcolor{ours}SurpMark & \cellcolor{ours} 86.69 \com 0.42 & \cellcolor{ours} 50 \com 0.2 \\
\midrule
\multirow{3}{*}{WritingPrompts}
& Fast-DetectGPT & 70.2 \com 0.84 & 44 \com 0.11 \\
& Lastde++ & 55.99 \com 0.44 & 34.67 \com 0.38 \\

& \cellcolor{ours}SurpMark & \cellcolor{ours}81.32 \com 0.68 & \cellcolor{ours}56 \com 0.47 \\
\midrule
\multirow{3}{*}{Yelp\_review}
& Fast-DetectGPT & 71.43 \com 0.62 & 56.67 \com 1.16 \\
& Lastde++ & 56.29 \com 0.56 & 40.67 \com 0.44 \\

& \cellcolor{ours}SurpMark & \cellcolor{ours}90.18 \com 0.67 & \cellcolor{ours}62.67 \com 0.87 \\
\bottomrule
\end{tabular}}
\end{table}

\paragraph{Mixed Domains.}
We further evaluate the robustness of different detectors under two mixed evaluation settings. 
Table~\ref{tab:mix_domain} presents the mixed-domain evaluation, where results are grouped by source model and each test set contains samples from multiple domains. 
This setting measures whether a detector can maintain stable performance when the textual domain varies while the source generator is fixed. 
SurpMark achieves the highest AUROC across all source models, with particularly strong gains on Llama-2-70B and Claude-instant. 
Notably, for Claude-instant, Fast-DetectGPT and Lastde++ obtain AUROC scores below 50, suggesting that their detection signals are unreliable under mixed-domain inputs, whereas SurpMark remains substantially more effective.
\vspace{-2mm}
\paragraph{Mixed Generators.}
Table~\ref{tab:mix_generator} reports the mixed-generator evaluation, where test samples are grouped by domain and each group contains outputs from multiple source generators. 
This setting evaluates whether a detector can maintain stable performance when the domain is fixed but the generator source varies. 
SurpMark achieves the strongest overall performance, obtaining the highest AUROC on multiple domains. 
These results indicate that SurpMark is more robust to heterogeneous generation sources under mixed-generator inputs.
% \paragraph{Mixed Domains.}






\vspace{-2mm}
\subsection{Out-of-Domain (OOD) Generalization}
\vspace{-2mm}
\paragraph{Domain Shift OOD.} We further evaluate SurpMark under a reference-domain shift setting on the DetectRL benchmark, where the reference texts are drawn from a different domain from the target test set. 
The machine-generated texts are sampled from multiple frontier generators, including Llama-2-70B, Claude-Instant, and GPT-3.5-Turbo, resulting in a mixed-generator and cross-reference evaluation. 
As shown in Table~\ref{table:ref_transfer}, SurpMark remains robust under out-of-domain references and consistently outperforms Fast-DetectGPT and Lastde++. 
On average, SurpMark improves over the strongest baseline by 13.26 AUROC and 11.06 TPR@5\%FPR in genuine cross-reference settings, indicating that it captures transferable detection signals rather than relying on domain-specific reference artifacts.
\begin{table}[t]
\centering
\caption{Detection performance on three DetectRL datasets under different reference settings. Results are reported using AUROC and TPR@5\%FPR. For SurpMark, self-ref uses references from the same target dataset, while other settings use references from another dataset/domain. WP denotes WritingPrompts.}
\label{table:ref_transfer}
{\scriptsize
\begin{tabular}{llcc}
\toprule
Dataset & Method & AUROC & TPR@5\%FPR \\
\midrule
\multirow{5}{*}{Yelp}
& Fast-DetectGPT & 71.43 & 56.67 \\
& Lastde++ & 56.29 & 40.67 \\
& SurpMark self-ref & 90.18 & 62.67 \\
& SurpMark WP-ref & 88.31 & 66.00 \\
& SurpMark arXiv-ref & 85.48 & 66.67 \\
\midrule
\multirow{5}{*}{WP}
& Fast-DetectGPT & 70.20 & 44.00 \\
& Lastde++ & 55.99 & 34.67 \\
& SurpMark self-ref & 81.32 & 56.00 \\
& SurpMark arXiv-ref & 76.48 & 50.00 \\
& SurpMark Yelp-ref & 82.17 & 57.33 \\
\midrule
\multirow{5}{*}{arXiv}
& Fast-DetectGPT & 71.24 & 44.00 \\
& Lastde++ & 71.29 & 40.67 \\
& SurpMark self-ref & 86.69 & 50.00 \\
& SurpMark WP-ref & 86.40 & 55.67 \\
& SurpMark Yelp-ref & 86.54 & 60.00 \\
\bottomrule
\end{tabular}}
\end{table}




% \begin{table*}[t]
% \centering
% \caption{Detection performance on three DetectRL datasets under different reference settings. Results are reported on Yelp-review, WritingPrompts, and arXiv using AUROC and TPR@5\%FPR. For SurpMark, self-ref denotes using human-written reference texts from the same target dataset as the test samples, while WP-as-ref, arxiv-as-ref, and yelp-as-ref denote using reference texts from a different DetectRL dataset/domain.}


% \label{table:ref_transfer}
% {\scriptsize
% \begin{tabular}{lcc|cc|cc}
% \toprule
% \multirow{2}{*}{Method}
% & \multicolumn{2}{c}{Yelp-review}
% & \multicolumn{2}{c}{WritingPrompts}
% & \multicolumn{2}{c}{arXiv} \\
% \cmidrule(lr){2-3} \cmidrule(lr){4-5} \cmidrule(lr){6-7}
% & AUROC & TPR@5\%FPR
% & AUROC & TPR@5\%FPR
% & AUROC & TPR@5\%FPR \\
% \midrule
% Fast-DetectGPT 
% & 71.43 & 56.67
% & 70.20 & 44.00
% & 71.24 & 44.00 \\

% Lastde++ 
% & 56.29 & 40.67
% & 55.99 & 34.67
% & 71.29 & 40.67 \\

% SurpMark self-ref
% & 90.18 & 62.67
% & 81.32 & 56.00
% & 86.69 & 50.00 \\

% SurpMark WP-as-ref
% & 88.31 & 66.00
% & -- & --
% & 86.40 & 55.67 \\

% SurpMark arxiv-as-ref
% & 85.48 & 66.67
% & 76.48 & 50.00
% & -- & -- \\

% SurpMark yelp-as-ref
% & -- & --
% & 82.17 & 57.33
% & 86.54 & 60.00 \\
% \bottomrule
% \end{tabular}}
% \end{table*}
\vspace{-2mm}
\paragraph{Cross-Source Corruption OOD.}
We also evaluate robustness under cross-source corruption, where the test samples are partially corrupted at the sentence level with content from the opposite source. 
As shown in Table~\ref{table:cross_source_corruption}, SurpMark achieves the best performance across all datasets,  demonstrating stronger robustness when human-written and machine-generated signals are locally mixed within the same text.

\begin{table}[t]
\centering
\scriptsize
\caption{Results under cross-source corruption on the DetectRL benchmark. The reference sets contain pure human-written and LLM-generated texts from DetectRL, while the test sets are partially corrupted at the sentence level with the opposite source. WP denotes WritingPrompts.}
\label{table:cross_source_corruption}
\begin{tabular}{llcc}
\toprule
Dataset & Method & AUROC & TPR@FPR=5\% \\
\midrule
\multirow{3}{*}{arXiv}
& Fast-DetectGPT & 68.53 & 41.33 \\
& Lastde++ & 85.46 & 41.33 \\
& \cellcolor{ours}SurpMark & \cellcolor{ours}\textbf{93.86} & \cellcolor{ours} \textbf{74.00} \\
\midrule
\multirow{3}{*}{WP}
& Fast-DetectGPT & 60.17 & 34.00 \\
& Lastde++ & 74.40 & 34.67 \\
& \cellcolor{ours} SurpMark & \cellcolor{ours} \textbf{84.15} & \cellcolor{ours} \textbf{50.00} \\
\midrule
\multirow{3}{*}{Yelp}
& Fast-DetectGPT & 69.93 & 46.00 \\
& Lastde++ & 70.60 & 44.67 \\
& \cellcolor{ours}SurpMark & \cellcolor{ours}\textbf{88.29} & \cellcolor{ours}\textbf{58.00} \\
\bottomrule
\end{tabular}
\end{table}
\vspace{-2mm}
\paragraph{Paraphrase OOD. }
We consider a paraphrase-based OOD setting in which the reference sets contain pure human-written and LLM-generated texts, while the test sets contain paraphrased versions of those texts. This creates an OOD shift by altering surface form while largely preserving meaning and source identity. We consider three paraphrasing settings: polish uses LM-based fluency rewriting, Back-translation uses round-trip translation, DIPPER uses a dedicated paraphrasing model to generate semantically equivalent rewrites with stronger lexical and syntactic variation. Notably, SurpMark maintains near-perfect AUROC and TPR@FPR=5\%, indicating strong robustness to lexical and syntactic variations introduced by paraphrasing.
\begin{table}[t]
\centering
\scriptsize
\caption{Performance under paraphrase-induced OOD shift on DetectRL. The reference sets contain the original source-pure human-written and LLM-generated texts, while the test sets contain their paraphrased versions. We consider three paraphrase settings: Polish, DIPPER, and Back-translation.}
\label{table:paraphrase_ood}
\begin{tabular}{llcc}
\toprule
Setting & Method & AUROC & TPR@FPR=5\% \\
\midrule
\multirow{3}{*}{Polish}
& Fast-DetectGPT & 83.47 & 42.00 \\
& Lastde++ & 82.37 & 43.33 \\
& \cellcolor{ours}SurpMark & \cellcolor{ours}\textbf{99.21} & \cellcolor{ours}\textbf{99.33} \\
\midrule
\multirow{3}{*}{DIPPER}
& Fast-DetectGPT & 94.70 & 85.33 \\
& Lastde++ & 92.82 & 84.00 \\
& \cellcolor{ours}SurpMark & \cellcolor{ours}\textbf{99.45} & \cellcolor{ours}\textbf{99.33} \\
\midrule
\multirow{3}{*}{Back-trans.}
& Fast-DetectGPT & 97.69 & 84.00 \\
& Lastde++ & 94.31 & 70.00 \\
& \cellcolor{ours}SurpMark & \cellcolor{ours}\textbf{99.79} & \cellcolor{ours}\textbf{99.33} \\
\bottomrule
\end{tabular}
\end{table}
% \begin{table}[htbp]
% \centering
% {\tiny
% \begin{tabular}{l c c c c}
% \toprule
% Test & self-ref & WritingPrompts-as-ref & XSum-as-ref & SQuAD-as-ref \\
% \midrule
% XSum@Llama2-13b              & 97.09 & 97.22 & --   & 96.45 \\
% WritingPrompts@Llama2-13b    & 99.53 & --    & 99.71 & 99.33 \\
% SQuAD@Llama2-13b             & 96.13 & 94.36 & 95.67 & -- \\
% XSum@Llama3-8b               & 97.09 & 97.25 & --    & 95.13 \\
% WritingPrompts@Llama3-8b     & 99.86 & --    & 99.95 & 99.88 \\
% SQuAD@Llama3-8b              & 94.17 & 93.77 & 93.90 & -- \\
% XSum@Gemma\_7b               & 72.97 & 68.32 & --    & 73.81 \\
% WritingPrompts@Gemma\_7b     & 89.42 & --    & 86.60 & 90.00 \\
% SQuAD@Gemma\_7b              & 68.27 & 72.92 & 70.44 & -- \\
% XSum@GPT-Neox-20b            & 82.45 & 82.52 & --    & 81.80 \\
% WritingPrompts@GPT-Neox-20b  & 94.31 & --    & 93.45 & 92.48 \\
% SQuAD@GPT-Neox-20b           & 82.45 & 81.43 & 83.02 & -- \\
% \bottomrule
% \end{tabular}}
% \caption{AUROC of SurpMark under different reference choices across datasets and models.}
% \label{tab:ref-choice2}
% \end{table}

% \begin{table}[t]
% \centering
% \setlength{\tabcolsep}{3pt}      % 默认约 6pt，缩小列间距
% \renewcommand{\arraystretch}{1.0}
% {\tiny
% \begin{tabular}{l l c c c c}
% \toprule
% Dataset & Model & self-ref & WP-as-ref & XSum-as-ref & SQuAD-as-ref \\
% \midrule
% XSum           & Llama2-13B      & 97.09 & 97.22 & --    & 96.45 \\
% WritingPrompts & Llama2-13B      & 99.53 & --    & 99.71 & 99.33 \\
% SQuAD          & Llama2-13B      & 96.13 & 94.36 & 95.67 & --    \\
% XSum           & Llama3-8B      & 97.09 & 97.25 & --    & 95.13 \\
% WritingPrompts & Llama3-8B       & 99.86 & --    & 99.95 & 99.88 \\
% SQuAD          & Llama3-8B       & 94.17 & 93.77 & 93.90 & --    \\
% XSum           & Gemma\_7B       & 72.97 & 68.32 & --    & 73.81 \\
% WritingPrompts & Gemma\_7B       & 89.42 & --    & 86.60 & 90.00 \\
% SQuAD          & Gemma\_7B       & 68.27 & 72.92 & 70.44 & --    \\
% XSum           & GPT-Neox-20B    & 82.45 & 82.52 & --    & 81.80 \\
% WritingPrompts & GPT-Neox-20B    & 94.31 & --    & 93.45 & 92.48 \\
% SQuAD          & GPT-Neox-20B    & 82.45 & 81.43 & 83.02 & --    \\
% \bottomrule
% \end{tabular}}
% \caption{AUROC of SurpMark under different reference choices across datasets and models.}
% \label{tab:ref-choice2}
% \end{table}
% \begin{figure*}[!t]
% \centering
% \includegraphics[width=0.9\textwidth]{exp-length.pdf}
% \caption{(a) AUROC vs. number of reference samples. The blue curve (“$k$-optimized”) picks the best $k$ at each number of reference. orange/green curves fix $k \in \{7,8\}$. (b) AUROC vs. test length $n$ under different reference lengths. Solid lines are k-optimized for each reference sample truncated to 50/100/200 tokens; shaded bands show the attainable range across $k$ at each $n$. (c) Detection results of 7 detection methods on 6 test lengths.}
% \label{fig:k-optmized-k7-k8}
% \end{figure*}
% \begin{figure*}[!t]
% \centering
% \includegraphics[width=1 \textwidth]{contour_proxy_model.pdf}
% \caption{(a-b) AUROC contour maps (WritingPrompts@Gemma-7B). Left: $k=7$; right: $k=8$. The x-axis is reference length (tokens) and the y-axis is test length (tokens). Colors encode AUROC. In both panels, contours tilt up-right, indicating a trade-off: larger reference length allows smaller test length at similar performance. (c) AUROC vs. proxy model.}
% \label{fig:contour-proxy}
% \end{figure*}

% \begin{table}[htbp]
% \centering
% {\tiny
% \begin{tabular}{l l c c c}
% \toprule
% Dataset@Source model      & Method              & $k=7$ & $k=8$ & Best \\
% \midrule
%          & k-means            & 80.42 & 79.32 & \textbf{81.79} \\
% XSum@GPT-4.1-mini          & equal-width  & 80.41 & 76.20  & 80.41 \\
%         & equal-mass   & 71.40 & 74.22 & 74.22 \\
% \midrule
%   & k-means             & 82.05 & 84.86 & \textbf{84.86} \\
% WritingPrompts@GPT5-chat  & equal-width & 62.90 & 72.06  & 72.06 \\
%   & equal-mass    & 84.42 & 83.59 & 84.42 \\
% \bottomrule
% \end{tabular}
% \caption{AUROC of different discretization schemes under varying number of states $k$.}
% \label{tab:discretization-k}}
% \end{table}
% \vspace{-2mm}
% \paragraph{Necessity of $k$-means.}
% \vspace{-2mm}
% In Table~\ref{tab:discretization-k}, we evaluate the effect of different discretization schemes, including $k$-means, equal-width, and equal-mass binning. Across all datasets and $k$ values, $k$-means is the most robust quantization scheme: it consistently reaches or matches the best AUROC, while equal-mass can degrade sharply on XSum@GPT-4.1-mini and equal-width is unstable and often much worse.

% \vspace{-3mm}
% \paragraph{Cross-Generator Generalization.} 
% To evaluate robustness without recomputing transitions per model, we propose SurpMark-MC, a multi-class variant. We first pool all surprisal sequences from the human reference texts and from several LM reference texts and run one k-means over this pooled set to obtain a shared set of discretization bins. Using this common quantizer, we then estimate one first-order Markov chain for the human class and one machine Markov chain for each generator. We then classify test texts by assigning them to the source with the minimal GJS divergence (i.e., a nearest-neighbor rule). In a multi-class setting with GPT-J-6B, GPT-4.1-mini, and Llama-13B, SurpMark-MC achieves 82.3\% overall accuracy, with 78.7\% accuracy for human texts, 78.0\% for GPT-J-6B, 84.0\% for GPT-4.1-mini, and 96.0\% for Llama-13B.  In Table~\ref{tab:wp-generators}, we further collapse the multi-generator setting into a binary classification task (LM-generated vs.\ human-written), defining the score as the difference between the GJS to the human reference and the minimum GJS to any machine reference. As shown in the bottom of Table~\ref{tab:wp-generators}, this approach generalizes effectively even to unseen models (e.g., Llama-3-8B, GPT-5-chat, GPT-neo-20B), matching or surpassing baselines like Lastde++ and Fast-DetectGPT.

% \begin{table}[htbp]
% \centering
% {\tiny
% \begin{tabular}{l c c c}
% \toprule
% WritingPrompts      & LLaMA2-13B & GPT-4.1-mini & GPT-J-6B & LLaMA3-8B & GPT-5-chat & GPT-neo-20B \\
% \midrule
% Lastde++            & 99.14      & 68.49        & \textbf{95.96} & 99.56      & 30.64      & 92.68      \\
% Fast-DetectGPT      & 99.56      & 70.23        & 93.80          & 99.84      & 30.01      & 92.22      \\
% SurpMark self-ref   & 99.59      & 87.27        & 96.85          & 99.87      & 83.56      & 93.93      \\
% SurpMark-MC         & \textbf{99.72} & \textbf{90.01} & 95.61      & \textbf{99.77} & \textbf{64.27} & \textbf{92.84} \\
% \bottomrule
% \end{tabular}}
% \caption{AUROC on WritingPrompts when the reference transition probabilities (LLaMA2-13B, GPT-4.1-mini and GPT-J-6B) are held fixed and not re-estimated for each test-time generator}
% \label{tab:wp-generators}
% \end{table}



% \begin{table}[htbp]
% \centering
% \resizebox{\linewidth}{!}{
% \begin{tabular}{l c c c}
% \toprule
% \multicolumn{4}{c}{\textbf{Blind Detection on Seen Generators}} \\
% \cmidrule(lr){1-4}
% WritingPrompts & Llama2-13B & GPT-4.1-mini & GPT-J-6B \\
% \midrule
% Lastde++ & 99.14 & 68.49 & {95.96} \\
% Fast-DetectGPT & 99.56 & 70.23 & 93.80 \\
% SurpMark-MC & {99.72} & {90.01} & 95.61 \\
% \midrule
% \multicolumn{4}{c}{\textbf{Zero-shot Detection on Unseen Generators}} \\
% \cmidrule(lr){1-4}
% WritingPrompts & Llama3-8B & GPT-5-chat & GPT-neo-20B \\
% \midrule
% Lastde++ & 99.56 & 30.64 & 92.68 \\
% Fast-DetectGPT & 99.84 & 30.01 & 92.22 \\
% SurpMark-MC & {99.77} &{64.27} & {92.84} \\
% \bottomrule
% \end{tabular}
% }
% \caption{\textbf{Generator-agnostic detection capabilities.} We evaluate using a \textbf{single fixed reference} constructed from the pooled training set. In the top setting, the test generator belongs to the pool but its specific identity is \textbf{not provided} to the detector. In the bottom setting, the test generator is completely unseen. SurpMark-MC consistently outperforms or matches baselines without requiring source-specific calibration.}
% \label{tab:wp-generators}
% \end{table}

% SurpMark-MC extends SurpMark to multi-class attribution using a shared quantizer and per-source Markov chains, assigning each test text to the source with minimum GJS; it attains strong multi-class accuracy and, when collapsed to binary detection via a human-vs-best-machine GJS. See Appendix~\ref{app:cross-generator} (Table~\ref{tab:wp-generators}) for algorithm details and results.

\vspace{-3mm}
\paragraph{More Results.} We provide additional experimental results in the Appendix, including: (1) TPR at low FPR (Appendix~\ref{sec:tpr}) (2) ablation study on decoding strategies (Appendix~\ref{app:decoding}) (3) evaluations under paraphrasing attack (Appendix~\ref{sec:paraphrasing-attack}) (4) evaluations under prompt-engineered adversarial attacks (Appendix~\ref{sec:prompts-eng-attack}) (5) ablation on the necessity of first-order markov chain (Appendix~\ref{sec:ablation-fo-markov}) (6) ablation on the necessity of GJS distance (Appendix~\ref{sec:ablation-GJS-l1-l2}) (7) discussion for threhold $\tau$ selection (Appendix~\ref{app:threshold}) (8) robustness against character-word perturbation (Appendix~\ref{app:cw-perturbation-ood}). 
\vspace{-2mm}
\section{Conclusion}
\vspace{-2mm}
We presented SurpMark, a reference-based detector for black-box detection
of machine-generated text. By quantizing token surprisals into interpretable states and modeling their dynamics as a Markov chain, SurpMark reduces each passage to a transition matrix and scores it via a GJS score against fixed human/machine references. Theoretically, we derive design guidance for how the discretization bins should scale with data and provide a principled justification for our test statistic. Empirically, across diverse datasets, source models, and scenarios, SurpMark consistently matches or surpasses strong baselines.
% \vspace{-3mm}
% \paragraph{Limitations.} 
% Our analysis does not claim that human or LLM text generation is well-approximated by a first-order Markov process at the token level. Rather, it establishes that, for the discretized proxy-surprisal state representation used by our detector, the incremental benefit of moving beyond first order is empirically tiny, while higher-order estimation variance is substantial under finite reference/test budgets. Extending the analysis to richer state representations and longer-range dependence models is an important direction for future work.
\vspace{-2mm}
\paragraph{Limitations.}
SurpMark does not assume token-level first-order Markov generation; the approximation is only applied to discretized proxy-surprisal states. 
Higher-order variants bring limited gains but require more samples, so richer state and dependency modeling is left for future work. 

Another limitation is that human texts seen during pretraining may have smoother likelihood trajectories \cite{shi2024detecting}, making them more LLM-like in our feature space and potentially increasing false positive rates. 
We provide additional analysis of this effect in Appendix~\ref{app:pretraining_contamination}.


\section*{Impact Statement}
SurpMark provides a probabilistic detection signal, not definitive evidence that a text is AI-generated. 
High AUROC does not imply perfect reliability on individual cases. 
Thus, its outputs should not be used as the sole basis for high-stakes decisions such as plagiarism accusations or academic penalties, and should be accompanied by human review and contextual evidence.
\section*{Acknowledgments}
Resources used in preparing this research were provided, in part, by the Province of Ontario, the Government of Canada through CIFAR, and companies sponsoring the Vector Institute www.vectorinstitute.ai/partnerships/.


\bibliography{example_paper}
\bibliographystyle{icml2026}

%%%%%%%%%%%%%%%%%%%%%%%%%%%%%%%%%%%%%%%%%%%%%%%%%%%%%%%%%%%%%%%%%%%%%%%%%%%%%%%
%%%%%%%%%%%%%%%%%%%%%%%%%%%%%%%%%%%%%%%%%%%%%%%%%%%%%%%%%%%%%%%%%%%%%%%%%%%%%%%
% APPENDIX
%%%%%%%%%%%%%%%%%%%%%%%%%%%%%%%%%%%%%%%%%%%%%%%%%%%%%%%%%%%%%%%%%%%%%%%%%%%%%%%
%%%%%%%%%%%%%%%%%%%%%%%%%%%%%%%%%%%%%%%%%%%%%%%%%%%%%%%%%%%%%%%%%%%%%%%%%%%%%%%
\newpage
\appendix
\onecolumn
\startcontents[appendices] % 开启名为 appendices 的局部目录
\section*{Table of Contents} % 手写一个标题

% 打印局部目录
% 参数说明: [标记名] {前缀} {起始深度} {格式设置}
\printcontents[appendices]{}{1}{\setcounter{tocdepth}{3}}

% 先打印“上一次编译”收集到的 appendix 目录

\newpage



\section{Algorithm}\label{sec:alg}
\begin{algorithm}[h]
\caption{SurpMark (Offline): Build Human/Machine Reference Transitions}
\label{alg:surpmark-offline}
\begin{algorithmic}[1]
\REQUIRE Proxy LM $F_\theta$; human corpus $\mathcal{D}_Q$; machine/LLM corpus $\mathcal{D}_P$; number of bins $k$
\ENSURE Shared surprisal quantizer $q_k$; reference transition matrices $\hat M_P, \hat M_Q$; total reference length $N$

\STATE \textbf{Score references.}
For every $t \in \mathcal{D}_Q \cup \mathcal{D}_P$, run $F_\theta$ to obtain token sequence $x_{1:N}$ and surprisals $s_{1:N}$, where
$s_t \leftarrow -\log p_\theta(x_t\mid x_{1:t-1})$.

\STATE \textbf{Fit shared quantizer.}
Pool all reference surprisals and fit $k$-means to obtain $q_k:\mathbb{R}\to\{1,\dots,k\}$.

\STATE \textbf{Discretize to states.}
Set $a_t \leftarrow q_k(s_t)$ for $t=1,\dots,N$.

\STATE \textbf{Estimate transitions.}
For each corpus $C\in\{P,Q\}$, estimate the empirical first-order transition matrix $\hat M_C$ by counts:
\STATE \quad $\displaystyle
\hat M_C(j\mid i)=\frac{\sum_{t=1}^{N-1}\mathbf{1}\{a_t=i,\,a_{t+1}=j\}}{\sum_{t=1}^{N-1}\mathbf{1}\{a_t=i\}},\quad i,j\in\{1,\dots,k\}.
$

\STATE \textbf{Record length.}
Let $N$ be the total length of \emph{reference} used to form $\hat M_P$ and $\hat M_Q$ (sum over sequences).

\STATE \textbf{return} $q_k,\ \hat M_P,\ \hat M_Q,\ N$.
\end{algorithmic}
\end{algorithm}
\begin{algorithm}[h]
\caption{SurpMark (Online): Decision via GJS score against References}
\label{alg:surpmark-online}
\begin{algorithmic}[1]
\REQUIRE Proxy LM $F_\theta$; test text $t$; shared quantizer $q_k$; reference transitions $\hat M_P,\hat M_Q$; reference length $N$
\ENSURE Score $\Delta\mathrm{GJS}_n$ and label $\Omega\in\{\textsc{Machine},\textsc{Human}\}$

\STATE \textbf{Score test text.}
Run $F_\theta$ on $t$ to get tokens $x_{1:n}$ and surprisals $s_{1:n}$.

\STATE \textbf{Discretize.}
Map to surprisal states $a_t \leftarrow q_k(s_t)$ for $t=1,\dots,n$ and estimate the test transition matrix $\hat M_T$ using the same formula as Offline.

\STATE \textbf{Set mixing weight.}
$\alpha \leftarrow N/n$.

\STATE \textbf{Compute divergence.}
\STATE \quad $\displaystyle
\Delta\mathrm{GJS}_n
= \mathrm{GJS}(\hat M_P,\hat M_T,\alpha)
- \mathrm{GJS}(\hat M_Q,\hat M_T,\alpha).
$

\STATE \textbf{Decision rule.}
\STATE \quad $\displaystyle
\Omega =
\begin{cases}
\textsc{Machine}, & \Delta\mathrm{GJS}_n \le \tau,\\
\textsc{Human}, & \Delta\mathrm{GJS}_n > \tau.
\end{cases}
$

\STATE \textbf{return} $\Delta\mathrm{GJS}_n,\ \Omega$.
\end{algorithmic}
\end{algorithm}




\section{Justification for First-Order Modeling} \label{sec:empirical-finding-for-FO}
\subsection{Empirical Findings}
\begin{table}[h]
\centering
\begin{minipage}{0.48\textwidth} 
\centering
\small
\begin{tabular}{ccccc}
\toprule
 & \multicolumn{4}{c}{Reference size} \\
Order & 300 & 600 & 900 & 1200 \\
\midrule
1 & 86.51 & 86.17 & 86.69 & 86.89 \\
2 & 81.49 & 82.08 & 82.77 & 82.80 \\
3 & 74.98 & 74.74 & 75.84 & 75.85 \\
4 & 64.72 & 69.51 & 71.20 & 73.00 \\
\bottomrule
\end{tabular} 
\caption*{(a) AUROC vs. reference size}
\end{minipage}
\hfill
\begin{minipage}{0.48\textwidth}
\centering
\small
\begin{tabular}{ccccc}
\toprule
 & \multicolumn{4}{c}{Test length} \\
Order & 150 & 200 & 250 & 300 \\
\midrule
1 & 90.59 & 92.66 & 94.60 & 94.58 \\
2 & 89.42 & 90.67 & 92.21 & 92.60 \\
3 & 88.67 & 90.57 & 91.40 & 91.40 \\
4 & 62.29 & 65.59 & 68.46 & 68.33 \\
\bottomrule
\end{tabular}
\caption*{(b) AUROC vs. test length}
\end{minipage}

\caption{Effect of reference size and test length on AUROC for different Markov orders.}

\label{tab:aurox-vs-test}
\end{table}
All AUROC values in Figure~\ref{fig:transition-heatmap-distribution}(b) are computed using exactly the same amount of reference data and test data. As in our main experiments, we use 300 human paragraphs and 300 machine-generated paragraphs, each with length about 100-200 tokens as reference data. We use 150 human paragraphs and 150 machine paragraphs as test data. Intuitively, increasing the Markov order makes the state space explode while the amount of reference data is fixed, so transition estimates become very sparse and noisy.

The degradation with larger order is a sparsity effect that arises both on the reference side and on the test side: (i) state-space explosion: A first-order chain with $k$ bins has $k^2$ transitions; an order-$v$ chain effectively has $k^{v+1}$ transitions. With only 300 human + 300 machine paragraphs of 100-200 tokens, many high-order contexts in the reference data are observed only a few times or not at all, so the estimated transitions become extremely noisy. (ii) Short test sequences. Each test paragraph is itself only 100-200 tokens long. Even if the reference transitions were perfectly estimated, an order-$v$ model on a 100-200-token sequence can observe only a very small number of distinct 
$v$-length contexts. The higher-order model is severely under-sampled on each individual test example.

To isolate the effect of reference sparsity, in the XSum@GPT-J-6B dataset, we fixed $k=6$ and the test set (150 human + 150 machine paragraphs, 100-200 tokens each), and increased the reference size from 300 to 1200 paragraphs per side. As shown in Table~\ref{tab:aurox-vs-test}(a), AUROC for higher-order models improves only slightly and remains clearly below the first-order model.

To further evaluate the effect of test sparsity, in another WritingPrompts@Genmma-7B dataset, we vary the test passage length from 150 to 300 tokens while keeping the reference size fixed (300 passages per side, each with fixed 300 tokens). As in Table~\ref{tab:aurox-vs-test}(b), AUROC consistently increases for all orders, but the first-order model remains clearly best, and higher orders still lag behind by several points. 

Taken together, these ablations reflect practical text-detection settings with limited reference data and short passages. In this regime, the first-order model offers the best bias-variance tradeoff, so we believe it is the most reasonable default choice.

\subsection{Choice of First-Order Markov Modeling} \label{sec:theo-justification}
We clarify our reasoning by (i) starting from Gray's Markov approximation theory, (ii) explaining how the gain from order
$K$ to $K+1$ is governed by conditional mutual information, (iii) mapping this theory to our discretized surprisal process, and (iv) presenting empirical measurements showing that the additional benefit of a second-order approximation over a first-order one is very small.

(1) Best finite-order Markov approximation in Gray's theory. 

Following Gray's Entropy and Information Theory [\cite{gray2011entropy}, Sec. 6.4, Cor. 6.4.1–6.4.2; Sec. 7.4, Cor. 7.4.2–7.4.3], consider a stationary discrete-time source $\{X_n\}$. Gray constructs, for each order $K$, a canonical $K$-th order Markov chain $M_K$ whose conditional distributions match those of the source given the last $K$ symbols. He shows that $M_K$ is optimal in the sense that it uniquely minimizes the relative entropy rate between the true source and any $K$-th order Markov chain on the same alphabet. In other words, the family of finite-order Markov chains $\{M_K\}$ provides a sequence of best approximations to the stationary process in the relative-entropy-rate sense. Formally,

\begin{align}
H_{p\|p^K} (\{X_n\}) = \inf_{M_K\in \mathcal{M}_K} H_{p\|M_K}(\{X_n\}) = I(X_0;X_{-\infty}^{-K-1}|X_{-K}^{-1})
\end{align}

where $p$ is the true stationary source, $p^K$ is the canonical $K$-th order Markov approximation to $p$, $\mathcal{M}_K$ is the class of stationary $K$-th order Markov sources on the same alphabet, $H_{p\|q} (\{X_n\})$ is the relative entropy rate of $p$ with respect to $q$, $X_{-\infty}^{-K-1}=(X_{-\infty}\dots,X_{-K-1})$ is the infinite past, $X_{-K}^{-1}=(X_{-K},\dots,X_{-1})$ is the block of the last $K$ symbols, and $I(\cdot;\cdot|\cdot)$ conditional mutual information. Applying the above indentity with $K+1$ instead of $K$, we get
\begin{align}
H_{p\|p^{K+1}}(\{X_n\}) = I(X_0;X_{-\infty}^{-K-2}|X_{-K-1}^{-1})
\end{align}

We are interested in the gain from going from order $K$ to $K+1$, so
\begin{align}
\Delta_K&= H_{p\|p^K} (\{X_n\})-H_{p\|p^{K+1}}(\{X_n\}) \\
& = I(X_0;X_{-\infty}^{-K-1}|X_{-K}^{-1}) - I(X_0;X_{-\infty}^{-K-2}|X_{-K-1}^{-1}) \\
& = I(X_0;X_{-K-1}|X_{-K}^{-1})
\end{align}

(2) Mapping to our discretized surprisal process

In our setting, $\{X_n\}$ is instantiated by the discretized surprisal process $\{a_t\}$, where $a_t$ corresponds to $X_0$, $X_{-K}^{-1}$ corresponds to $a_{t-K}^{t-1}==(a_{t-K},\cdots, a_{t-1})$, $X_{-K-1}$ corresponds to $a_{t-K-1}$. For the case $K=1$, the gain from first-order to second order is precisely $I(a_t;a_{t-2}|a_{t-1})$. We directly estimate the relevant conditional mutual information term on our data. We first fit a first-order $\hat{P}_1(a_t|a_{t-1})$ and a second-order model $\hat{P}_2(a_t|a_{t-1},a_{t-2})$ from transition counts on the reference set. We then compute plug-in estimates on test set
\begin{align}
\hat{H}_1 &= -\frac{1}{n-1}\sum_{t=2}^{n}\log_2 \hat{P}_1(a_t|a_{t-1}) \\
\hat{H}_2 &= -\frac{1}{n-2}\sum_{t=3}^{n}\log_2 \hat{P}_2(a_t|a_{t-1},a_{t-2})
\end{align}
Their difference is the plug-in estimate of the conditional mutual information $\hat{I}=\hat{H}_1-\hat{H}_2=I(a_t;a_{t-2}|a_{t-1})$, in bits per token, i.e., the extra predictive information contributed by the second-order context beyond the immediate past. On our data, we obtain empirical estimates of conditional mutual information and perplexity in Table~\ref{tab:cond-entropy-perplexity}. In our experiments, $\hat I$ is at most 0.0076 bits/token, which corresponds to a perplexity reduction around 0.5\%.  Thus, in terms of average predictive performance, the second-order Markov model brings only a sub-percent improvement over the first-order model. Combined with Gray's best Markov approximation theory, this indicates that a first-order Markov chain already captures most useful temporal dependence in the discretized surprisal dynamics, and provides a theoretically justified and empirically sufficient model for the discretized surprisal dynamics in our detector.

\begin{table}[htbp]
\centering
{\scriptsize
\begin{tabular}{l l c c c c}
\toprule
Source      & Order pair      & $\hat H_k$ (bits/token) & $\hat H_1 - \hat H_2$ (bits/token) & Perplexity & Rel.\ PP change vs.\ 1st \\
\midrule
GPT-5-chat  & 1st (baseline)  & 2.7882                  & 0.0000                              & 6.9075     & 0.000\%                 \\
GPT-5-chat  & 2nd order       & 2.7805                  & 0.0076                              & 6.8711     & +0.528\%                \\
Human       & 1st (baseline)  & 2.8089                  & 0.0000                              & 7.0074     & 0.000\%                 \\
Human       & 2nd order       & 2.8043                  & 0.0045                              & 6.9854     & +0.314\%                \\
\bottomrule
\end{tabular}}
\caption{Conditional entropies and perplexities for discretized surprisal states.}
\label{tab:cond-entropy-perplexity}
\end{table}

\section{Related Work}\label{app:related-work}
Prior work on text detection can be broadly categorized into classifier-based and statistics-based methods. Classifier-based detectors train task-specific classifiers to distinguish between human-written and machine-generated text \citep{guo2023closechatgpthumanexperts,GPTZero,guo2024detective}. While effective with sufficient training data, they are costly to build and must be retrained whenever the domain or generator shifts.

Statistics-based approaches can be divided into two groups based on their design of decision statistics. The first global-statistic methods rely on overall features of the text such as likelihood \citep{solaiman2019releasestrategiessocialimpacts}, LogRank \citep{solaiman2019releasestrategiessocialimpacts} that measures the log of each token’s rank in a model’s predicted distribution , or entropy \citep{gehrmann-etal-2019-gltr} that measures the uncertainty of a model’s next-token distribution. Binoculars \citep{hans2024spotting} enhances this by computing the contrastive cross-perplexity between two different LLMs, identifying machine-generated text through the relative "surprisal" captured by an observer model. Distributional-statistic methods generate a neighborhood around the test passage via perturbation, continuation, or sampling, and then measure divergence between the test instance and this synthetic distribution. DetectGPT \citep{mitchell2023detectgptzeroshotmachinegeneratedtext} leverages the local curvature of log-probability function, comparing original passages with perturbed variants to enable detection of machine-generated text. Fast-DetectGPT \citep{bao2024fastdetectgpt} introduces conditional probability curvature for faster detection. DNA-GPT \citep{yang2023dnagptdivergentngramanalysis} truncates passages, and analyzes n-gram divergences of the regeneration. DetectLLM-NPR \citep{su2023detectllm} leverages normalized perturbed log-rank statistics, showing that machine-generated texts are more sensitive to small perturbations. Lastde++ \citep{xu2025trainingfree} combines global likelihood with local diversity entropy, where discretization of token probabilities stabilizes the entropy feature. Expanding beyond time-domain statistics, FourierGPT \citep{xu2024detectingsubtledifferenceshuman} utilizes spectral analysis to detect low-frequency structural patterns in token probability trajectories, which are characteristic of LLM outputs. In contrast, our framework discretizes token surprisals to build surprisal-state Markov transitions, enabling likelihood-free hypothesis test. Our method lies between global- and distributional-statistic approaches: it scores each text in a single pass without regeneration, yet makes comparative decisions by measuring alignment with fixed human and machine references.


Recent work has explored kernel-based statistical tests for machine-generated text detection \citep{zhang2024detecting,song2025deep}. \cite{song2025deep} introduced R-Detect, a relative test framework that reduces false positives by comparing whether a test text is closer to human-written or machine-generated distributions. Our method shares a common foundation with \cite{song2025deep} in that it can also be viewed as a relative test framework. Notably, while the decision rules of these kernel-based approaches are non-parametric and do not rely on supervised classifiers, their optimized variants require training kernel parameters on reference corpora, which increases computational cost. Our approach only requires a lightweight data discretization stage. 
\section{Theoretical Analysis}
\subsection{Problem Setup}
Let $\{s_t^P\}_{t=1}^{N}$ and $\{s_t^Q\}_{t=1}^{N}$ be the surprisal sequences produced by a fixed proxy LM on reference corpora from $P$ and $Q$. Each sequence is modeled as an ergodic first-order Markov process on $\mathbb{R}$.
For an integer $k \ge 2$, let $q_k: \mathbb{R} \rightarrow \mathcal{A}=\{1,\dots,k\}$ be a shared quantizer with boundaries $b_1<\dots <b_{k-1}$, and discretized states $a_t^P = q_k(s_t^P)$ and $a_t^Q = q_k(s_t^Q)$. Let $\mathcal{S}_P,\mathcal{S}_Q$ denote the underlying Markov transition kernels on the real-valued surprisal sequences before discretization.  The induced transition kernels on the $k$-state alphabet are $M_P(j|i)=\Pr[a_{t+1}^P=j|a_{t}^P=i]$ and likewise $M_Q$. Their plug-in estimators $\hat M_P, \hat M_Q$ are formed from transition counts with $\hat{M}_P(a|s) = \frac{N_P(s,a)}{N_P(s)}$, where $N_P(s)$ is the number of occurrences of state $s$ in $a_{1:N}^P$, and $N_P(s,a)$ is the number of times $s$ is followed by $a$; analogously for $Q$. Let $\pi_P, \pi_Q$ are stationary distributions of $M_P$ and $M_Q$, we define
$\pi_{\min}:=\min \{\min_{s\in \mathcal{A}}\pi_P(s),\min_{s\in \mathcal{A}}\pi_Q(s)\} $.

We observe an independent test surprisal-state sequence $a^T_{1:N}:=\{a_t^T\}_{t=1}^{n} \sim M_T$, where the test source $M_T$ is either $M_P$ (null $H_0$) or $M_Q$ (alternative $H_1$). All three sequences are discretized by the same $q_k$.

Throughout the analysis we impose the following conditions on the induced chains
$M_P$ and $M_Q$. These assumptions are standard in the study of Markov
concentration inequalities and are required in order to apply the auxiliary
results recalled below.
\begin{assumption}\label{assumption:markov-chain}
    We impose the following standing conditions on the induced chains $M_P, M_Q$. 
$M_P$ and $ M_Q$ are irreducible, aperiodic Markov chain on the finite alphabet $\mathcal{A}$ with unique stationary distribution $\pi_P$ and $\pi_Q$ and maximum hitting time $T(M_P)$ and $T(M_Q)$ respectively. We assume $\pi_{\min}:=\min \{\min_{s\in \mathcal{A}}\pi_P(s),\min_{s\in \mathcal{A}}\pi_Q(s)\} \gtrsim 1/k $, and $T(M_{\bullet})=O(1)$. 
\end{assumption}
\subsection{Discretization Effect}
\subsubsection{Auxiliary Results from Literature}
\paragraph{The GJS Divergence as $f$-divergence.} The GJS divergence is a specific instance of a broader class of divergences known as $f$-divergences. An $f$-divergence between two discrete probability distributions $p$ and $q$ is defined by a convex generator function $f$ where $f(1)=0$. The GJS divergence is equivalent to the $w$-skew Jensen-Shannon Divergence with $w=\alpha/(1+\alpha)$, which is an $f$-divergence generated by the function $f_{JS}^{w}(t)$.
\begin{align}
 f_{JS}^{w}(t) = \alpha t \log(\frac{t}{\alpha t +1 - \alpha})+(1-\alpha) \log (\frac{1}{\alpha t +1 - \alpha})   
\end{align}


For notational convenience, we abbreviate $f_{JS}^{\alpha}$ as $f$. This connection allows us to leverage the following theoretical tools developed for general $f$-divergences. 

\begin{assumption}[Assumption 9 in \cite{pillutla2023mauvescoresgenerativemodels}]\label{assump:3}
    We assume that the generator function $f$ of the $f$-divergence must satisfy the following three conditions:
    \begin{itemize}
        \item (A1) The function $f$ and its conjugate generator $f^*$ must be bounded at zero. Formally, $f(0)<\infty$ and $f^*(0)<\infty$.
        \item (A2) The first derivatives of $f$ and $f^*$ must not grow faster than a logarithmic function. For any $t\in (0,1)$, there must exits constants $C_1$ and $C_1^*$ such that $|f'(t)|\leq C_1 (\max(1,\log(1/t)))$ and $|(f^*)'(t)|\leq C_1^* (\max(1,\log(1/t)))$.
        \item The second derivatives of $f$ and $f^* 
$ must not grow faster than $\frac{1}{t}$ as $t\rightarrow 0$. Formally, there must exist  constants $C_2$ and $C_2^*$such that for any$t\in (0, \infty), \frac{t}{2}f''(t)\leq C_2$, and $\frac{t}{2}(f^*)''(t)\leq C_2^*$.
    \end{itemize}
\end{assumption}
\begin{lemma}[Approximate Lipschitz Property of the $f$-divergence, Lemma 20 in \cite{pillutla2023mauvescoresgenerativemodels}]\label{lemma:Lip}
     Let $f$ be a generator function satisfying Assumption~\ref{assump:3}. Consider the bivariate scalar function $\psi:[0,1] \times [0,1] \rightarrow [0,\infty)$ defined as $\psi(p, q)=q f(\frac{p}{q})$. For all probability values $p,p',q,q' \in [0,1]$ with $\max(p,p')>0$ and $\max(q,q')>0$, the following inequalities hold:
     \begin{align}
         |\psi(p',q)-\psi(p,q)|&\leq \left(C_1\max\left(1,\log\frac{1}{\max(p,p')}\right)+\max(C_0^*,C_2)\right)|p-p'| \\
         |\psi(p,q')-\psi(p,q)|&\leq \left(C_1^*\max\left(1,\log\frac{1}{\max(q,q')}\right)+\max(C_0,C_2^*)\right)|q-q'|
     \end{align}
\end{lemma}
\begin{assumption}[Assumption 3(b) in \cite{kara2023qlearningmdpsgeneralspaces}] \label{assumption:tx-lip}
    Let $P(\cdot|x)$ be a probability measure on $(\mathcal{X},\mathcal{F})$. There exit $L_P<\infty$ such that 
    \begin{align}
        \text{TV}(P(\cdot|x)-P(\cdot|x'))\le L_P |x-x'|, \quad \forall x,x'\in \mathcal{X}.
    \end{align}
\end{assumption}
\begin{proposition}[Quantization Error of f-Divergence, Proposition 13 in \cite{pillutla2023mauvescoresgenerativemodels}]\label{prop:quantization-error}
Let $P$ and $Q$ be two probability distributions over a common sample space $\mathcal{X}$.

Let ${S} = \{S_1, S_2, \dots, S_m\}$ be a partition of the space $\mathcal{X}$ into $m$ disjoint sets. The corresponding quantized distributions, $P_{\mathcal{S}}$ and $Q_{\mathcal{S}}$, are defined as multinomial distributions over the indices $\{1, \dots, m\}$.

Then, for any integer $k \ge 1$, and $f$-divergence functional $\mathcal{D}_f$, there exists a partition ${S}$ of size $m \le 2k$ such that the absolute difference between the original and the quantized f-divergence is bounded as follows:
\[
|\mathcal{D}_f(P,Q) - \mathcal{D}_{f}(P_{\mathcal{S}},Q_{\mathcal{S}})| \le \frac{f(0) + f^{*}(0)}{k}
\]
\end{proposition}



Theorem~\ref{conditional-row-tv-bound} is adapted from Theorem 3.1 and Lemma 3.1 of \cite{wolfer2023empiricalinstancedependentestimationmarkov}, which provide high-probability bounds on the row-wise total variation error of the empirical transition matrix for a finite-state, irreducible, aperiodic Markov chain observed over a single trajectory. The bound holds uniformly over all states and depends explicitly on the number of states and the trajectory length, while accounting for the chain’s dependence structure.

\begin{lemma}[Row-wise TV bound, \citet{wolfer2023empiricalinstancedependentestimationmarkov}]\label{conditional-row-tv-bound}
Let $(X_1, \dots, X_N)$ be an irreducible, aperiodic, stationary Markov chain on a finite state space $\mathcal{A}$ with $|\mathcal{A}|=k$, transition matrix $M$ and stationary distribution $\pi$. Then there exists a universal constant $C>0$ such that, for any $0 < \delta < 1$, the following holds with probability at least $1-\delta$:
\[
\max_{s\in\mathcal{A}} \sum_{a \in \mathcal{A}}
\left| \hat{M}(a|s)
 -  M(a\mid s)
\right|
\;\le\;
C \, \sqrt{\,\frac{\tau_{\mathrm{mix}} k\, \log\left(\frac{k N}{\delta}\right)}{\,N}} ,
\]
where $\tau_{\mathrm{mix}}$ is a mixing-time–type constant depending only on $M$ (for reversible chains one has $\tau_{\mathrm{mix}} \asymp 1/\gamma_{\mathrm{ps}}$, with $\gamma_{\mathrm{ps}}$ denoting the pseudo–spectral gap).
\end{lemma}


We will use the missing-mass bound from \cite{skorski2020missingmassconcentrationmarkov} to handle unseen transitions.
\begin{lemma}[Missing Mass Bound, Theorem 1 in \cite{skorski2020missingmassconcentrationmarkov}] \label{thm:missing-mass}
Let $(X_1, \dots, X_N)$ be an irreducible Markov chain over a finite state space $\mathcal{A}$ with stationary distribution $\pi_P$ and true transition matrix $M_P$. Define the \emph{transition missing mass} as
\[
\mathrm{Mmass }=\sum_{s \in \mathcal{A}} \sum_{s \in \mathcal{A}} \pi_P(s) M_P(a|s) \cdot \mathbf{1}\{\hat{M}_P(a|s) = 0\}.
\]

Let $T$ be the maximum hitting time of any set of states with stationary probability at least $0.5$. Then there exists an absolute constant $c > 0$ and independent Bernoulli random variables
\[
Q_{s,a} \sim \mathrm{Bernoulli}\left(e^{-c \cdot N \cdot \pi_P(s)M_P(a|s) / T}\right)
\]
such that for any subset $\mathcal{E} \subseteq \{(s,a) :s, a \in \mathcal{A}\}$ and any $n \ge 1$,
\[
\Pr\left[ \bigwedge_{(s,a) \in \mathcal{E}} \{\hat{M}_P(a|s) = 0\} \right]
\le \prod_{(s,a) \in \mathcal{E}} \Pr[ Q_{s,a} = 1 ].
\]
In particular, for any $t > 0$ it holds that
\[
\mathbb{E} \exp\left( t \cdot \mathrm{Mmass}\right)
\le
\mathbb{E} \exp\left( t \cdot \sum_{s \in \mathcal{A}} \sum_{s \in \mathcal{A}}  \pi_P(s)M_P(a|s) \, Q_{s,a} \right).
\]
For bounding deviations of weighted sums over Markov chains, we rely on the inequality of \cite{chung2012chernoffhoeffdingboundsmarkovchains}.
\end{lemma}
\begin{lemma}[Theorem 3.1 of \cite{chung2012chernoffhoeffdingboundsmarkovchains}] \label{thm:estimation-state}
Let \(M\) be an ergodic Markov chain on state space $\mathcal{A}$ with stationary distribution \(\pi\).
For \(\varepsilon \le 1/8\), let \(T(\varepsilon)\) denote its total-variation mixing time.
Consider a length-\(N\) chain $(X_1, \dots, X_N)$ on \(M\) with $X_1 \sim \varphi$.
For each \(s\in \mathcal{A}\), let \(f_s:\mathcal{A}\to[0,1]\) be a weight function with
\(\mathbb{E}_{X\sim\pi}[f_s(X)] = \pi(s)\).
Define the total weight \(N(s)=\sum_{i=1}^N f_s(X_i)\).
Then there exists an absolute constant \(c\) such that:
\[
\Pr\!\big[N(s) \ge (1+\delta)\pi(s) N\big]
\ \le\
c\,\|\varphi\|_{\pi}\times
\begin{cases}
\exp\!\big(-\,\delta^{2}\pi(s) N/(72\,T(\epsilon))\big), & 0\le \delta \le 1,\\[4pt]
\exp\!\big(-\,\delta\,\pi(s)N/(72\,T(\epsilon))\big), & \delta>1,
\end{cases}
\]
and, for \(0\le \delta\le 1\),
\[
\Pr\!\big[N(s) \le (1-\delta)\pi(s) N\big]
\ \le\
c\,\|\varphi\|_{\pi}\,\exp\!\big(-\,\delta^{2}\pi(s) N/(72\,T(\epsilon))\big).
\]
Here \(\langle u,v\rangle_{\pi}=\sum_{x}u_x v_x/\pi(x)\) and
\(\|u\|_{\pi}=\sqrt{\langle u,u\rangle_{\pi}}\).
\end{lemma}

\subsubsection{Auxiliary Results}
\begin{lemma}
\label{lem:uniform-bound}
For all $\alpha > 0$ and $p \in (0,1]$, it holds that
\begin{equation}
  p \, \max\{1, \log(1/p)\} \, e^{-\alpha p}
  \;\;\le\;\; \frac{2 + \log(1+\alpha)}{e\,\alpha}.
  \label{eq:uniform-ineq}
\end{equation}
\end{lemma}

\begin{proof}
Let $y = \alpha p \in (0,\alpha]$ and $A = \log \alpha$. Then we can rewrite
\[
p \, \max\{1,\log(1/p)\} e^{-\alpha p}
= \frac{1}{\alpha} \, y e^{-y} \, \max\{1, A - \log y\}.
\]
Next observe the inequality
\[
\max\{1, A-\log y\} \;\le\; 1 + A_{+} + (-\log y)_{+},
\]
where $x_{+}=\max\{0,x\}$ and $A_{+}=\max\{0,A\}$.

Therefore,
\[
y e^{-y} \max\{1,A-\log y\}
\;\le\; (1+A_{+}) \cdot y e^{-y} + y e^{-y}(-\log y)_{+}.
\]

Now use the following standard bounds:
\[
\sup_{y>0} y e^{-y} = \frac{1}{e}, 
\qquad
\sup_{0<y\le 1} y(-\log y) = \frac{1}{e}.
\]
Hence
\[
\sup_{y>0} y e^{-y} \max\{1, A-\log y\}
\;\le\; \frac{1+A_{+}}{e} + \frac{1}{e}
= \frac{2 + \log\alpha_{+}}{e},
\]
where $\log\alpha_{+} = \max\{0,\log \alpha\} \le \log(1+\alpha)$.

Substituting back into the expression, we obtain
\[
p \, \max\{1, \log(1/p)\} e^{-\alpha p}
\le \frac{1}{\alpha}\cdot \frac{2+\log(1+\alpha)}{e}.
\]
This proves Eq.~\ref{eq:uniform-ineq}.
\end{proof}

\begin{lemma}[Stationarity of Quantized Kernels]\label{lem:quant-stationary}
Let $S_P$ be the population first-order Markov transition kernel on the continuous surprisal space $\mathbb{R}$ with stationary law $\rho_P$. 
Fix a shared $k$-bin quantizer $q_k:\mathbb{R}\to\mathcal A=\{1,\dots,k\}$ with boundaries $b_1<\dots <b_{k-1}$ partitions space into bins $B_i=[b_i,b_{i+1})$. 
Define the row-stationary weights and the edge measure
\[
\pi_P(i)\;:=\;\rho_P(B_i),\qquad 
Z_P(i,j)\;:=\;\int_{B_i}\rho_P(\mathrm{d}x)\,S_P(B_j|x),\quad i,j\in\mathcal A,
\]
and the induced $k$-state transition kernel
\[
M_P(j\mid i)\;:=\;\frac{Z_P(i,j)}{\pi_P(i)}\qquad(\text{for }\pi_P(i)>0).
\]
Then $\pi_P$ is a stationary distribution of $M_P$, i.e.\ $\sum_{i}\pi_P(i)M_P(j\mid i)=\pi_P(j)$ for all $j\in\mathcal A$.
\end{lemma}

\begin{proof}
By definition,
\[
\sum_{i\in\mathcal A}\pi_P(i)M_P(j\mid i)
=\sum_{i\in\mathcal A}Z_P(i,j)
=\int_{\mathbb{R}}\rho_P(\mathrm{d}x)\,S_P(B_j|x)
=\rho_P(B_j)
=\pi_P(j),
\]
where the penultimate equality uses the stationarity of $\rho_P$ for $S_P$.
\end{proof}
\subsubsection{Proof of Theorem~\ref{thm:sta-error}} \label{sec:proof-sta-error}
In this step, we aim to bound the expected absolute difference between the estimated GJS divergence and the GJS divergence for the induced Markov kernels after discretization. The statistical error of our estimator is:
\begin{align}
    E_1 =  |\mathcal{D}_f(\hat M_P, \hat M_Q)-\mathcal{D}_f(M_P,M_Q)|
\end{align}
 The analysis will reveal how this error depends on the number of bins $k$ and the sequence length $N$. To analyze the statistical error, we will extend the logic used in  \citet{pillutla2023mauvescoresgenerativemodels}. We will apply Lemma~\ref{lemma:Lip} (Lemma 20 in  \citet{pillutla2023mauvescoresgenerativemodels}), which establishes an approximate Lipschitz property for the core component of any $f$-divergence. 

 \begin{proof}[Proof of Theorem~\ref{thm:sta-error}]
     To bound the statistical error $E_1$, we first decompose it and then expand the GJS function into a sum of its core components, allowing for the application of Lemma~\ref{lemma:Lip}.
Using the triangle inequality, we can bound the total statistical error by the sum of the errors arising from the estimation of each matrix individually:
     \begin{align}
    E_1 \leq \underbrace{|\mathcal{D}_f(\hat M_P, \hat M_Q)-\mathcal{D}_f(M_P,\hat M_Q)| }_{=:\mathcal{T}_1}+ \underbrace{|\mathcal{D}_f( M_P, \hat M_Q)-\mathcal{D}_f(M_P,M_Q)| }_{=:\mathcal{T}_2}\label{eq:E1-decomposed}
\end{align}

The f-divergence between two Markov chains, $M_A$ and $M_B$, is defined as the expected divergence of their row-wise conditional probability distributions, weighted by the stationary distribution of the second chain. Let $\pi_B(s)$ be the stationary probability of state $s$ for chain $M_B$. The f-divergence is:
\begin{align}
    D_f(M_A,M_B) = \sum_{s\in \mathcal{A}} \pi_B(s) \sum_{a \in \mathcal{A}} \psi (M_A(a|s),M_B(a|s))
\end{align}

Applying this to the first term of our decomposed error Eq.~\eqref{eq:E1-decomposed}, with $f=f_{JS}^{w}$, we get
\begin{align}
    &\mathcal{T}_1=|\mathcal{D}_f(\hat M_P, \hat M_Q)-\mathcal{D}_f(M_P,\hat M_Q)| \\
    &= \bigg| \sum_{s\in \mathcal{A}} \hat \pi_Q(s) \sum_{a \in \mathcal{A}} \psi (\hat M_P(a|s),\hat M_Q(a|s)) - \sum_{s\in \mathcal{A}} \hat \pi_Q(s) \sum_{a \in \mathcal{A}} \psi ( M_P(a|s), \hat M_Q(a|s)) \bigg| \\
    &=  \bigg| \sum_{s\in \mathcal{A}} \hat \pi_Q(s) \sum_{a \in \mathcal{A}} \big[\psi (\hat M_P(a|s),\hat M_Q(a|s))-\psi ( M_P(a|s),\hat M_Q(a|s))\big]  \bigg| \\
    & \leq \bigg| \sum_{s\in \mathcal{A}} \sum_{a \in \mathcal{A}} \big[\psi (\hat M_P(a|s),\hat M_Q(a|s))-\psi ( M_P(a|s),\hat M_Q(a|s))\big]  \bigg| \\
    & \leq \sum_{s\in \mathcal{A}} \sum_{a \in \mathcal{A}}   \bigg|\psi (\hat M_P(a|s),\hat M_Q(a|s))-\psi ( M_P(a|s),\hat M_Q(a|s))  \bigg|
\end{align}
\paragraph{Case 1: Observed Transitions} For a transition that appears in the human-written text sample, its empirical probability is $\hat{M}_P(a|s) = \frac{N_P(s,a)}{N_P(s)} \geq \frac{1}{N_P(s)}$, where $N_P(s)$ is the number of times state $s$ was visited in the sequence of length $N$. We apply the first inequality of Lemma~\ref{lemma:Lip} with $p' = \hat M_P(a|s), p=M_P(a|s)$, and $q=\hat M_Q(a|s)$. The term $\max\big(1,\log \frac{1}{\max (p,p')}\big)$ is bounded by $\log N_P(s)$ as long as $N_P(s)\geq 3$. Thus, the error for a single observed transition is bounded by:
\begin{align}
    \big|\psi (\hat M_P(a|s),\hat M_Q(a|s))-\psi ( M_P(a|s),\hat M_Q(a|s))  \big|&\leq (C_1\log N_P(s)+C')\big| \hat M_P(a|s) -M_P(a|s)  \big| \\
    &\leq (C_1\log N+C')\big| \hat M_P(a|s) -M_P(a|s)  \big|
\end{align}
where $C'$ is a constant absorbing $C_0^*$ and $C_2$. Summing over all observed transitions gives a bound proportional to the Total Variation (TV) distance between the estimated and true transition matrices, multiplied by a logarithmic factor.
\paragraph{Case 2: Missing Transitions} This case addresses transitions that have a non-zero true probability ($M_P(a|s)$) but were not observed in the finite sample, resulting in an empirical probability of $\hat M_P(a|s)=0$. This scenario is formally known as the missing mass problem for Markov chains, a non-trivial extension of the classic IID case due to the dependencies between samples. To analyze the error contribution, we directly bound the error for a single missing transition using Lemma~\ref{lemma:Lip}. Let $p'= \hat M_P(a|s)=0$ and $p=M_P(a|s)$. The error is now $\big|\psi (0 ,\hat M_Q(a|s))-\psi ( M_P(a|s),\hat M_Q(a|s))  \big|$. Applying the first inequality of Lemma~\ref{lemma:Lip}, we get:
\begin{align}
    \big|\psi (0 ,\hat M_Q(a|s))-\psi ( M_P(a|s),\hat M_Q(a|s))  \big| &\leq  (C_1\max \bigg(1,\log \frac{1}{ M_P(a|s)}\bigg)+C')\big| 0 -M_P(a|s)  \big| 
    \\ &= (C_1\max \bigg(1,\log \frac{1}{ M_P(a|s)}\bigg)+C')M_P(a|s) 
\end{align}

This bound shows that the error from a missing transition is proportional to its true probability $M_P(a|s)$, scaled by its information content. The total error from this case is the sum of these individual bounds over all unobserved transitions. This sum constitutes the missing transition mass of the Markov chain.


We summarize the following: 
\begin{align}
    \mathbb E[\mathcal{T}_1] \leq\big(C_1\log N+C'\big)\cdot\sum_{s\in \mathcal{A}}\alpha_{N_P(s)}(M_P(\cdot|s))+\big(C_1+C'\big)\sum_{s\in \mathcal{A}}\beta_{N_P(s)}(M_P(\cdot|s)) \label{eq:Expectation-T1}
\end{align}
where $M_P(\cdot|s)$ is a $k$-dimensional probability distribution corresponding to state $s$, and we formally define the row-wise error terms:


 \begin{itemize}
    \item Row-wise TV term $\alpha_{N_P(s)}(M_P(\cdot|s))$: This term sums the error from observed transitions in state $s$.
    \begin{align}
        \mathbb E[\alpha_{N_P(s)}(M_P(\cdot|s))]  = \mathbb E \bigg[ \sum_{\substack{a \in \mathcal{A}, \\ \text{s.t.} \hat M_P(a|s)>0} } \big|\hat M_P(a|s)-M_P(a|s)\big|\bigg] \label{eq:alpha-state}
    \end{align}
    \item Row-wise Missing Mass term $\beta_{N_P(s)}(M_P(\cdot|s))$ This term sums the error from unobserved transitions in state $s$.
    \begin{align}
        \mathbb E[\beta_{N_P(s)}(M_P(\cdot|s))]= \mathbb E \bigg[ \sum_{\substack{a \in \mathcal{A}, \\ \text{s.t.} \hat M_P(a|s)=0} } M_P(a|s) \cdot \max \Big(1,\log \frac{1}{M_P(a|s)}\Big)\bigg]
    \end{align}
\end{itemize}
Then we use Lemma~\ref{conditional-row-tv-bound} to upper bound Eq~\ref{eq:alpha-state}.
\begin{align}
    \mathbb E[\alpha_{N_P(s)}(M_P(\cdot|s))] &= \mathbb E \bigg[ \sum_{\substack{a \in \mathcal{A}, \\ \text{s.t.} \hat M_P(a|s)>0} } \big|\hat M_P(a|s)-M_P(a|s)\big|\bigg]\\ &\leq \mathbb E \bigg[ \sum_{a \in \mathcal{A} } \big|\hat M_P(a|s)-M_P(a|s)\big|\bigg] \\
    & = O(\sqrt{\,\frac{k \, \log\left(k N\right)}{\,N}} ) \label{apply-B3}
\end{align}
where Eq.~\ref{apply-B3} follows Lemma~\ref{conditional-row-tv-bound} by inverting its tail bound and integrating to expectation;
the mixing-time constant is absorbed into $O(1)$ under Assumption~\ref{assumption:markov-chain}.

Lemma~\ref{thm:missing-mass} gives an exponential tail for the event $\hat{M}_P(a|s) = 0$: for some absolute constant $c>0$ and $T$ the maximum hitting time of any set with stationary probability at least 0.5,
\begin{align}
    \mathbb P[\hat M_P(a|s)=0]\leq \exp\Big(-\frac{cN}{T}\pi_P(s)M_P(a|s)\Big)
\end{align}





Then we upper bound the missing mass term $\mathbb E[\beta_{N_P(s)}(M_P(\cdot|s))]$. Let $p_a = M_P(a|s)$ and $\Gamma = \frac{cN}{T}\pi_P(s)$.
\begin{align}
    \mathbb E[\beta_{N_P(s)}(M_P(\cdot|s))] &= \sum_{a\in \mathcal{A}} p_a \max \Big(1,\frac{1}{p_a}\Big) \mathbb P[\hat M_P(a|s)=0] \\
    &= \sum_{a\in \mathcal{A}} p_a \max \Big(1,\frac{1}{p_a}\Big) e^{-\Gamma p_a} \\
    & \leq  \sum_{a\in \mathcal{A}} \frac{2+\log(1+\Gamma)}{e \Gamma} \label{apply-inequality}\\
    & = \frac{kT}{ecN \pi_P(s)} \Big(2+\log \Big(1+\frac{cN\pi_P(s)}{T}\Big)\Big)
\end{align}
where Eq.~\ref{apply-inequality} follows Lemma~\ref{lem:uniform-bound} for all $\Gamma>0$ and $p_a\in (0,1]$. Assuming $\pi_P(s)\geq \frac{c_0}{k}$ for some constant $c_0>0$ and $T=O(1)$, we obtain
\begin{align}
    \mathbb E[\beta_{N_P(s)}(M_P(\cdot|s))]  = O\Big(\frac{k^2}{N}\log (1+\frac{N}{k}) \Big) \label{eq:beta-final-bound}
\end{align}
By Eq.~\ref{eq:Expectation-T1}, Eq.~\ref{apply-B3}, and Eq.~\ref{eq:beta-final-bound} we obtain
\begin{align}
    \mathcal{T}_1 = O \Big(\log N\cdot \sqrt{\frac{k^3\log(kN)}{N}}+\frac{k^3}{N}\log\big(1+\frac{N}{k}\big)\Big)
\end{align}
Next we bound $ \mathcal{T}_2$. 
\begin{align}
    \mathcal{T}_2 &= |\mathcal{D}_f( M_P, \hat M_Q)-\mathcal{D}_f(M_P,M_Q)|  \\
    & = \bigg| \sum_{s\in \mathcal{A}} \hat \pi_Q(s) \sum_{a \in \mathcal{A}} \psi (M_P(a|s),\hat M_Q(a|s)) - \sum_{s\in \mathcal{A}} \pi_Q(s) \sum_{a \in \mathcal{A}} \psi ( M_P(a|s),  M_Q(a|s)) \bigg| \\
    & =  \bigg| \sum_{s\in \mathcal{A}} \hat \pi_Q(s) \sum_{a \in \mathcal{A}} \psi (M_P(a|s),\hat M_Q(a|s)) - \sum_{s\in \mathcal{A}} \hat \pi_Q(s) \sum_{a \in \mathcal{A}} \psi (M_P(a|s),M_Q(a|s))\nonumber\\&+ \sum_{s\in \mathcal{A}} \hat \pi_Q(s) \sum_{a \in \mathcal{A}} \psi (M_P(a|s),M_Q(a|s))- \sum_{s\in \mathcal{A}} \pi_Q(s) \sum_{a \in \mathcal{A}} \psi ( M_P(a|s),  M_Q(a|s)) \bigg| \\
    &\leq \bigg| \sum_{s\in \mathcal{A}} \hat \pi_Q(s) \sum_{a \in \mathcal{A}} \psi (M_P(a|s),\hat M_Q(a|s)) - \sum_{s\in \mathcal{A}} \hat \pi_Q(s) \sum_{a \in \mathcal{A}} \psi (M_P(a|s),M_Q(a|s))\bigg|\nonumber\\&+  \bigg|\sum_{s\in \mathcal{A}} \hat \pi_Q(s) \sum_{a \in \mathcal{A}} \psi (M_P(a|s),M_Q(a|s))- \sum_{s\in \mathcal{A}} \pi_Q(s) \sum_{a \in \mathcal{A}} \psi ( M_P(a|s),  M_Q(a|s)) \bigg| \\
        & \leq  \bigg|\sum_{s\in \mathcal{A}}\sum_{a\in \mathcal{A}}[\psi ( M_P(a|s),\hat M_Q(a|s))-\psi ( M_P(a|s),M_Q(a|s)) ] \bigg| \nonumber\\& +  \bigg| \sum_{s\in \mathcal{A}}\big(\hat \pi_Q(s)-\pi_Q(s)\big)\sum_{a\in \mathcal{A}}\psi ( M_P(a|s),M_Q(a|s))\bigg| \\
         & \leq \underbrace{\sum_{s\in \mathcal{A}}\sum_{a\in \mathcal{A}} \bigg|\psi ( M_P(a|s),\hat M_Q(a|s))-\psi ( M_P(a|s),M_Q(a|s))  \bigg| }_{=:\mathcal{T}_{2,1}}\nonumber\\& + \underbrace{ \bigg| \sum_{s\in \mathcal{A}}\big(\hat \pi_Q(s)-\pi_Q(s)\big)\sum_{a\in \mathcal{A}}\psi ( M_P(a|s),M_Q(a|s))\bigg|}_{=:\mathcal{T}_{2,2}}
\end{align}
By symmetry, bounding $\mathcal{T}_{2,1}$ proceeds identically to $\mathcal{T}_{1}$, and yields the same rate as $\mathcal{T}_{1}$. To upper bound $\mathcal{T}_{2,2}$, we consider 
\begin{align}
    \sum_{a\in \mathcal{A}}\psi ( M_P(a|s),M_Q(a|s)) =  \sum_{a\in \mathcal{A}} M_Q(a|s)f_{JS}^w(M_P(a|s)/M_Q(a|s)) \le H(w) \le \log 2
\end{align}
where $H(w)=-[w\log(w)+(1-w)\log (1-w)]$ with $w=\frac{\alpha}{1+\alpha}\in[0,1]$ is the binary entropy function of which the absolute maximum possible value is $\log 2$. To upper bound $\mathcal{T}_{2,2}$,
\begin{align}
    \mathcal{T}_{2,2} \leq \log 2 \cdot \mathbb E \big| \hat \pi_Q-\pi_Q\big|
\end{align}
We apply Lemma~\ref{thm:estimation-state} to upper bound $\mathcal{T}_{2,2}$. Consider $\hat \pi_Q(s)=\frac{N_Q(s)}{N}$, for any $\delta>0$, we have 
\[
\Pr\!\big[N_Q(s) \ge (1+\delta)\pi_Q(s) N\big]
\ \le\
c\,\|\varphi\|_{\pi_Q}\times
\begin{cases}
\exp\!\big(-\,\delta^{2}\pi_Q(s) N/(72\,T)\big), & 0\le \delta \le 1,\\[4pt]
\exp\!\big(-\,\delta\,\pi_Q(s)N/(72\,T)\big), & \delta>1,
\end{cases}
\]
and similarly for the lower tail with $0<\delta<1$. With $\epsilon=\delta \pi_Q(S)$, we have
\begin{align}
    \Pr\!\big[|\hat \pi_Q(s)-\pi_Q(s)| \ge \epsilon\big]
\ \le\
2c\,\|\varphi\|_{\pi_Q}\times
\begin{cases}
\exp\!\big(-\,\epsilon^{2} N/(72\,T\pi_Q(s))\big), & 0\le \epsilon \le \pi_Q(s),\\[4pt]
\exp\!\big(-\,\epsilon\,N/(72\,T)\big), & \epsilon>\pi_Q(s),
\end{cases}
\end{align}
Using $\mathbb E|Z|=\int_{0}^{\infty}\Pr (|Z|\ge\epsilon)$ and splitting the integral at $\pi_Q(s)$,
\begin{align}
    \mathbb E\big[|\hat \pi_Q(s)-\pi_Q(s)| \big]&\le 2c\|\varphi\|_{\pi_Q}\Big( \int_{0}^{\pi_Q(s)}e^{-\frac{N\epsilon^2}{72T\pi_Q(s)}}d\epsilon+\int_{\pi_Q(s)}^{\infty}e^{-\frac{N\epsilon}{72T}}d\epsilon\Big) \\
    &\leq 2c\|\varphi\|_{\pi_Q}\bigg(C\sqrt{\frac{T\pi_Q(s)}{N}} +\frac{72T}{N}\exp\big(-\frac{N\pi_Q(s)}{72T}\big)\bigg)\\
    & = O\bigg(\|\varphi\|_{\pi_Q}\sqrt{\frac{T\pi_Q(s)}{N}}\bigg)
\end{align}
Thus we obtain
\begin{align}
     \mathbb E\big[|\hat \pi_Q-\pi_Q| \big]&=\sum_{s\in \mathcal{A}} \mathbb E\big[|\hat \pi_Q(s)-\pi_Q(s)| \big] \\
     & \le \sum_{s\in \mathcal{A}}  C \|\varphi\|_{\pi_Q}\sqrt{\frac{T\pi_Q(s)}{N}} \\
     & = C \|\varphi\|_{\pi_Q}\sqrt{\frac{T}{N}}   \sum_{s\in \mathcal{A}} \sqrt{\pi_Q(s)}\\
     &
     \le   C \|\varphi\|_{\pi_Q}\sqrt{\frac{Tk}{N}} \\
     &= O\Big(\frac{k}{\sqrt{N}}\Big) \label{eq:exp-pi}
\end{align}
where Eq.~\ref{eq:exp-pi} holds since $\|\varphi\|_{\pi_Q}=\frac{1}{\sqrt{\pi_Q(s_0)}}\le \frac{1}{\sqrt{\min_{s\in\mathcal{A}}\pi_Q(s)}}=O(\sqrt{k})$ for the first state $s_0$, and $T=O(1)$. To sum up, $\mathcal{T}_{2,2}=O\Big(\frac{k}{\sqrt{N}}\Big)$, and the rate for the total statistical error is
\begin{align}
    E_1 &\le \mathcal{T}_{1}+\mathcal{T}_{2,1}+\mathcal{T}_{2,2}\\
    &= O\Big(\log N\cdot \sqrt{\frac{k^3\log(kN)}{N}}+\frac{k^3}{N}\log\big(1+\frac{N}{k}\big)\Big)\nonumber \\&+O\Big(\log N\cdot \sqrt{\frac{k^3\log(kN)}{N}}+\frac{k^3}{N}\log\big(1+\frac{N}{k}\big)\Big)+O\Big(\frac{k}{\sqrt{N}}\Big) \\
    & = O\Big(\log N\cdot \sqrt{\frac{k^3\log(kN)}{N}}+\frac{k^3}{N}\log\big(1+\frac{N}{k}\big)+\frac{k}{\sqrt{N}} \Big) \label{eq:final-sta-error}
\end{align}



 \end{proof}

 \subsubsection{Proof of Proposition~\ref{prop:quant-error}}\label{sec:proof-quant-error}
\begin{proof}[Proof of Proposition~\ref{prop:quant-error}]
    Let $\rho_P$ and $\rho_Q$ be the continuous stationary distributions of $\mathcal{S}_P$ and $\mathcal{S}_Q$ respectively. We expand $\mathcal{D}_f(\mathcal{S}_P,\mathcal{S}_Q)$ and $\mathcal{D}_f(M_P,M_Q)$,
    \begin{align}
    \mathcal{D}_f(\mathcal{S}_P,\mathcal{S}_Q) &= \int_{\mathbb{R}}\rho_Q(\text{d}x) \mathcal{D}_f (\mathcal{S}_P(\cdot|x),(\mathcal{S}_Q(\cdot|x))\\
        \mathcal{D}_f(M_P,M_Q) &= \sum_{i\in \mathcal{A}} \pi_Q(i) \mathcal{D}_f(M_P(\cdot|i),M_Q(\cdot|i))
    \end{align}
     The quantizer $q_k: \mathbb{R}\rightarrow \mathcal{A}=[k]$ with boundaries $b_1<\dots <b_{k-1}$ partitions space into bins $B_i=[b_i,b_{i+1})$. Let $\rho_Q(B_i)=\int_{B_i}\text{d}\rho_Q(x)$, then $\pi_Q(i)=\rho_Q(B_i)$.

     
    Define two intermediate objects $U_P$ and $U_Q$ to be markov kernel such that each has a discrete state index $i\in \mathcal{A}$, within a given state $i$, the observable variable $x$ lives in a continuous space $\mathbb{R}$, The corresponding stationary distributions over states are $\pi_P$ for $P$ and $\pi_Q$ for $Q$. Thus 
    \begin{align}
    \mathcal{D}_f(U_P,U_Q) = \sum_{i \in \mathcal{A}} \pi_Q(i) \mathcal{D}_f(\mathcal{S}_P(\cdot|i),\mathcal{S}_Q(\cdot|i)) 
    \end{align}
    where $\mathcal{S}_P(\cdot|i)=\mathbb{E}_{x\sim \rho_Q(B_i)}[\mathcal{S}_P(\cdot|x)]$ and similarly for $\mathcal{S}_Q(\cdot|i)$. We have
    \begin{align}
    |\mathcal{D}_f(\mathcal{S}_P,\mathcal{S}_Q)-\mathcal{D}_f(M_P,M_Q)| \le |\mathcal{D}_f(\mathcal{S}_P,\mathcal{S}_Q)-\mathcal{D}_f(U_P,U_Q)|+|\mathcal{D}_f(U_P,U_Q)-\mathcal{D}_f(M_P,M_Q)|
    \end{align}
    The second term is bounded as
    \begin{align}
        &|\mathcal{D}_f(U_P,U_Q)-\mathcal{D}_f(M_P,M_Q)| \\&= \Big|\sum_{i \in \mathcal{A}} \pi_Q(i) \mathcal{D}_f(\mathcal{S}_P(\cdot|i),\mathcal{S}_Q(\cdot|i)) -\sum_{i\in \mathcal{A}} \pi_Q(i) \mathcal{D}_f(M_P(\cdot|i),M_Q(\cdot|i))\Big| \\
        &\le \sum_{i\in \mathcal{A}} \pi_Q(i)\Big|\mathcal{D}_f(\mathcal{S}_P(\cdot|i),\mathcal{S}_Q(\cdot|i)) -\mathcal{D}_f(M_P(\cdot|i),M_Q(\cdot|i))\Big| \label{eq:U-M-result-2}\\
        & = O(\frac{1}{k}) \label{eq:U-M-result}
    \end{align}
    Eq.~\ref{eq:U-M-result} holds by applying Proposition~\ref{prop:quantization-error} to each term in Eq.~\ref{eq:U-M-result-2}, yielding an $O(1/k)$ bound per term. Since the weighted sum of $O(1/k)$ terms remains $O(1/k)$, the overall bound follows. The first term is 
    \begin{align}
     &\mathcal{D}_f(\mathcal{S}_P,\mathcal{S}_Q)-\mathcal{D}_f(U_P,U_Q) \\&=  \int_{\mathbb{R}}\rho_Q(\text{d}x) \mathcal{D}_f (\mathcal{S}_P(\cdot|x),(\mathcal{S}_Q(\cdot|x)) - \sum_{i \in \mathcal{A}} \pi_Q(i) \mathcal{D}_f(\mathcal{S}_P(\cdot|i),\mathcal{S}_Q(\cdot|i)) \\
     & = \sum_{i=1}^k \int_{B_i} \rho_Q(\text{d}x) \mathcal{D}_f (\mathcal{S}_P(\cdot|x),(\mathcal{S}_Q(\cdot|x))- \sum_{i \in \mathcal{A}} \pi_Q(i) \mathcal{D}_f(\mathcal{S}_P(\cdot|i),\mathcal{S}_Q(\cdot|i)) \\
     &=\sum_{i\in \mathcal{A}}  \rho_Q(B_i) \mathbb{E}_{x\sim \rho_Q(B_i)}[\mathcal{D}_f (\mathcal{S}_P(\cdot|x),(\mathcal{S}_Q(\cdot|x))]- \sum_{i \in \mathcal{A}} \pi_Q(i) \mathcal{D}_f(\mathcal{S}_P(\cdot|i),\mathcal{S}_Q(\cdot|i)) \\
     & =\sum_{i\in \mathcal{A}}  \pi_Q(i) \mathbb{E}_{x\sim \rho_Q(B_i)}[ \mathcal{D}_f (\mathcal{S}_P(\cdot|x),(\mathcal{S}_Q(\cdot|x))]- \sum_{i \in \mathcal{A}} \pi_Q(i) \mathcal{D}_f(\mathcal{S}_P(\cdot|i),\mathcal{S}_Q(\cdot|i))\\
     & = \sum_{i\in \mathcal{A}}  \pi_Q(i)\big[\mathbb{E}_{x\sim \rho_Q(B_i)}[ \mathcal{D}_f (\mathcal{S}_P(\cdot|x),(\mathcal{S}_Q(\cdot|x))] - \mathcal{D}_f(\mathcal{S}_P(\cdot|i),\mathcal{S}_Q(\cdot|i))\big] \\
     & =: \sum_{i\in \mathcal{A}}  \pi_Q(i) J_i
    \end{align}
    Because $\mathcal{D}_f$ is jointly convex,
    \begin{align}
\mathcal{D}_f(\mathbb{E}_{x\sim \rho_Q(B_i)}[\mathcal{S}_P(\cdot|x)],\mathbb{E}_{x\sim \rho_Q(B_i)}[\mathcal{S}_Q(\cdot|x)]\big) \le \mathbb{E}_{x\sim \rho_Q(B_i)}[ \mathcal{D}_f (\mathcal{S}_P(\cdot|x),(\mathcal{S}_Q(\cdot|x))]
    \end{align}
    Therefore,
    \begin{align}
    |\mathcal{D}_f(\mathcal{S}_P,\mathcal{S}_Q)-\mathcal{D}_f(U_P,U_Q)| = \sum_{i\in \mathcal{A}}  \pi_Q(i) J_i
    \end{align}
Lemma~\ref{lemma:Lip} implies a Lipschitz-type continuity bound in total variation distance, that is
\begin{align}
    |\mathcal{D}_f(P,Q)-\mathcal{D}_f(P',Q') | \le 2L_f (\text{TV}(P,P')+\text{TV}(Q,Q')) \label{lip-tv}
\end{align}
where $L_f$ depends on $C_1$, $C_1^*$, $C_2$, $C_2^*$ in Lemma~\ref{lemma:Lip}. Applying Eq.~\ref{lip-tv} to $J_i$ yields
\begin{align}
    J_i &= \mathbb{E}_{x\sim \rho_Q(B_i)}[ \mathcal{D}_f (\mathcal{S}_P(\cdot|x),(\mathcal{S}_Q(\cdot|x))] - \mathcal{D}_f(\mathcal{S}_P(\cdot|i),\mathcal{S}_Q(\cdot|i)) \\
    &\le \mathbb{E}_{x\sim \rho_Q(B_i)}[\mathcal{D}_f (\mathcal{S}_P(\cdot|x),(\mathcal{S}_Q(\cdot|x)) - \mathcal{D}_f(\mathcal{S}_P(\cdot|i),\mathcal{S}_Q(\cdot|i))] \\
    & \le 2L_f \mathbb{E}_{x\sim \rho_Q(B_i)} [\text{TV}(\mathcal{S}_P(\cdot|x),\mathcal{S}_P(\cdot|i)) + \text{TV}(\mathcal{S}_Q(\cdot|x),\mathcal{S}_Q(\cdot|i))]
\end{align}
By Assumption~\ref{assumption:tx-lip}, 
\begin{align}
    \text{TV}(\mathcal{S}_P(\cdot|x),\mathcal{S}_P(\cdot|i)) + \text{TV}(\mathcal{S}_Q(\cdot|x),\mathcal{S}_Q(\cdot|i))\le (L_P+L_Q)\mathbb{E}_{x'\sim\rho_Q( B_i)}|x-x'|
\end{align}
Let $c_i$ be the centroid of $B_i$ and define the mean radius $r_i = \mathbb{E}_{x\sim \rho_Q(B_i)}|x-c_i|$. For any $x\in B_i$,
\begin{align}
    \mathbb{E}_{x'\sim \rho_Q(B_i)}|x-x'| \le |x-c_i|+ \mathbb{E}_{x'\sim \rho_Q(B_i)}|x'-c_i|=|x-c_i|+r_i
\end{align}
Then,
\begin{align}
    (L_P+L_Q)\mathbb{E}_{x'\sim \rho_Q(B_i)}|x-x'| \le (L_P+L_Q)\mathbb{E}_{x'\sim \rho_Q(B_i)}|x-c_i|+r_i = 2(L_P+L_Q)r_i
\end{align}
Then,
\begin{align}
    J_i \le 4L_f(L_P+L_Q)r_i
\end{align}
Summing over buckets with weight $\pi_P(i)$ gives:
\begin{align}
    |\mathcal{D}_f(\mathcal{S}_P,\mathcal{S}_Q)-\mathcal{D}_f(U_P,U_Q)| &= \sum_{i\in \mathcal{A}}  \pi_Q(i) J_i \\
    &\le 4L_f(L_P+L_Q)\sum_{i\in \mathcal{A}}\pi_Q(i)  r_i\\
    & = 4L_f(L_P+L_Q) \mathbb{E}_{x \sim \rho_Q} [x-q_k(x)]\\
    &=O(1/k) \label{S-U}
\end{align}
By Eq.~\ref{eq:U-M-result} and Eq.~\ref{S-U}, 
\begin{align}
|\mathcal{D}_f(\mathcal{S}_P,\mathcal{S}_Q)-\mathcal{D}_f(M_P,M_Q)|\le \frac{c}{k} \label{eq:quant-error}
\end{align}

\subsubsection{Balancing Two Errors}
A clear choice for k is found by balancing the dominant statistical error (Eq.~\ref{eq:final-sta-error}) with the quantization error (Eq.~\ref{eq:quant-error}) in rate form, ignoring logarithmic factors. The leading statistical term scales as $c_1 k^{\frac{3}{2}}N^{-\frac{1}{2}}$ and the quantization term as $\frac{c_2}{k}$. Minimizing their sum $f(k)=c_1 k^{\frac{3}{2}}N^{-\frac{1}{2}}+\frac{c_2}{k}$  by first-order condition $f'(k)=0$ yields that 
\begin{align}
    k^* = \Big( \frac{4c_2}{3c_1}\Big)^{\frac{2}{7}}N^{\frac{1}{5}}
\end{align}
Thus, up to constants and polylog factors, the optimal bin count is $k^* = \Theta(N^{\frac{1}{5}})$.
\end{proof}




\subsection{Decision Statistic Analysis}
\subsubsection{Auxiliary Results from Literature}
\begin{lemma}[Second-Order Taylor Expansion of Generalized Jensen Shannon Divergence, \cite{zhou2018secondorderasymptoticallyoptimalstatistical}] \label{GJS-exp}
Let $P_1, P_2 \in \mathcal{P}(\mathcal{X})$ be two distinct probability distributions over a finite alphabet $\mathcal{X}$, representing a point of expansion. Let $\hat{P}_1, \hat{P}_2 \in \mathcal{P}(\mathcal{X})$ be two other probability distributions in a neighborhood of $(P_1, P_2)$. Let $\alpha$ be a fixed positive constant. The Generalized Jensen-Shannon (GJS) divergence, viewed as a function $GJS(\hat{P}_1, \hat{P}_2, \alpha)$, has the following second-order Taylor approximation around the point $(P_1, P_2)$.
\begin{align}
    % GJS(\hat{P}_1, \hat{P}_2, \alpha) = & GJS(P_1, P_2, \alpha) 
    % + \sum_{x \in \mathcal{X}} (\hat{P}_1(x) - P_1(x))\alpha\iota_1(x) 
    %  + \sum_{x \in \mathcal{X}} (\hat{P}_2(x) - P_2(x))\iota_2(x) + R
        \text{GJS}(\hat{P}_1, \hat{P}_2, \alpha) = & \underbrace{\text{GJS}(P_1, P_2, \alpha)}_{\text{Zeroth-Order Term}} 
    \underbrace{+ \sum_{x \in \mathcal{X}} (\hat{P}_1(x) - P_1(x))\alpha\iota_1(x) + \sum_{x \in \mathcal{X}} (\hat{P}_2(x) - P_2(x))\iota_2(x)}_{\text{First-Order Term}} \nonumber \\
    & \underbrace{+  O\left(||\hat{P}_1 - P_1||^2 + ||\hat{P}_2 - P_2||^2\right)}_{\text{Remainder Term}}
\end{align}
where the remainder term is of the order of the squared Euclidean distance between the points, $\text{GJS}(P_1, P_2, \alpha)$ is the zeroth-order term, the GJS function evaluated at the point of expansion $(P_1, P_2)$. The first-order term is a linear function of the differences $(\hat{P}_1 - P_1)$ and $(\hat{P}_2 - P_2)$. The summation is taken over all symbols $x$ in the alphabet $\mathcal{X}$. The partial derivatives of the GJS function, evaluated at $(P_1, P_2)$, are given by the information densities.
\begin{align}
    \iota_1(x) :=\iota_1(x|P_1, P_2, \alpha)= \log\frac{(1+\alpha)P_1(x)}{\alpha P_1(x) + P_2(x)}\\
        \iota_2(x) := \iota_2(x|P_1, P_2, \alpha) = \log\frac{(1+\alpha)P_2(x)}{\alpha P_1(x) + P_2(x)}
\end{align}
   
\end{lemma}


\begin{lemma}[Central Limit Theorem for Additive Functionals, \cite{HOLZMANN20051518}] \label{lemma:clt-AF}
    Let $(X_1, \dots, X_N)$ be a stationary, ergodic, discrete-time Markov chain with state space $\mathcal{S}$, transition operator $M$, and unique stationary distribution $\pi$. Let $f: \mathcal{S} \rightarrow \mathbb{R}$ be a real-valued function defined on the state space, and assume its expectation with respect to the stationary distribution is zero, i.e., $\mathbb{E}_{\pi} [f(x)]=0.$ Consider the additive functional $S_N(f)=\sum_{i=1}^Nf(X_i)$. If a martingale approximation to $S_N(f)$ exits, then the Central Limit Theorem holds, i.e.:
    \begin{align}
        \frac{S_N(f)}{\sqrt{N}} \xrightarrow{d} N(0,\sigma^2(f)) 
    \end{align}
    The term $\sigma^2(f)$ is the asymptotic variance of the process.
\end{lemma}


\begin{lemma}[Asymptotic Variance for Markov Chains, \cite{HOLZMANN20051518}] \label{lemma:var-markov-chain}
    Under the same conditions as Lemma~\ref{lemma:clt-AF}, the asymptotic variance $\sigma^2(f)$ of the additive functional $S_N(f)$ is given by:
\begin{align}
    \sigma^2(f) = 2\lim_{\epsilon\rightarrow 0} \langle g_{\epsilon},f \rangle - \|f\|^2
\end{align}
where $g_{\epsilon}$ is the solution to the following equation $((1+\epsilon)I-M)^{-1}$, which is a function defined on the state space $\mathcal{A}$. $\langle g_{\epsilon},f \rangle$ is the inner product in the Hilbert space $L_2(\pi)$, calculated as $\langle g_{\epsilon},f \rangle =\sum_{x\in \mathcal{A}}\pi(x)g_{\epsilon}(x)f(x)$. $\|f\|^2$ is the squared norm of the function $f$ in the space $L_2(\pi)$ , which is its variance with respect to the stationary distribution.

\end{lemma}
\subsubsection{Proof of Proposition~\ref{prop:GJS-n-as-LLR}} \label{sec:proof-llr}
\begin{proof}[Proof of Proposition~\ref{prop:GJS-n-as-LLR}]
Let $\mathcal{F}_k$ be the family of stationary first-order Markov models on $\mathcal{A}:=[k]$. 
Consider the following likelihood ratio,
     \begin{align}
      \Lambda_{n,N} &=   \frac{1}{n}\log\!
\frac{\displaystyle \sup_{M,M'\in\mathcal{F}_k}
M\!\bigl((a^P_{1:N},a^T_{1:n})\bigr)\,M'\!\bigl(a^Q_{1:N}\bigr)}
{\displaystyle \sup_{M,M'\in\mathcal{F}_k}
M\!\bigl(a^P_{1:N}\bigr)\,M'\!\bigl((a^Q_{1:N},a^T_{1:n})\bigr)} \\
&=\frac{1}{n}\log \frac{\hat{M}_{\alpha 1}((a^P_{1:N},a^T_{1:n}))\hat{M}_Q(a^Q_{1:N})}{\hat{M}_P(a^P_{1:N})\hat{M}_{\alpha 2}((a^Q_{1:N},a^T_{1:n}))}
     \end{align}
where $(a^P_{1:N},a^T_{1:n})$ denotes the concatenation of $a^P_{1:N}$ and $a^T_{1:n}$, $\hat{M}_{\alpha 1}=\frac{\alpha \hat{M}_P+ \hat{M}_T}{1+\alpha}$, and $\hat{M}_{\alpha 2}=\frac{\alpha \hat{M}_Q+ \hat{M}_T}{1+\alpha}$. By Eq. (4)-(6) in \citet{32134}, we have
\begin{align}
\sup_{M \in\mathcal{F}_k}M\bigl((a^P_{1:N},a^T_{1:n})\bigr)
   = 2^{-(N+n)\,H((a^P_{1:N},a^T_{1:n}))}, 
\sup_{M'\in\mathcal{F}_k}M'\bigl(a^Q_{1:N}\bigr)
   = 2^{-N\,H(a^Q_{1:N})}, \\
   \sup_{M'\in\mathcal{F}_k}M'\bigl((a^Q_{1:N},a^T_{1:n})\bigr)
   = 2^{-(N+n)\,H((a^Q_{1:N},a^T_{1:n}))}, 
\sup_{M\in\mathcal{F}_k}M\bigl(a^P_{1:N}\bigr)
   = 2^{-N\,H(a^P_{1:N})},
\end{align}
where $H(\cdot)$ is the empirical conditional entropy per transition in the corresponding sequence. Plugging into the ratio gives

\begin{align}
    \,
\Lambda_{n,N}
  = \frac{N+n}{n}H((a^P_{1:N},a^T_{1:n}))-\frac{N}{n}\,H(a^P_{1:N})
    -\bigl[\frac{N+n}{n}H((a^Q_{1:N},a^T_{1:n}))-\frac{N}{n}\,H(a^Q_{1:N})\bigr]
\,
\end{align}
With weight \(\alpha=N/n\),
\begin{align}
    \Delta\mathrm{GJS}_n
  = \frac{N+n}{n}H((a^P_{1:N},a^T_{1:n}))-H(a^T_{1:n})-\frac{N}{n}H(a^P_{1:n})
    \nonumber\\\;-\;
    \Bigl[\frac{N+n}{n}H((a^Q_{1:N},a^T_{1:n}))-H(a^T_{1:n})-\frac{N}{n}H(a^Q_{1:N})\Bigr] \label{eq:gjs-entropy}
\end{align}
The two terms \(\pm H(a^T_{1:n})\) cancel. Thus we obtain $\Delta \text{GJS}_n = \Lambda_{n,N}$

    
\end{proof}
\subsubsection{Asymptotic Normality of $\Delta\mathrm{GJS}_n$} \label{sec:proof-asymp-gaussian}
\begin{theorem}[Asymptotic normality of $\Delta\mathrm{GJS}_n$ ]\label{thm:asymptotic-normality}
Assume the setting of Section~\ref{sec:setup} with $\alpha=N/n$ and standard ergodicity, $\Delta \rm{GJS}_n$ is asymptotically normal. Under $H_0: M_T=M_P$,
$\mu_{H_0}=-\mathrm{GJS}(M_Q,M_P,\alpha)<0$,
and
$
\sigma_{H_0}^2 \;=\; \frac{\alpha^2}{N^2}\,\sigma_{1,0}^2 \;+\; \frac{1}{n^2}\,\sigma_{2,0}^2$,
where $\sigma_{1,0}^2$ is the long-run variance of the $P$-reference-side information-density sum and
$\sigma_{2,0}^2$ is the long-run variance of the test-side information-density sum (details in Appendix~D).
Under $H_1: M_T=M_Q$, 
$
\mu_{H_1}=+\mathrm{GJS}(M_P,M_Q,\alpha)>0,
$
and
$
\sigma_{H_1}^2 \;=\; \frac{\alpha^2}{N^2}\,\sigma_{1,1}^2 \;+\; \frac{1}{n^2}\,\sigma_{2,1}^2,
$
where $\sigma_{1,1}^2$ is the $Q$-reference-side long-run variance, and $\sigma_{2,1}^2$ is the test-side long-run variance under $H_1$. 

In both cases,
\[
\frac{\sqrt{n}(\Delta\mathrm{GJS}_n-\mu_{H_\bullet})}{\sqrt{\sigma^2_{H_\bullet}}}
\;\xRightarrow[]{d}\; \mathcal{N}(0,1),
\]
where the bullet $\bullet\in\{0,1\}$ denotes the active hypothesis. 
\end{theorem}

\begin{proof}[Proof of Theorem~\ref{thm:asymptotic-normality}]
We need to establish asymptotic normality of the test statistic $\Delta \text{GJS}_n$ by performing a second-order Taylor Expansion of it and determining the asymptotic mean and asymptotic variance.

Since Lemma~\ref{GJS-exp}, adapted from \cite{zhou2018secondorderasymptoticallyoptimalstatistical}, is a purely mathematical statement about the local properties of the GJS function itself, irrespective of how its input variables are generated, this lemma is equally applicable to Markov sources.

Thus, we can obtain Taylor Expansion of Generalized Jensen Shannon Divergence when it is applied to Markov source. Consider two distinct transition matrices of two Markov sources $M_1, M_2$.  Let $\hat{M}_1$ and $\hat{M}_2$ be two other empirical transition matrices in a neighborhood of $(M_1, M_2)$. Let $\alpha$ be a fixed positive constant. The GJS divergence has the following second-order Taylor approximation around the point $(M_1, M_2)$.
 \begin{align}
   & \text{GJS}(\hat{M}_1, \hat{M}_2, \alpha) = \text{GJS}(M_1, M_2, \alpha) \nonumber\\
    &+ \sum_{s\in \mathcal{A}}\pi_1(s)\sum_{a \in \mathcal{A}} (\hat{M}_1(a|s) - M_1(a|s))\alpha\iota_1(a|s) +\sum_{s\in \mathcal{A}}\pi_2(s) \sum_{a \in \mathcal{A}} (\hat{M}_2(a|s) - M_2(a|s))\iota_2(a|s) \nonumber\\
    & +  O\left(||\hat{M}_1 - M_1||^2 + ||\hat{M}_2 - M_2||^2\right)
\end{align}
where $\pi_1$ and $\pi_2$ denote the stationary distributions of $M_1$ and $M_2$, respectively. And $\iota_1(a|s)$ and $\iota_2(a|s)$ are information densities:
\begin{align}
    \iota_1(a|s) := \iota_1((a|s)\big|M_1,M_2, \alpha)=\log\frac{(1+\alpha)M_1(a|s)}{\alpha M_1(a|s) + M_2(a|s)}\\
        \iota_2(a|s) := \iota_2((a|s)\big|M_1,M_2, \alpha) = \log\frac{(1+\alpha)M_2(a|s)}{\alpha M_1(a|s) + M_2(a|s)}
\end{align}
Furthermore, because $\Delta \text{GJS}_n=\text{GJS} \left( \hat{M}_{P}, \hat{M}_{t}, \alpha \right) - \text{GJS} \left( \hat{M}_{Q}, \hat{M}_{t},\alpha \right)$ is constructed as the difference of two GJS functions, we can directly apply the Lemma~\ref{GJS-exp} to derive the Taylor expansion $\Delta \text{GJS}_n$ itself.


First, we define the following typical set, given any $M \in \mathcal{F}_k$,. 
\begin{align}
    \mathcal{C}_n(M):=\left\{ a_{1:n}\in \mathcal{A}^n: \max_{s\in \mathcal{A}, a\in \mathcal{A}}|\hat{M}_{a_{1:n}}(a|s)-M(a|s)|\le \sqrt{\frac{\log n}{n}}\right\}
\end{align}
This is  a direct generalization of the IID  case discussed in \cite{zhou2018secondorderasymptoticallyoptimalstatistical}, and can be justified in Lemma 3.1 of \cite{wolfer2023empiricalinstancedependentestimationmarkov}, which provides a precise asymptotic analysis of the confidence interval width for estimating the transition matrix. Next we establish an upper bound on the probability of atypical sequences. We need a two-step approach: first, ensure the number of visits $N_s$ in sequence $a_{1:n}$ to each state is sufficient, and then apply a concentration inequality under that condition.

\begin{align}
    &\mathbb P\left\{ a_{1:n} \notin \mathcal{C}_n(M) \right\} = \mathbb P\left\{   \max_{s\in \mathcal{A}, a\in \mathcal{A}}|\hat{M}_{a_{1:n}}(a|s)-M(a|s)|> \sqrt{\frac{\log n}{n}} \right\} \\
    & \le \sum_{s \in \mathcal{A}} \mathbb P\left\{ \max_{a \in \mathcal{A}} |\hat{M}_{a_{1:n}}(a|s)-M(a|s)| >\sqrt{\frac{\log n}{n}} \right\} \\
    & \le \sum_{s \in \mathcal{A}} \left[\mathbb P\left\{N_s < \frac{n\pi (s)}{2} \right\} + \mathbb P\left\{\max_{a \in \mathcal{A}} |\hat{M}_{a_{1:n}}(a|s)-M(a|s)| >\sqrt{\frac{\log n}{n}} \bigg|N_s \ge \frac{n\pi (s)}{2}\right\}\right] \\
    & \le   \sum_{s \in \mathcal{A}} \left[c_1 \exp(-c_2n\pi(s))+  2k \exp(-2\frac{n\pi(s)}{2}\cdot\frac{\log n}{n})   \right] \label{hoeffding}\\
    & =    \sum_{s \in \mathcal{A}} \left[c_1 \exp(-c_2n\pi(s))+2k \cdot n^{-\pi(s)}\right] \\
    &\le k \left[c_1 \exp(-c_2n\pi(s))+2k\cdot n^{-\pi(s)}\right] \\
    &:=\tau(n,M)
\end{align}
where $\pi(s)$ denotes the stationary probability of state $s$, the first term of Eq.~\ref{hoeffding} follows Chernoff-Hoeffding inequality for Markov Chains (Corollary 8.1 of \cite{wolfer2023empiricalinstancedependentestimationmarkov}), and the second term of Eq.~\ref{hoeffding} follows McDiarmid's inequality, as its conditions of independence of variables and the bounded differences property are met. This is because the analysis is performed on the sub-problem of transitions from state $s$, conditional on the number of visits $N_s=k$ (where $k\ge \frac{n \pi(s)}{2}$), which ensures the subsequent $k$ transitions can be treated as IID samples. A similar application of this technique is detailed in \cite{wolfer2023empiricalinstancedependentestimationmarkov}. Moreover, the constant $c_1$ depends on the initial state of the chain, measuring its deviation from the steady state, while $c_2$ depends on the mixing speed of the chain, measuring how quickly it converges to its steady state. Thus,
\begin{align}
    &\mathbb P\left\{ a^P_{1:N} \notin \mathcal{C}_N (M_P) \text{\quad or \quad} a^T_{1:n} \notin \mathcal{C}_n (M_P) \text{\quad or \quad} a^Q_{1:N} \notin \mathcal{C}_N (M_Q)\right\} \\
    &\leq \mathbb P\left\{ a^P_{1:N} \notin \mathcal{C}_N (M_P) \right\}+ P\left\{ a^T_{1:n} \notin \mathcal{C}_n (M_P) \right\}+\mathbb P\left\{ a^Q_{1:N} \notin \mathcal{C}_N (M_Q)\right\} \\
    & = \tau(\alpha n,M_P)+ \tau(n,M_P)+ \tau(\alpha n,M_Q)
\end{align}
This means as long as the observed Markov chain sequences are sufficiently long, the probability of sequences being atypical can be made arbitrarily small.


Then, under $H_0$, we derive the Taylor expansion of $\Delta \text{GJS}_n=\text{GJS} \left( \hat{M}_{P}, \hat{M}_{T}, \alpha \right) - \text{GJS} \left( \hat{M}_{Q}, \hat{M}_{T},\alpha \right)$ around the true transition matrices $(M_P, M_Q)$. The first term is expanded as
\begin{align}
    &\text{GJS} \left( \hat{M}_{P}, \hat{M}_{T}, \alpha \right) = \text{GJS} (M_P,M_P,\alpha) \nonumber \\
    &+ \sum_{s\in \mathcal{A}}\pi_P(s)\sum_{a \in \mathcal{A}} (\hat{M}_P(a|s) - M_P(a|s))\alpha\iota_1(a|s) +\sum_{s\in \mathcal{A}}\pi_P(s) \sum_{a \in \mathcal{A}} (\hat{M}_T(a|s) - M_P(a|s))\iota_2(a|s) \nonumber\\
    &+   O\left(||\hat{M}_P - M_P||^2 + ||\hat{M}_T - M_P||^2\right)
\end{align}
where $\text{GJS} (M_P,M_P,\alpha)=0$, and for a given symbol $a$ and state $s$,


\begin{align}
    \iota_1(a|s)  := \iota_1((a|s)\big|M_P,M_P, \alpha)=\log \frac{(1+\alpha)M_P(a|s)}{\alpha M_P(a|s)+M_P(a|s)}=0\\
    \iota_2(a|s) := \iota_2((a|s)\big|M_P,M_P, \alpha)= \log \frac{(1+\alpha)M_P(a|s)}{\alpha M_P(a|s)+M_P(a|s)}=0
\end{align}


Thus $\text{GJS} \left( \hat{M}_{P}, \hat{M}_{T}, \alpha \right) = O\left(||\hat{M}_P - M_P||^2 + ||\hat{M}_T - M_P||^2\right)$. Then, the second term of $\Delta \text{GJS}_n$ is expanded as 
\begin{align}
   & \text{GJS} \left( \hat{M}_{Q}, \hat{M}_{T},\alpha \right) = \text{GJS} (M_Q,M_P,\alpha) \nonumber \\
    &+ \sum_{s\in \mathcal{A}}\pi_Q(s)\sum_{a \in \mathcal{A}} (\hat{M}_Q(a|s) - M_Q(a|s))\alpha\iota_1(a|s) +\sum_{s\in \mathcal{A}}\pi_P(s) \sum_{a \in \mathcal{A}} (\hat{M}_T(a|s) - M_P(a|s))\iota_2(a|s) \nonumber\\
    &+   O\left(||\hat{M}_Q - M_Q||^2 + ||\hat{M}_T - M_P||^2\right)
\end{align}
where
\begin{align}
    \iota_1(a|s)  := \iota_1((a|s)\big|M_Q,M_P, \alpha)=\log \frac{(1+\alpha)M_Q(a|s)}{\alpha M_Q(a|s)+M_P(a|s)} \label{iota_1}\\
    \iota_2(a|s) := \iota_2((a|s)\big|M_Q,M_P, \alpha)= \log \frac{(1+\alpha)M_P(a|s)}{\alpha M_Q(a|s)+M_P(a|s)} \label{iota_2}
\end{align}


Therefore, we obtain the expansion for $\Delta \text{GJS}_n$ and
\begin{align}
    &\Delta \text{GJS}_n = - \text{GJS} (M_Q,M_P,\alpha) \nonumber\\
   & -\sum_{s\in \mathcal{A}}\pi_Q(s)\sum_{a \in \mathcal{A}} (\hat{M}_Q(a|s) - M_Q(a|s))\alpha\iota_1(a|s) -\sum_{s\in \mathcal{A}}\pi_P(s) \sum_{a \in \mathcal{A}} (\hat{M}_t(a|s) - M_P(a|s))\iota_2(a|s) \nonumber  \\
   & + O\left(\frac{\log n}{n}\right) \label{delta-GJS-exp}
\end{align}

Here we connect GJS to information densities,
\begin{align}
    &\text{GJS} (M_Q,M_P,\alpha) = \alpha\text{D}_{KL}(M_Q,\frac{\alpha M_Q+M_P}{1+\alpha})+\text{D}_{KL}(M_P,\frac{\alpha M_Q+M_P}{1+\alpha})\\
    &=\alpha\sum_{s\in \mathcal{S}}\pi_Q(s)\sum_{a \in \mathcal{A}} M_Q(a|s)\log \frac{M_Q(a|s)}{\frac{\alpha M_Q(a|s)+M_P(a|s)}{1+\alpha}} +\sum_{s\in \mathcal{S}}\pi_P(s)\sum_{a \in \mathcal{A}} M_P(a|s)\frac{M_Q(a|s)}{\frac{\alpha M_Q(a|s)+M_P(a|s)}{1+\alpha}} \\
    & = \alpha\sum_{s\in \mathcal{S}}\pi_Q(s)\sum_{a \in \mathcal{A}} M_Q(a|s)\log \frac{(1+\alpha)M_Q(a|s)}{\alpha M_Q(a|s)+M_P(a|s)} +\sum_{s\in \mathcal{S}}\pi_P(s)\sum_{a \in \mathcal{A}} M_P(a|s)\frac{(1+\alpha)M_Q(a|s)}{{\alpha M_Q(a|s)+M_P(a|s)}}  \\
    & = \alpha\sum_{s\in \mathcal{S}}\pi_Q(s)\sum_{a \in \mathcal{A}} M_Q(a|s)\iota_1(a|s) + \sum_{s\in \mathcal{S}}\pi_P(s)\sum_{a \in \mathcal{A}} M_P(a|s)\iota_2(a|s) \label{GJS-info-dens}
\end{align}
where $\iota_1(a|s)$ and $\iota_2(a|s)$ are defined in Eq.~\ref{iota_1} and Eq.~\ref{iota_2}. We subsititute Eq.~\ref{GJS-info-dens} into Eq.~\ref{delta-GJS-exp} and obtain

\begin{align}
    \Delta \text{GJS}_n &= - \alpha\sum_{s\in \mathcal{A}}\pi_Q(s)\sum_{a \in \mathcal{A}} \hat{M}_Q(a|s) \iota_1(a|s) -\sum_{s\in \mathcal{A}}\pi_P(s) \sum_{a \in \mathcal{A}} \hat{M}_T(a|s) \iota_2(a|s) + O\left(\frac{\log n}{n}\right) \label{delta-gjs-2-exp}
\end{align}
Recall that $\hat{M}_Q(a|s)=\frac{N_Q(s,a)}{N_Q(s)}$, where $N_Q(s)$ is the number of occurences of state $s$ in $a^Q_{1:N}$, and $N_Q(s,a)$ the number of times $s$ is followed by $a$ in $a^Q_{1:N}$. According to Ergodic Theorem (Strong Law of Large Numbers, e.g. \cite{levin2017markov}, Theorem C.1), we consider a long Markov chain to be time-homogeneous, that is for a state $s$, we have $N_Q(s) \approx N\cdot \pi_Q(s)$. Based on this, we simplify the first term of Eq.\ref{delta-gjs-2-exp}.
\begin{align}
    \sum_{s \in \mathcal{A}} \pi_Q(s) \alpha \sum_{a \in \mathcal{A}} \hat{M}_Q(a|s)\iota_1(a|s) &= \alpha \sum_{s \in \mathcal{A}} \pi_Q(s) \sum_{a \in \mathcal{A}} \frac{N_Q(s,a)}{N_Q(s)}\iota_1(a|s) \\
    & = \frac{\alpha}{N} \sum_{s \in \mathcal{A}}\sum_{a \in \mathcal{A}}N_Q(s,a)\iota_1(a|s) \\
    & = \frac{\alpha}{N} \sum_{i=2}^{N} \iota_1(a_i^Q|a_{i-1}^Q) \label{iota-1-sum}
\end{align}
Similarly, the second term of Eq.\ref{delta-gjs-2-exp} is simplified as:
\begin{align}
    \sum_{s\in \mathcal{A}}\pi_P(s) \sum_{a \in \mathcal{A}} \hat{M}_T(a|s) \iota_2(a|s) = \frac{1}{n} \sum_{i=2}^{n} \iota_2(a^T_{i}|a_{i-1}^T) \label{iota-2-sum}
\end{align}
Combining Eq.\ref{iota-1-sum} and Eq.\ref{iota-2-sum}, we get
\begin{align}
    \Delta \text{GJS}_n &= -\frac{\alpha}{N} \sum_{i=2}^{N} \iota_1(a_i^Q|a_{i-1}^Q)-\frac{1}{n} \sum_{i=2}^{n} \iota_2(a^T_{i}|a_{i-1}^T)+ O\left(\frac{\log n}{n}\right) \label{iota-1-2}
\end{align}
Then we compute the asymptotic mean and asymptotic variance of Eq.~\ref{iota-1-2}.
By comparing Eq.~\ref{GJS-info-dens} and Eq.~\ref{delta-gjs-2-exp}, we obtain the asymptotic mean.
\begin{align}
    \mathbb{E}[\Delta \text{GJS}_n] = -\text{GJS} (M_Q,M_P,\alpha)
\end{align}

Eq.~\ref{iota-1-2} shows that the random behavior of $\Delta \text{GJS}_n$ is primarily determined by two additive functionals on Markov chains. Since the two reference sequences, $a^Q_{1:N}$ and $a^T_{1:n}$ are mutually independent, the total variance is the sum of their individual variances. 
\begin{align}
    \text{Var}(\Delta \text{GJS}_n)& = \text{Var} (-\frac{\alpha}{N} \sum_{i=2}^{N} \iota_1(a_i^Q|a_{i-1}^Q))+\text{Var}(-\frac{1}{n} \sum_{i=2}^{n} \iota_2(a^T_{i}|a_{i-1}^T)) \\
    & = \frac{\alpha^2}{N^2}\text{Var} ( \sum_{i=2}^{N} \iota_1(a_i^Q|a_{i-1}^Q)) + \frac{1}{n^2}\text{Var} (\sum_{i=2}^{n} \iota_2(a^T_{i}|a_{i-1}^T))
\end{align}

Here we use Lemma~\ref{lemma:clt-AF} and \ref{lemma:var-markov-chain} to compute the asymptotic variance for $\Delta \text{GJS}_n$. We begin by defining a new Markov chain whose state at time $i$ is given by $b_{i}:=(a^Q_{i-1},a^Q_i)$. Then we can define a function $f_1$ that acts on the state $b_i$, $f_1(b_{i})=:=\iota_1(a^Q_i|a^Q_{i-1})$. With these definitions, we have successfully converted the original sum over transitions into a sum over the states of the new chain, which perfectly fits the framework of Lemma~\ref{lemma:clt-AF} and \ref{lemma:var-markov-chain}.
\begin{align}
    \sum_{i=2}^{N} \iota_1(a^Q_i|a^Q_{i-1}) \Leftrightarrow \sum_{i=2}^{N} f_1(b_i)
\end{align}
According to Lemma~\ref{lemma:var-markov-chain}, the asymptotic variance $\sigma_1^2$ of the additive functional $\sum_{i=2}^{N}f_1(b_i)$ is given by
\begin{align}
    \sigma_{1,0}^2 = 2 \lim_{\epsilon\rightarrow 0} \langle g_{1,\epsilon},f_1 \rangle - \|f_1\|^2 \label{eq:sigma-1-formula}
\end{align}
Now we need to calculate the two main components of this formula.
The stationary distribution $\pi'$ of the new chain is determined by $\pi'=\pi_Q(s)\cdot M_Q(a|s)$.
By Eq.~\ref{GJS-info-dens}, we get \begin{align}
    \mu_1=\mathbb{E}_{\pi'}[f_1(b)]&=\sum_{(s,a) \in \mathcal{A}\times \mathcal{A}} \pi'(s,a)f_1(s,a) 
    \\
    &= \sum_{s \in \mathcal{A}}\pi_Q(s)\sum_{a\in \mathcal{A}}M_Q(a|s)\iota_1(a|s)\\
    &=\text{D}_{KL}(M_Q,\frac{\alpha M_Q+M_P}{1+\alpha})
\end{align}
We obtain the centered function
\begin{align}
    \tilde{f}_1(s,a)=f_1(s,a)-\mu_1=\iota_1(a|s)-\mu_1
\end{align}
Then according to Lemma~\ref{lemma:var-markov-chain}, we calculate the squared norm $\|\tilde{f}_1\|^2$, which is the variance of $\tilde{f}_1$ under the stationary distribution $\pi'$.
\begin{align}
    \|\tilde{f}_1\|^2= \text{Var}_{\pi'}(f_1) = \mathbb{E}_{\pi'}[(\tilde{f}_1(b))^2]=\sum_{(s,a) \in \mathcal{A}\times \mathcal{A}} \pi'(s,a)(\iota_1(a|s)-\mu_1)^2 \label{eq:tilde_f_1}
\end{align}
Calculating the inner product $\langle g_{1,\epsilon},\tilde{f}_1 \rangle$ requires first finding $g_{1,\epsilon}$ by solving the resolvent equation:
\begin{align}
    g_{1,\epsilon} = ((1+\epsilon)I-M_b)^{-1} \tilde{f}_1 \label{eq:var-framework-2}
\end{align}
where $M_b$ is the transition operator of the new chain and can be constructed from $M_Q$. Each element of the $M_b$ matrix, $M_b((s,a),(s',a'))$,  represents the probability of the new chain transitioning from state $(s, a)$ to state $(s', a')$. 
\begin{align}
    M_b((s,a),(s',a')) = 
    \begin{cases} 
    M_Q(a'|s') \text{\quad If $s'=shift(s,a)$} \\
    0 \text{\quad   \quad    \quad\quad   \quad       otherwise}
    \end{cases}
\end{align}
where $shift(s,a)$ denotes an operation that removes the first element of the sequences $s$ and appends $a$ to the end. After solving $g_{1,\epsilon}$, we compute the inner product:
\begin{align}
    \langle g_{1,\epsilon},f_1 \rangle = \sum_{(s,a)\in \mathcal{A}\times \mathcal{A}} \pi'(s,a) g_{1,\epsilon}(s,a)\tilde{f}_1(s,a)
\end{align}
We take the limit $\lim_{\epsilon \rightarrow 0}\langle g_{1,\epsilon},f_1 \rangle$, then substitute the limit and the value of Eq.~\ref{eq:tilde_f_1} into Eq.~\ref{eq:sigma-1-formula} get the final asymptotic variance $\sigma_{1,0}^2$. Similarly, we use the same method to calculate the asymptotic variance $\sigma^{2}_{2,0}=\text{Var}(\sum_{i=2}^{n} \iota_2(a^T_{i}|a^T_{i-1}))$. While the asymptotic variance does not generally admit a closed-form expression, Lemma~\ref{lemma:clt-AF} and ~\ref{lemma:var-markov-chain} provide us with constructive representations. They can be used to compute or approximate the asymptotic variance in practice.

Now we have proved that under $H_0$, the asymptotic normality of $\Delta \text{GJS}_n$, that is 
\begin{align}
    \frac{\sqrt{n}(\Delta \text{GJS}_n-\mu)}{\sigma_{H_0}}\xrightarrow{d}\mathcal{N}(0,1)
\end{align}
where $\mu_{H_0}=\mathbb{E}[\Delta \text{GJS}_n]=-\text{GJS} (M_Q,M_P,\alpha)$ and variance $\sigma_{H_0}^2=\frac{\alpha^2}{N^2}\sigma_{1,0}^2+\frac{1}{n^2}\sigma_{2,0}^2$.

Analogously, under $H_1$, we can prove  the asymptotic normality of $\Delta \text{GJS}_n$ with $\mu_{H_1}=\text{GJS} (M_P,M_Q,\alpha)$ and variance $\sigma_{H_1}^2=\frac{\alpha^2}{N^2}\sigma_{1,1}^2+\frac{1}{n^2}\sigma_{2,1}^2$, where $\sigma_{1,1}^2 = \text{Var} ( \sum_{i=2}^{N} \iota_1(a_i^P|a_{i-1}^P))$ and $\sigma_{2,1}^2 = \text{Var} (\sum_{i=2}^{n} \iota_2(a^T_{i}|a_{i-1}^T))$. As discussed in the variance framework above, they can be represented by the resolvent formulation as in Eq.~\ref{eq:sigma-1-formula} and Eq.~\ref{eq:var-framework-2}.




\end{proof}

\newpage

\section{Experiments: Configurations and More Results}
\subsection{Implementation and Configurations} \label{sec:exp-configuration}
Our implementation is adapted from MAUVE \citep{pillutla2023mauvescoresgenerativemodels} and Lastde \citep{xu2025trainingfree}. All detection experiments were conducted on one RTX 4090, while data generation ran on an A40 GPU.
We use 9 open-source models and 3 closed-source models for generating text.
Open-source models include GPT-XL \citep{radford2019gpt2}, GPT-J-6B \citep{wang2021gptj6b}, GPT-Neo-2.7B \citep{gptneo2021}, GPT-NeoX-20B \citep{black2022gptneox20b}, OPT-2.7B \citep{zhang2022opt}, Llama-2-13B \citep{touvron2023llama2}, Llama-3-8B \citep{dubey2024llama3}, Llama-3.2-3B \citep{llama32docs2024}, and Gemma-7B \citep{gemma2024}. Closed-source models include Gemini-1.5-Flash \citep{gemini15_2024}, GPT-4.1-mini \citep{openai2025gpt41}, and GPT-5-Chat \citep{openai2025gpt5chat}. 


\paragraph{Generation Pipeline} In our generation pipeline, for each dataset, we filtered out samples with text length less than 150 words and always condition only on the first 30 tokens of the human text. Each machine passage is generated between 100 and 200 tokens. After generation, we pair each human passage with its corresponding machine passage and truncate both to the shorter side (measured in words). Thus every human-machine pair used for detection has the same length and there is no systematic length advantage for either class.

\paragraph{Default Decoding Strategy}
In our experiments, unless otherwise specified, for each model family we use a fixed default decoding configuration. Concretely, for open-source models on HuggingFace we use the standard decoding configuration temperature = 1.0, top-p = 1.0, top-k = 50. For GPT-4.1-mini and GPT-5-chat (OpenAI API), we follow the default settings temperature = 1.0, top-p = 1.0 (no top-k parameter). For Gemini, we use the default settings of the Gemini API, temperature = 1.0, top-p = 0.95, top-k = 64.

\subsection{More Results}
\subsubsection{Expansion of Table~\ref{tab:opensource} and Table~\ref{tab:closesource}}

Table~\ref{tab:closesource-xsum},\ref{tab:closesource-writing},\ref{tab:closesource-squad},
\ref{tab:opensource-xsum},\ref{tab:opensource-writing}, and \ref{tab:opensource-squad} show the detection results on XSum, WritingPrompts,and SQuAD datasets. The performance is the average over three detections, where each detection is conducted on a randomly sampled test set.

\begin{table}[htbp]
\centering
{\scriptsize
\begin{tabular}{lcccc}
\toprule
& Gemini-1.5-Flash & GPT-4.1-mini & GPT-5-Chat & Avg  \\
\midrule
Likelihood & 53.2 \com 1.31 & 55.54 \com 1.09 & 43.03 \com 2.69 &  50.59 \\

LogRank & 52.01 \com 2.53 & 57.96 \com 2.81 & 45.86 \com 3.88 & 51.94 \\

Entropy & 63.19 \com 1.78 & 51.7 \com 1.02 & 56.8 \com 2.02 &  57.23\\

DetectLRR & 49.85 \com 2.54 & 62.26 \com 0.91 & 54.14 \com 3.6 & 55.42 \\

Lastde & 59.26 \com 3.39 & 55.97\com 2.18& 45.3 \com 1.34 & 53.51 \\

Lastde++ & \textbf{76.9} \com 1.62& 69.29 \com 2.00 & 48.14 \com 3.28& 64.78 \\

DNA-GPT &60.85 \com 1.41 & 55.7 \com 0.46 & 45.4 \com 0.77 &  53.98\\

Fast-DetectGPT & \underline{75.52} \com 1.58 & 66.7 \com 1.45 & 48.51 \com 2.01 & 63.58 \\

DetectGPT &62.58 \com 1.31 & 61.25 \com 3.08& 50.17 \com 0.29 & 58 \\

DetectNPR & 58.77 \com 2.47 & 62.17 \com 1.50 & 53.32 \com 0.97 & 58.09\\

R-Detect & 63.68  \com 0.77 & 63.43 \com 2.31 & 58.74 \com 1.62 & 61.95\\

Binoculars & 74.84  \com 2.12 & 61.12 \com 1.47 & 45.94 \com 0.67 & 60.63\\

FourierGPT & 52.06  \com 0.39 & 55.53 \com 2.31 & 61.1 \com 1.1 &  56.23\\

\rowcolor{ours}$\text{SurpMark}_{k=6}$ & 70.24 \com 0.77 & \underline{84.07} \com 2.21 & 84.16 \com 1.01 & 79.49 \\

\rowcolor{ours}$\text{SurpMark}_{k=7}$ & 71.22 \com 0.32 & 82.52 \com 1.11 & \textbf{87.02} \com 1.4 & \underline{80.25}\\
\rowcolor{ours}$\text{SurpMark}_{k=8}$ & 69.03 \com 1.74 & \textbf{85.78} \com 0.76 & \underline{86.38} \com 0.94 & \textbf{80.40}\\
\bottomrule
\end{tabular}}
\caption{Detection results on XSum for text generated by 3 closed-source models under the black-box setting.}
\label{tab:closesource-xsum}
\end{table}

\begin{table}[htbp]
\centering
{\scriptsize
\begin{tabular}{lcccc}
\toprule
& Gemini-1.5-Flash & GPT-4.1-mini & GPT-5-Chat & Avg  \\
\midrule
Likelihood & 80.53 \com 1.29 & 82.95 \com 1.23 & 62.00 \com 2.95 & 75.16 \\

LogRank & 74.73 \com 2.64 & 80.66 \com 2.81 & 58.01 \com 4.04 & 71.13 \\

Entropy & 46.34 \com 3.11 & 19.00 \com 6.43 & 25.23 \com 4.08 & 30.19 \\

DetectLRR & 48.22 \com 2.7 & 68.50 \com 1.06& 43.92 \com 2.48 & 53.55 \\

Lastde & 41.09 \com 2.88 & 55.72 \com 2.62& 30.64 \com 1.59 & 42.48 \\

Lastde++ & 76.90 \com 1.05 & 68.49 \com 2 & 30.64\com 3.23 & 58.68 \\

DNA-GPT & 78.19 \com 0.87 & 63.70 \com 1.73 & 45.60 \com 3.2 & 62.50 \\

Fast-DetectGPT & \underline{91.96} \com 0.31 & 70.23 \com 1.91& 30.01 \com 4.07 & 64.07 \\

DetectGPT & 87.12 \com 0.49& 78.04 \com 0.9 & 58.72 \com 2.01  & 74.63 \\

DetectNPR & 80.47  \com 1.23 & 75.80 \com 0.97 & 55.97 \com 2.31& 70.75 \\

R-Detect &83.31  \com 0.89 & 78.79 \com 1.92 & 77.06 \com 0.48& 79.72 \\

Binoculars & \textbf{95.35}  \com 0.1 &  80.55\com 0.34 & 42.26\com 0.67 & 72.72 \\

FourierGPT &  77.8 \com 0.36 &  77.96 \com 1.05 & 74.45 \com 1.72 & 76.74 \\

\rowcolor{ours}$\text{SurpMark}_{k=6}$ & 86.64 \com 2.33 & \underline{85.80} \com 0.57 & \underline{82.25} \com 1.03 & \underline{84.90} \\

\rowcolor{ours}$\text{SurpMark}_{k=7}$ & 86.68 \com 1.4 & 83.64 \com 0.33 & 83.73 \com 0.52 & 84.68 \\

\rowcolor{ours}$\text{SurpMark}_{k=8}$ & {89.43} \com 0.35 & \textbf{87.27} \com 0.14 & \textbf{83.56} \com 0.67 & \textbf{86.75} \\
\bottomrule
\end{tabular}}
\caption{Detection results on WritingPrompts for text generated by 3 closed-source models under the black-box setting.}
\label{tab:closesource-writing}
\end{table}
\begin{table}[htbp]
\centering
{\scriptsize

\begin{tabular}{lcccc}
\toprule
& Gemini-1.5-Flash & GPT-4.1-mini & GPT-5-Chat & Avg  \\
\midrule
Likelihood & 35.74 \com 3.46 & 61.82 \com 3.21 & 43.83 \com 2.01  & 47.13 \\

LogRank & 34.86 \com 2.61 & 61.78 \com 3.52 & 45.62 \com 3.66 & 47.42 \\

Entropy & 65.55 \com 1.08 & 45.46 \com 1.43 & 58.94 \com 0.65 & 56.65 \\

DetectLRR & 35.46 \com 1.84 & 59.10 \com 2.11 & 51.42 \com 2.50 & 48.66 \\

Lastde & 44.03 \com 1.55 & 60.15 \com 2.92 & 49.95 \com 3.65  & 51.38 \\

Lastde++ & 52.47 \com 1.86& 66.90 \com 2.18 & 51.76 \com 3.02 & 57.04 \\

DNA-GPT & 47.15 \com 0.93 & 50.74 \com 2.88 & 58.45 \com 1.18 & 52.11 \\

Fast-DetectGPT & 49.98 \com 1.33 & 68.04 \com 1.19 & 51.64 \com 1.98 & 56.55 \\

DetectGPT &57.87 \com 2.65 & {70.95} \com 0.82 & 54.90\com 0.83 & 61.24 \\

DetectNPR & 55.63 \com 2.91& \textbf{74.53} \com 1.29 & 55.67 \com 2.13 & 61.94 \\

R-Detect & 60.86 \com 1.33& {72.69} \com 1.41 & 67.45 \com 2.37 & 67 \\

Binoculars &  53.34 \com 2.53 & \underline{73.69}\com 0.55 & 60.76 \com 0.67 & 62.6 \\

FourierGPT &   53.89 \com 2.57 &  55.66 \com 2.25 & 58.92 \com 2.24 & 56.16\\

\rowcolor{ours}$\text{SurpMark}_{k=6}$ & \underline{66.84} \com 1.11 & 70.87 \com 0.86 & 68.57 \com 1.48& 68.76 \\

\rowcolor{ours}$\text{SurpMark}_{k=7}$ & \textbf{67.51} \com 1.3 & 69.27 \com 1.83 & \underline{73.23} \com 0.87 & \textbf{70.00} \\

\rowcolor{ours}$\text{SurpMark}_{k=8}$ & 59.53 \com 1.49 & 72.27 \com 1.32 & \textbf{74.81} \com 1.02 & \underline{68.87} \\
\bottomrule
\end{tabular}}
\caption{Detection results on SQuAD for text generated by 3 closed-source models under the black-box setting. }
\label{tab:closesource-squad}
\end{table}

\begin{table}[htbp]
\centering
{
\tiny 
\resizebox{\textwidth}{!}{
\begin{tabular}{lcccccccccc}
\toprule
& GPT2-XL & GPT-J-6B & GPT-Neo-2.7B & GPT-NeoX-20B & OPT-2.7B & Llama-2-13B & Llama-3-8B & Llama-3.2-3B & Gemma-7B  & Avg \\
\midrule
Likelihood & 76.5 \com 0.63 & 62.74 \com 1.07 & 58.36 \com 1.62 & 60.58 \com 1.8  & 68.51 \com 1.37 & 92.22 \com 0.48 & 93.41 \com 0.82 & 51.61 \com 0.62 & 55.13 \com 1.18 & 68.78\\

LogRank & 80.16 \com 0.89 & 67.83 \com 1.13 & 64.54 \com 0.98 & 63.58 \com 1.25 & 72.33 \com 1.56 & 94.56 \com 0.32 & 95.05 \com 0.17 & 59.35 \com 0.08 & 59.13 \com 0.68 & 76.89 \\

Entropy & 59.65 \com 1.52 & 56.37 \com 0.66 & 63.76 \com 1.43 & 55.32 \com 1.11 & 52.88 \com 0.68 & 42.33 \com 2.58 & 29.31 \com 3.19 & 55\com 2.89 & 53.2\com 1.48 & 50.40 \\

DetectLRR & 83.2 \com 0.83 & 76.5 \com 0.88 & 76.94 \com 1.09 & 68.4 \com 1.35 & 77.49 \com 0.54 & 95.74 \com 0.23 & 94.85 \com 0.08  & \textbf{75.05} \com 0.31 & 66.42 \com 1.42 & 81.42 \\

Lastde & 91.97 \com 0.44 & 77.99 \com 0.89 & 82.49 \com 0.85 & 72.12 \com 1.63 & 77.85 \com 0.68 & 92.01 \com 0.89 & 94.29 \com 0.38  & 59.52 \com 0.05& 61.09\com 1.27 & 82.57 \\

Lastde++ & \textbf{98.99} \com 0.21 & 85.38\com 0.63 & 87.5\com 0.11 & 80.3 \com 0.92 & 87.93 \com 0.54 & 92.52 \com 0.43 & 95.9 \com 0.14 & 59.9 \com 0.08 & 65.68 \com 0.97 & 87.51 \\

DNA-GPT & 71.43 \com 1.33 & 55.47 \com 2.85 & 54.43 \com 3.2 & 56.31 \com 1.86 & 58.2 \com 1.72 & 93.69 \com 0.36 & 96.54 \com 0.12 & 50.37 \com 0.07 & 55.29 \com 1.04 & 70.70 \\

Fast-DetectGPT & 95.54 \com 0.34 & 78.6 \com 0.56 & 81.84\com 0.88 & \textbf{83.76} \com 1.28 & 90.55 \com 0.77 & \textbf{97.77}\com 0.05 & 96.78 \com 0.21 & 61.86 \com 1.42 & 63.2 \com 1.18 & 84.71 \\

DetectGPT & 92.88 \com 1.3 & 71.86 \com 1.79 & 76.67 \com 2.01 & 78.06 \com 0.87 & 82.88 \com 1.23 & 82.79 \com 0.62 & 83.61 \com 1.25 & 56.06\com 2.65 & 61.6 \com 2.94 & 77.18 \\

DetectNPR & 91.87 \com 1.13 & 72.36 \com 1.46 & 78.83 \com 0.66 & 76.76 \com 1.48 & 84.06 \com 1.21 & 94.29 \com 0.86 & 92.31  \com 0.3& 59.62 \com 1.77 & 60.52 \com 1.78 & 80.05 \\

R-Detect & 72.87 \com 1.49 & 59.86 \com 1.11 & 67.59 \com 0.48 & 63.45 \com 2.45 &  69.75 \com 0.71 & 72.11 \com 0.93 & 81.06  \com 0.84 & 62.43 \com 0.82 & 46.75 \com 0.73 & 66.21  \\

Binoculars &  \underline{98.87} \com 0.13 &  74.66 \com 0.48 &  78.05\com 1.27 & 76.18\com 1.22 & 79.89 \com 0.79 &  96.78 \com 0.21 & 96.19  \com 0.16& 48.22 \com 0.71 & 63.71 \com0.72 & 79.17 \\

FourierGPT & 51.8  \com 1.39 &  52.52 \com 2.02 &  50.44 \com 2.96 &  59.17 \com 0.28 & 48.16 \com 3.01 & 63.38 \com 2.42 & 59.74  \com 3.4 & 51.98 \com 1.73 & 53.62 \com 0.78 & 54.53 \\

\rowcolor{ours}$\text{SurpMark}_{k=6}$ & 96.95 \com 0.43 & \underline{88.35} \com 1.02 & \underline{92.26} \com 0.65  & 81.58 \com 0.72 & \underline{90.88} \com 0.1 & 96.87 \com 0.26 & \textbf{97.77} \com 0.35 & {73.96} \com 0.86 & \textbf{73.01} \com 0.98 & \underline{87.96} \\

\rowcolor{ours}$\text{SurpMark}_{k=7}$ & {97} \com 0.8 & \textbf{89.26} \com 0.48 & \textbf{92.92} \com 0.06 & \underline{82.45} \com 1.03 & \textbf{91.16} \com 1.08 & \underline{97.09} \com 0.45  & \underline{97.48} \com 0.31 & \underline{73.07} \com 0.6 & \underline{72.97} \com 0.85 & \textbf{88.16} \\
\rowcolor{ours}$\text{SurpMark}_{k=8}$ & 95.55 \com 0.21 & 85.49 \com 0.63 & 88.33 \com 0.83 & 82.35 \com 0.49 & 90.19 \com 0.41 & 96.83 \com 0.16 & 97.24\com 0.08 & 72.92 \com 1.02 & 70.11 \com 0.98 & 86.56 \\
\bottomrule
\end{tabular}}}
\caption{Detection results on XSum for text generated by 9 open-source models under the black-box setting.}
\label{tab:opensource-xsum}
\end{table}

\begin{table}[htbp]
\centering
{
\tiny 
\resizebox{\textwidth}{!}{
\begin{tabular}{lcccccccccc}
\toprule
& GPT2-XL & GPT-J-6B & GPT-Neo-2.7B & GPT-NeoX-20B & OPT-2.7B & Llama-2-13B & Llama-3-8B & Llama-3.2-3B & Gemma-7B  & Avg \\
\midrule
Likelihood & 94.55 \com 0.63 & 88.73 \com 1.11 & 89.67 \com 0.84 & 87.12 \com 1.13 & 85.15 \com 2.55 & 99.48 \com 0.2 & 99.61 \com 0.08& 85.95 \com 0.35 & 83.16 \com 1.45 & 90.38 \\

LogRank & 96.04\com 0.43& 91.78 \com 1.18 & 92.20 \com 1.22 & 89.68 \com 0.57 & 89.96 \com 0.62 & 99.59 \com 0.01 & 99.81 \com 0.11 & 89.09 \com 1.05 & 86.00 \com 0.86 & 92.68 \\

Entropy & 34.72 \com 2.75& 33.64 \com 2.81 & 32.82 \com 2.13 & 32.63 \com 1.74  & 40.88 \com 2.17 & 5.83 \com 3.74 & 8.42 \com 4.86 & 53.00 \com 2.55 & 37.16 \com 2.4 & 31.01 \\

DetectLRR & 96.96 \com 0.31 & 95.31 \com 0.42 & 94.85 \com 0.16 & 92.03 \com 0.32 & 95.68 \com 0.64 & 98.57 \com 0.12 & 99.81  \com 0.03 & 92.44 \com 0.17 & 89.19 \com 0.03 & 94.98 \\

Lastde & 98.50 \com 0.2 & 93.94 \com 0.12 & 95.97 \com 0.33 & 90.36 \com 0.82& 96.05 \com 0.18 & 97.97\com 0.48 & 98.69 \com 0.23 & 92.04 \com 0.1 & 84.96 \com 0.56  & 94.28 \\

Lastde++ & \underline{99.68} \com 0.11 & 95.96 \com 0.51 & 98.86 \com 0.1 & 92.68 \com 0.74 & \textbf{98.39} \com 0.12 & 99.14 \com  0.08 & 99.56 \com 0.06 & \textbf{95.04} \com 0.3 & \textbf{92.59} \com 0.65 & \textbf{96.88} \\

DNA-GPT & 90.53 \com 1.62 & 85.34 \com 1.13 & 85.72 \com 0.7 & 83.01 \com 1.41 & 85.05 \com 1.29 & 98.88 \com 0.12 & 99.65 \com 0.03 & 84.47 \com 0.65 & 80.60 \com 0.81 & 88.14 \\

Fast-DetectGPT & 99.67 \com 0.02 & 93.80 \com 0.6 & 96.62 \com 0.31 & 92.22 \com 0.27 & 94.99 \com 0.52 & \underline{99.56} \com 0.01 & 99.84 \com 0.04 & 93.55 \com 0.53 & 89.36\com 1.03& 95.51 \\

DetectGPT & 95.88 \com 0.2 & 85.83 \com 1.15 & 91.12 \com 1.52 & 85.17 \com 1.84 & 90.13 \com 1.21  & 92.67 \com 0.63& 93.10 \com0.61 & 80.08 \com 1.07 & 83.10 \com 2.3 & 88.56 \\

DetectNPR & 98.29 \com 0.2 & 89.77 \com 0.33 & 93.02 \com 0.92 & 87.96 \com 0.55 & 92.36 \com 1.43 & 98.20 \com 0.51 & 98.52 \com 0.18 & 85.22\com 0.5 & 86.71 \com 1.03 & 92.23 \\

R-Detect &  86.68 \com 1.35 &  75.93\com 1.06 & 75.23 \com 0.59 & 73.83 \com 1.1 &   51.03\com 2.57 & 79.69 \com 0.88 & 82.79  \com 0.93& 71.2 \com 2.36 & 72.62 \com 0.89 & 74.33 \\

Binoculars & 99.6 \com 0.03 & 93.7 \com 0.51 &  94.96 \com 0.21 & 93.22 \com 0.21 & 91.33 \com 0.86 &  98.9 \com 0.16 & 99  \com 0.06& 93.4 \com 0.27 & 89.22 \com 0.82 & 94.81 \\

FourierGPT & 60.23 \com 4.8 & 59.81 \com 1.62 &68.08  \com 1.46 & 60.6 \com 0.29 & 56.95\com 3.04 & 91.4 \com 0.74 &  91.61 \com 1.14 & 58.68 \com 1.28 & 61.52 \com 0.72 & 67.65 \\

\rowcolor{ours}$\text{SurpMark}_{k=6}$ & 99.44 \com 0.06 & \textbf{97.60}  \com 0.22 & \textbf{98.32} \com 0.57 & \textbf{94.38} \com 0.16 & \underline{97.22} \com 0.16 & 99.47 \com 0.07 & 99.65 \com 0.1 & 92.71 \com 1.45 & 89.28  \com 1.69 & \underline{96.45} \\

\rowcolor{ours}$\text{SurpMark}_{k=7}$ & 99.27 \com 0.12 & \underline{97.29} \com 0.61 & \underline{97.63} \com 0.17 & \underline{94.31} \com 0.12 & 96.79 \com 0.52 & 99.53 \com 0.06 & \underline{99.86} \com 0.02 & \underline{93.61} \com 0.41 & 89.42 \com 0.95 & 96.41 \\

\rowcolor{ours}$\text{SurpMark}_{k=8}$ & \textbf{99.9} \com 0.01 & 96.85 \com 1.06 & 97.61 \com 0.38 & 93.93 \com 0.24 & 96.48 \com 0.4 & \textbf{99.59} \com 0.03 & \textbf{99.87} \com 0.03 & 91.65\com 0.37 & \underline{90.37} \com 1.43 & 96.25 \\
\bottomrule
\end{tabular}}}
\caption{Detection results on WritingPrompts for text generated by 9 open-source models under the black-box setting.}
\label{tab:opensource-writing}
\end{table}

\begin{table}[htbp]
\centering
{
\tiny 
\resizebox{\textwidth}{!}{
\begin{tabular}{lcccccccccc}
\toprule
& GPT2-XL & GPT-J-6B & GPT-Neo-2.7B & GPT-NeoX-20B & OPT-2.7B & Llama-2-13B & Llama-3-8B & Llama-3.2-3B & Gemma-7B  & Avg \\
\midrule
Likelihood & 84.00 \com 2.33 & 73.00 \com 3.12 & 71.93 \com 2.95 & 68.40 \com 1.32 & 78.01 \com 1.25 & 91.47 \com 1.43 & 88.77 \com 1.01 & 58.11 \com 1.86 & 59.10 \com 1.58 & 74.75 \\

LogRank & 88.39 \com 2.06 & 78.14 \com 0.96 & 78.13 \com 2.26 & 72.85 \com 1.45 & 83.68 \com 1.2& 93.55\com 0.59 & 90.48 \com 1.3 & 64.69\com 0.64 & 62.41 \com 1.72 & 79.15 \\

Entropy & 58.93 \com 3.11 & 51.43 \com 2.6 & 56.24 \com 2.91 & 49.86 \com 1.68 & 52.88 \com 3.1 & 38.92 \com 2.37& 38.72 \com 2.71& 51.00 \com 2.26 & 50.18 \com 1.82 & 49.80 \\

DetectLRR & 93.05 \com 0.11 & 85.61 \com 1.24 & 89.56 \com 1.01 & 80.38 \com 1.19 & 92.28 \com 1.05 & 94.98 \com 0.35 & 91.47 \com 1.45 & 77.14 \com 1.09 & 70.89 \com 2.31& 86.15 \\

Lastde & 97.45 \com 0.37 & 85.71 \com 1.45 & 88.82 \com 0.44 & 78.01 \com 1.87 & 92.78 \com 1.18 & 89.88 \com 1.03 & 90.89 \com 0.72 & 67.41 \com 2.9 & 62.40 \com 2.55 & 83.71 \\

Lastde++ & \textbf{99.72} \com 0.05 & \textbf{93.27} \com 0.42 & \textbf{96.51} \com 0.05 & 82.42 \com 0.3  & 96.13 \com 0.21 & 94.85 \com 0.14 & 94.72 \com 0.02 & \underline{77.47} \com 0.32 & \textbf{72.43}\com 0.24 & \textbf{89.72} \\

DNA-GPT & 83.97 \com 2.21 & 71.23 \com 2.17 & 78.21\com 1.45  & 71.93 \com 1.86 & 78.33 \com 1.43  & 95.15\com 0.49 & \underline{95.00} \com 0.32  & 59.52 \com 1.61 & 60.06 \com 1.67 & 77.04 \\

Fast-DetectGPT & 98.60 \com 0.05 & 88.09 \com 1.05 & 89.00 \com 1.18 & 81.79 \com 1.58 & 92.89 \com 0.6 & \textbf{97.32} \com 0.28 & \textbf{97.32} \com 0.05 & 67.56 \com 2.47 & 69.29 \com 0.61 & 86.87 \\

DetectGPT & 94.59 \com 0.43 & 80.95 \com 2.04  & 86.34 \com 1.21  & 69.04 \com 2.6  & 80.45 \com 2.84 & 84.08 \com 1.65 & 82.13 \com 1.72 & 56.56 \com 3.7 & 62.44 \com 1.54 & 77.40 \\

DetectNPR & 94.64 \com 0.26 & 83.59 \com 1.24& 87.34 \com 1.29& 75.01\com 2.13  & 83.07 \com 1.78 & 93.09 \com 0.69 & 90.18 \com 1.05 & 63.52 \com 2.43 & 67.25 \com 1.7 & 81.97 \\

R-Detect & 63.58 \com 0.97 & 55.04 \com 0.64 & 60.28 \com 1.67 & 52.77 \com 2.64 &  51.03 \com 0.72 & 88.15 \com 0.69 &  81.06 \com 0.87 & 53.03 \com 2.77 & 47.02 \com 3.4 & 61.33 \\

Binoculars & 99.09 \com 0.04 & 88.91 \com 1.03 &  89.49 \com 0.46 & 76.66 \com 1.21 & 89.49 \com 0.27 & 95.1 \com 0.02 & 94.04  \com 0.3& 63.46  \com 0.46 & 67.77 \com 2.58 & 84.89\\

FourierGPT &  52.12\com 3.12 & 50.5 \com 2.56 & 56.5 \com 2.79 & 49.76 \com 1.82 &  52.3 \com 2.61 & 62.49 \com 0.86 &  64.82 \com 1.22& 53.83 \com 0.72 & 52.38 \com 2.49 & 54.97 \\

\rowcolor{ours}$\text{SurpMark}_{k=6}$ & 97.88 \com 0.55 & \underline{92.93} \com 0.82  & 94.99 \com 0.3 & \textbf{84.39}\com 0.18 & 95.37 \com 0.6 & 95.89 \com 0.49 & 93.76 \com 0.35 & \textbf{78.54}\com 1.97 & \underline{69.92} \com 0.54 & \underline{89.30} \\

\rowcolor{ours}$\text{SurpMark}_{k=7}$ & \underline{98.77} \com 0.72 & 92.74 \com 0.45 & \underline{95.72} \com 0.38 & \underline{82.45} \com 1.03 & \underline{96.68} \com 0.65 & \underline{96.13} \com 0.3 & 94.17 \com 0.57 & 75.55 \com 1.21 & 68.27 \com 0.95 & 88.94 \\

\rowcolor{ours}$\text{SurpMark}_{k=8}$ & 98.76 \com 0.66 & 90.78 \com 0.23 & 94.56\com 0.1 & 79.36 \com 1.67 & \textbf{97.26} \com 0.21 & 94.81 \com 0.41 & 93.32 \com 0.16& 76.55 \com 1.2 & 67.47 \com 0.83 & 88.10 \\
\bottomrule
\end{tabular}}}
\caption{Detection results on SQuAD for text generated by 9 open-source models under the black-box setting.}
\label{tab:opensource-squad}
\end{table}
\newpage 
\subsubsection{Empirical Calibration of the Bin-count Scaling Constant}\label{app:choice-of-k}

Our intention in Section 4.2 is justify the scaling law $k=\Theta(N^{1/5})$. In practice, for each dataset we treat the theorem as providing the functional form $k=CN^{1/5}$ and then select $k$ by a small grid search. To handle constant $C$, we examined the ratio $\frac{k}{N^{1/5}}$ across several reference size and found it to be consistently around 0.8. This suggests that in our regime the implicit constant is approximately $C \approx 0.8$, and that the empirically chosen $k$ is well aligned with the theoretical scaling law.


\begin{table}[htbp]
\centering{\small
\begin{tabular}{lcccc}
\toprule
Number of ref samples 
& $N$ (approx.\ total length) 
& Empirical best $k$ 
& $N^{1/5}$ 
& $\frac{k}{N^{1/5}}$ \\
\midrule
100  & 15{,}000  & 6 & 6.84  & 0.88 \\
300  & 45{,}000  & 7 & 8.52  & 0.82 \\
400  & 60{,}000  & 7 & 9.03  & 0.78 \\
600  & 90{,}000  & 7 & 9.80  & 0.71 \\
900  & 135{,}000 & 9 & 10.62 & 0.85 \\
\bottomrule
\end{tabular}}
\caption{Scaling of the empirically optimal number of bins $k$ with  $N$.}
\label{tab:k-vs-N}
\end{table}

\subsubsection{Empirical Validation of the Asymptotic Normal Approximation} \label{sec:shapiro-wilk}
While the asymptotic variance in Theorem 4.4 does not provide a simple closed-form expression, Appendix~\ref{sec:proof-asymp-gaussian} along with Lemma~\ref{lemma:var-markov-chain} give an explicit numerical procedure to solve it. To quantitatively compare this theoretical variance with empirical fluctuations, we proceed as follows. We first compute the theoretical variance using the estimated Markov kernels from reference data. Then we estimate the empirical variance of $\Delta \text{GJS}_n$ in detection procedure. Table~\ref{tab:var-theory-vs-emp} reports theoretical variance and empirical variance with test length 250. Overall, the theoretical variance captures the right order of magnitude $\Delta \text{GJS}_n$ fluctuations, so we interpret it as a conservative asymptotic scale parameter rather than a precise finite-sample variance estimator.

\begin{table}[htbp]
\centering{\tiny
\begin{tabular}{lcccccc}
\toprule
Model 
& Emp $\sigma^2$ (Human) 
& Emp $\sigma^2$ (LM) 
& Th $\sigma^2$ (Human) 
& Th $\sigma^2$ (LM) 
& Ratio Th/Emp (Human) 
& Ratio Th/Emp (LM) \\
\midrule
Llama3-8B   & $1.15\times 10^{-5}$ & $1.06\times 10^{-5}$ & $2.48\times 10^{-5}$ & $7.50\times 10^{-5}$ & 2.16 & 7.08 \\
Llama3.2-3B & $1.12\times 10^{-5}$ & $9.93\times 10^{-6}$ & $2.73\times 10^{-5}$ & $9.11\times 10^{-5}$ & 2.44 & 9.17 \\
Gemma-7B    & $1.50\times 10^{-6}$ & $5.96\times 10^{-7}$ & $3.56\times 10^{-6}$ & $2.43\times 10^{-6}$ & 2.37 & 4.08 \\
\bottomrule
\end{tabular}}
\caption{Comparison between empirical and theoretical variances of $\Delta \mathrm{GJS}$ under human and LM text.}
\label{tab:var-theory-vs-emp}
\end{table}
Finally, to assess the distributional shape, we ran Shapiro-Wilk tests on the obtained $\Delta \text{GJS}_n$ score, as shown in Table~\ref{tab:shapiro-wilk}, the Shapiro-Wilk statistics are close to 1 and the p-values are not small (larger than 0.05). This indicates no evidence against normality and empirically supports the central-limit-theorem-based approximation in Theorem~\ref{sec:proof-asymp-gaussian}, consistent with the variance comparison above.

\begin{table}[t]
\centering{\scriptsize
\begin{tabular}{lccc}
\toprule
Setting 
& SQuAD@GPT-5-chat 
& WritingPrompts@Llama3-8B 
& XSum@Qwen3-8B \\
\midrule
{$H_1$ (LM text)} stat     & 0.9952 & 0.9856 & 0.9974 \\
{$H_1$(LM text)} p-value  & 0.9078 & 0.1203 & 0.9969 \\
{$H_0$ (human text)} stat  & 0.9876 & 0.9854 & 0.9929 \\
{$H_0$ (human text)} p-value & 0.2032 & 0.1143 & 0.6632 \\
\bottomrule
\end{tabular}}
\caption{Shapiro-Wilk test statistics and p-values for $\Delta \mathrm{GJS}_n$ under LM-generated ($H_1$) and human ($H_0$) text.}
\label{tab:shapiro-wilk}
\end{table}

\subsubsection{Score Distribution}
\begin{figure}[htbp]
\centering
\includegraphics[width=0.3 \textwidth]{xsum-llama-3-8b-score-distribution.png}
\includegraphics[width=0.3\textwidth]{writing-llama-3-8b-score-distribution.png}
\includegraphics[width=0.3\textwidth]{writing-llama-3-3.2b-score-distribution.png}
\includegraphics[width=0.3\textwidth]{writing-gemma-7b-score-distribution.png}
\includegraphics[width=0.3\textwidth]{squad-opt-2.7b-score-distribution.png}
\includegraphics[width=0.3\textwidth]{hc3-psy-score-distribution.png}
\caption{SurpMark's score distribution.}
\label{fig:appendix-score-distribution}
\end{figure}

\subsubsection{Effect of Test Length} 
\begin{figure}[htbp]
\centering
\includegraphics[width=0.4 \textwidth]{writing_llama3_8b_auroc_vs_n.png}
\includegraphics[width=0.39\textwidth]{writing_gptneox20b_auroc_vs_n.png}
\caption{AUROC vs test length.}
\label{fig:appendix-test-length}
\end{figure}

\newpage
\subsubsection{TPR} \label{sec:tpr}
In Table~\ref{tab:tpr-fpr}, we include TPR@FPR=1\% and 5\% for SurpMark and two strong baselines (Lastde++ and Fast-DetectGPT) across evaluation settings. Overall, these results indicate that SurpMark is particularly effective in the low-false-positive regime.

\begin{table}[htbp]
\centering
{\tiny
\begin{tabular}{l l c c c c c c }
\toprule
Method          &            & 
\makecell{XSum@\\GPT-5-Chat} & 
\makecell{WritingPrompts@\\GPT-4.1-mini} & 
\makecell{XSum@\\Llama2-13B} & 
\makecell{SQuAD@\\GPT-Neo-2.7B} & 
\makecell{WritingPrompts@\\Llama3-8B} & 
\makecell{HC3-Chi-\\Psy}  \\
\midrule
Lastde++        & TPR@FPR=1\%    & 4.00            & 6.00                         & 76.67           & \textbf{53.33}     & 97.33                     & 22.00        \\
       & TPR@FPR=5\%    & 12.67           & 18.67                        & 82.67           & 81.33              & 99.33                     & 31.33     \\
Fast-DetectGPT  & TPR@FPR=1\%    & 2.00            & 3.30                         & \textbf{80.67}  & 47.33              & 94.67                     & 26.00        \\
  & TPR@FPR=5\%    & 4.00            & 22.00                        & 86.67           & 77.33              & 98.00                     & 40.00       \\
SurpMark        & TPR@FPR=1\%    & \textbf{31.33}  & \textbf{31.33}               & 75.33           & 41.33              & \textbf{100.00}           & \textbf{90.00}  \\
       & TPR@FPR=5\%    & \textbf{37.33}  & \textbf{50.00}               & \textbf{90.00}  & \textbf{90.00}     & \textbf{100.00}           & \textbf{97.33}  \\
\bottomrule
\end{tabular}
\caption{TPR at fixed FPR levels (1\% and 5\%) for different detectors and datasets.}
\label{tab:tpr-fpr}}
\end{table}


% \subsubsection{Cross-Domain Generalization} \label{app:cross-domain}

% \begin{table}[htbp]
% \centering
% {\tiny
% \begin{tabular}{l c c c c}
% \toprule
% Test & self-ref & WritingPrompts-as-ref & XSum-as-ref & SQuAD-as-ref \\
% \midrule
% XSum@Llama2-13b              & 97.09 & 97.22 & --   & 96.45 \\
% WritingPrompts@Llama2-13b    & 99.53 & --    & 99.71 & 99.33 \\
% SQuAD@Llama2-13b             & 96.13 & 94.36 & 95.67 & -- \\
% XSum@Llama3-8b               & 97.09 & 97.25 & --    & 95.13 \\
% WritingPrompts@Llama3-8b     & 99.86 & --    & 99.95 & 99.88 \\
% SQuAD@Llama3-8b              & 94.17 & 93.77 & 93.90 & -- \\
% XSum@Gemma\_7b               & 72.97 & 68.32 & --    & 73.81 \\
% WritingPrompts@Gemma\_7b     & 89.42 & --    & 86.60 & 90.00 \\
% SQuAD@Gemma\_7b              & 68.27 & 72.92 & 70.44 & -- \\
% XSum@GPT-Neox-20b            & 82.45 & 82.52 & --    & 81.80 \\
% WritingPrompts@GPT-Neox-20b  & 94.31 & --    & 93.45 & 92.48 \\
% SQuAD@GPT-Neox-20b           & 82.45 & 81.43 & 83.02 & -- \\
% \bottomrule
% \end{tabular}}
% \caption{AUROC of SurpMark under different reference choices across datasets and models.}
% \label{tab:ref-choice2}
% \end{table}


% \subsubsection{Cross-Generator Generalization} \label{app:cross-generator}
% To evaluate robustness without recomputing transitions per model, we propose SurpMark-MC, a multi-class variant using a single shared quantizer (derived from pooled corpora) and distinct Markov chains for human and each reference generator. We classify test texts by assigning them to the source with the minimal GJS divergence (i.e., a nearest-neighbor rule). In a multi-class setting with GPT-J-6B, GPT-4.1-mini, and LLaMA2-13B, SurpMark-MC achieves 82.3\% overall accuracy, with 78.7\% accuracy for human texts, 78.0\% for GPT-J-6B, 84.0\% for GPT-4.1-mini, and 96.0\% for LLaMA2-13B.  We further collapse the multi-generator setting into a binary classification task (LM-generated vs.\ human-written), defining the score as the difference between the GJS to the human reference and the minimum GJS to any machine reference. As shown in Table~\ref{tab:wp-generators}, this approach generalizes effectively even to unseen models (e.g., Llama-3-8B, GPT-5-chat, GPT-neo-20B), matching or surpassing baselines like Lastde++ and Fast-DetectGPT.

% \begin{table}[htbp]
% \centering
% {\tiny
% \begin{tabular}{l c c c | c c c}
% \toprule
% WritingPrompts      & LLaMA2-13B & GPT-4.1-mini & GPT-J-6B & LLaMA3-8B & GPT-5-chat & GPT-neo-20B \\
% \midrule
% Lastde++            & 99.14      & 68.49        & \textbf{95.96} & 99.56      & 30.64      & 92.68      \\
% Fast-DetectGPT      & 99.56      & 70.23        & 93.80          & 99.84      & 30.01      & 92.22      \\
% SurpMark self-ref   & 99.59      & 87.27        & 96.85          & 99.87      & 83.56      & 93.93      \\
% SurpMark-MC         & \textbf{99.72} & \textbf{90.01} & 95.61      & \textbf{99.77} & \textbf{64.27} & \textbf{92.84} \\
% \bottomrule
% \end{tabular}}
% \caption{AUROC on WritingPrompts when the reference transition probabilities (LLaMA2-13B, GPT-4.1-mini and GPT-J-6B) are held fixed and not re-estimated for each test-time generator}
% \label{tab:wp-generators}
% \end{table}


\subsubsection{Decoding Strategies} \label{app:decoding}
In Table~\ref{tab:method-decoding}, to evaluate the effect of decoding stratigy, we use standard decoding strategies as described in Appendix~\ref{sec:exp-configuration}, varying one hyperparameter at a time while keeping the others at their default values. For open-source models on HuggingFace and Gemini, we (i) set top-p = 0.96, (ii) set top-k = 40, (iii) set temperature = 0.7. For GPT-5-chat, we vary one parameter at a time: top-p = 0.96 or temperature = 0.7 (no top-k parameter is exposed). Across all three models and decoding strategies, SurpMark either matches or exceeds the best baseline, and is especially strong under top-p/top-k sampling.
\begin{table}[htbp]
\centering
{\scriptsize
\begin{tabular}{l c c c c c c c c}
\toprule
\multirow{2}{*}{Method / Data@Model} 
  & \multicolumn{3}{c}{XSum@OPT-2.7B} 
  & \multicolumn{3}{c}{XSum@Gemma-7B} 
  & \multicolumn{2}{c}{WritingPrompts@GPT-5-chat} \\

    
  & top-p & top-k & temperature 
  & top-p & top-k & temperature 
  & top-p & temperature \\
\midrule
Likelihood     
  & 79.24 & 67.95 & 93.53
  & 66.73 & 55.56 & 87.56
  & 57.29 & 66.99 \\
LogRank        
  & 82.01 & 72.11 & 94.72
  & 68.28 & 59.25 & 88.69
  & 55.03 & 65.60 \\
Entropy        
  & 46.87 & 56.53 & 45.27
  & 49.16 & 52.12 & 45.06
  & 36.24 & 33.68 \\
LRR            
  & 83.43 & 77.88 & 93.48
  & 68.66 & 67.86 & 86.22
  & 47.28 & 59.38 \\
Lastde         
  & 86.09 & 81.26 & 94.19
  & 68.34 & 58.24 & 85.03
  & 39.38 & 47.74 \\
Lastde++       
  & 92.64 & 87.38 & 97.24
  & 81.43 & 69.42 & 93.15
  & 45.15 & 57.26 \\
Fast-DetectGPT 
  & 90.64 & 85.03 & \textbf{98.28}
  & 80.41 & 68.70 & \textbf{95.99}
  & 41.61 & 57.54 \\
SurpMark $k=6$ 
  & 92.41 & \textbf{87.81} & 96.65
  & \textbf{82.13} & 72.38 & 93.79
  & 75.80 & \textbf{77.08} \\
SurpMark $k=7$ 
  & \textbf{93.90} & 87.20 & 95.96
  & 80.90 & \textbf{77.88} & 93.57
  & \textbf{77.32} & \textbf{77.08} \\
\bottomrule
\end{tabular}}
\caption{AUROC of different detectors across decoding parameters, datasets, and models.}
\label{tab:method-decoding}
\end{table}


\subsubsection{Paraphrasing Attack} \label{sec:paraphrasing-attack}
Here we examine the robustness of detection methods to the paraphrasing attack. For SurpMark, we consider three paraphrase scenarios. Ref-P applies paraphrasing only to the offline references. Test-P paraphrases only the incoming text, which is the most realistic case in practice. Both-P paraphrases both sides. We follow the setup of Lastde++ and Fast-DetectGPT, and use T5-Paraphraser to perform paraphrasing attacks on texts. Under the practically most relevant Test-P case, the losses are minimal. Under Ref-P, the changes are modest. Under Both-P the drop is larger but still competitive. It shows that SurpMark’s surprisal-dynamics features are largely invariant to semantics-preserving rewrites.
\begin{table}[htbp]
{\scriptsize
    \centering
    \begin{tabular}{ccc|cc|cc}
    \toprule
    & \multicolumn{2}{c}{Xsum@Llama-3-8B}& \multicolumn{2}{c}{WritingPrompts@GPT-NeoX-20B}& \multicolumn{2}{c}{SQuAD@Llama-2-13B} \\
    \midrule
        & Original & Paraphrased & Original & Paraphrased& Original & Paraphrased \\
        \midrule
      Fast-DetectGPT  & 96.78& 95.3 ($\downarrow$1.48)&92.22&89.51 ($\downarrow$2.71)&94.85&92.78 ($\downarrow$2.07)\\
      Lastde++  & 93.42 &  91.3 ($\downarrow$2.12)&92.68&91.94 ($\downarrow$0.74)&97.32&92.12 ($\downarrow$5.2)\\
      SurpMark Ref-P   & 97.77 & 97.06 ($\downarrow$0.61)&94.31&93.12 ($\downarrow$1.19)&96.13& 94.89 ($\downarrow$1.24)\\
      SurpMark Test-P   &97.77 & 97.33 ($\downarrow$0.44)&94.31&94.05 ($\downarrow$0.26)&96.13&95.46 ($\downarrow$0.67)\\
      SurpMark Both-P   & 97.77 &97.17 ($\downarrow$0.6)&94.31&92.22 ($\downarrow$2.09)&96.13&93.98 ($\downarrow$2.15)\\
         \bottomrule
    \end{tabular}
    \caption{Robustness to paraphrase attacks. AUROC on three settings—XSum@Llama-3-8B, WritingPrompts@GPT-NeoX-20B, and SQuAD@Llama-2-13B. For SurpMark, Ref-P/Test-P/Both-P denote paraphrasing the reference set, the test text, or both.}}
    \label{tab:paraphrase-attack}
\end{table}


\subsubsection{Prompt-Engineered Adversarial Attacks} \label{sec:prompts-eng-attack}In this section, we run experiments with simple prompt-engineered attacks beyond plain paraphrasing. Specifically, for the XSum and WritingPrompts datasets, we design two types of attacks: (\textbf{attack 1}) prompts that ask the model to mimic human writing style, using instructions such as “Messy casual summary of the news article.” or “Short story in a quick, slightly messy human style.”; and (\textbf{attack 2}) prompts that explicitly instruct the model to evade detection, such as “Write a summary of the article that is designed to evade AI-text detectors.” or “Continue the story in a way that is hard for AI-text detectors.” See Table~\ref{tab:attack-gptj} for comparison. “SurpMark ref-attack” applies the adversarial prompts only when generating the reference machine texts, “SurpMark test-attack” applies them only to the test texts, and “SurpMark both-attack” applies the same adversarial prompts to both the reference and test texts. Across both datasets, SurpMark variants (especially the test-attack and both-attack settings) experience much smaller accuracy drops under all three attacks, showing the strongest overall robustness.

\begin{table}[h]
\centering
{\scriptsize
\begin{tabular}{l c c c c c c}
\toprule
& \multicolumn{3}{c}{WritingPrompts@GPT-J-6B} 
& \multicolumn{3}{c}{XSum@GPT-J-6B} \\


& Original & Attack 1 & Attack 2 
& Original & Attack 1 & Attack 2 \\
\midrule
Lastde++             
& 96.96 
& 84.24 ($\downarrow 12.72$) 
& 85.42 ($\downarrow 11.52$)
& 85.38 
& 69.79 ($\downarrow 15.59$) 
& 73.55 ($\downarrow 11.83$) \\
Fast-DetectGPT       
& 93.80 
& 85.95 ($\downarrow 7.85$) 
& 79.26 ($\downarrow 14.54$)
& 78.60 
& 75.44 ($\downarrow 3.16$) 
& 74.09 ($\downarrow 4.51$) \\
SurpMark ref-attack  
& 97.60 
& 95.06 ($\downarrow 2.54$) 
& 94.67 ($\downarrow 2.93$)
& 88.35 
& 83.84 ($\downarrow 4.51$) 
& 83.85 ($\downarrow 4.50$) \\
SurpMark test-attack 
& 97.60 
& 95.62 ($\downarrow 1.98$) 
& 92.59 ($\downarrow 5.01$)
& 88.35 
& 86.37 ($\downarrow 1.98$) 
& 84.86 ($\downarrow 3.49$) \\
SurpMark both-attack 
& 97.60 
& 94.30 ($\downarrow 3.30$) 
& 92.74 ($\downarrow 4.86$)
& 88.35 
& 84.44 ($\downarrow 3.91$) 
& 85.23 ($\downarrow 3.12$) \\
\bottomrule
\end{tabular}}
\caption{AUROC under adversarial attacks for different detectors on GPT-J-6B.}
\label{tab:attack-gptj}
\end{table}


\subsubsection{\quad Ablation on Necessity of First-Order Markov Chain} \label{sec:ablation-fo-markov}

In Table~\ref{tab:unigram-markov-comparison}, we evaluate the necessity of the use of first-order markov chain by comparing against the 1-gram distribution of surprisal states. Across the datasets, the first-order Markov features outperform the 1-gram distribution, with especially large gains on GPT-5-chat. This shows that modeling surprisal transitions, rather than only the stationary distribution, is particularly important for harder-to-detect models.
\begin{table}[htbp]
\centering
{\scriptsize
\begin{tabular}{l cccccc}
\toprule
                          & \multicolumn{3}{c}{GPT-J-6B}      & \multicolumn{3}{c}{GPT-5-chat} \\
Metric / Dataset          & XSum        & WritingPrompts & SQuAD    & XSum        & WritingPrompts & SQuAD    \\
\midrule
1-gram distribution       & 86.07       & 96.60          & 91.62    & 55.89       & 78.43          & 54.58    \\
First-order Markov chain  & 88.35       & 97.60          & 92.93    & 84.16       & 82.25          & 68.57    \\
\bottomrule
\end{tabular}}
\caption{AUROC of unigram vs. first-order Markov detectors across models and datasets.}
\label{tab:unigram-markov-comparison}
\end{table}


\subsubsection{\quad Ablation on Necessity of GJS distance} \label{sec:ablation-GJS-l1-l2}

In Table~\ref{tab:metric-model-dataset}, we evaluate the effect of different distance metrics including GJS divergence, $L_1$ and $L_2$ norm distance. GJS achieves the best AUROC on most dataset and source model. This suggests that GJS is a more robust measure than $L_1$ and $L_2$.

\begin{table}[htbp]
    \centering
    {\small
    \begin{tabular}{l c c c}
\toprule
                     & \multicolumn{3}{c}{GPT-4.1 mini} \\
          & XSum        & WritingPrompts & SQuAD      \\
\midrule
GJS                & 82.52       & 83.64          & 69.27      \\
$L_1$              & 73.51       & 82.17          & 62.28      \\
$L_2$              & 73.58       & 83.04          & 59.14      \\
\bottomrule
\end{tabular}
\caption{Comparison of different distance metrics across datasets.}
\label{tab:metric-model-dataset}}
\end{table}


\subsubsection{ \quad Analysis of Performance Disparity: Marginal vs. Transition Surprisal} \label{marginal-transition}
We investigate the performance disparity observed between closed-source (e.g., GPT-5-chat) and open-source models. Our analysis suggests that the distinguishing factor lies in the divergence between the generator and human text at the marginal surprisal level versus the transitional level.

For many open-source models, the marginal surprisal gap—the difference in the stationary distribution of token surprisals—is sufficiently large. Consequently, detectors relying on marginal statistics (e.g., Likelihood, LogRank, Entropy) perform well, and the relative gain from SurpMark is moderate. Conversely, for advanced closed-source models, this marginal gap is nearly negligible, rendering unigram-based methods ineffective. However, a significant transition gap persists in the surprisal dynamics. SurpMark captures these temporal dependencies, explaining its substantial performance advantage on proprietary models.

To quantify this, we compute the Jensen-Shannon (JS) divergence for both marginal surprisal distributions (JS-marginal) and first-order transition distributions (JS-transition) between human and machine text. As shown in Table~\ref{tab:marginal_vs_transition}, for GPT-5-chat, the ratio of transition divergence to marginal divergence is approximately 30, indicating that the signal primarily resides in the dynamics. In contrast, for GPT-J-6B, this ratio is close to 1, suggesting that marginal statistics alone are nearly as informative as transition statistics.

\begin{table}[htbp]
\centering


{\scriptsize
\begin{tabular}{llccc}
\toprule
\textbf{Generator} & \textbf{Dataset} & \textbf{JS-marginal} & \textbf{JS-transition} & \textbf{Ratio (Transition / Marginal)} \\
\midrule
GPT-J-6B     & XSum  & 0.00180   & 0.00228 & $\approx 1.27$ \\
   & SQuAD & 0.00358   & 0.00392 & $\approx 1.09$ \\
\midrule
GPT-5-chat   & XSum  & 0.00006   & 0.00170 & $\mathbf{\approx 29.97}$ \\
  & SQuAD & 0.00024   & 0.00100 & $\approx 4.17$ \\
\midrule
GPT-4.1-mini & XSum  & 0.00030   & 0.00160 & $\approx 5.33$ \\
& SQuAD & 0.00052   & 0.00150 & $\approx 2.88$ \\
\bottomrule
\end{tabular}}

\caption{Comparison of Jensen-Shannon (JS) divergence on marginal surprisal distributions versus first-order transition distributions. The high ratio for closed-source models (e.g., GPT-5-chat) indicates that detection signals are dominated by transition dynamics rather than marginal statistics.}
\label{tab:marginal_vs_transition}
\end{table}



\subsubsection{\quad Threshold Selection} \label{app:threshold}
The natural decision rule is simply the sign test by setting $\tau=0$. Our detector is built around the difference between two GJS divergences. Intuitively, $\Delta \text{GJS}$ is positive when the test sequence is closer to the machine reference than to the human reference, and negative in the opposite case. Also, $\Delta \text{GJS}$ can be viewed as a log-likelihood ratio $\Lambda_{n,N}$. In the classical Neyman-Pearson framework, the optimal likelihood-ratio test with equal class priors and symmetric costs is precisely $\Lambda_{n,N} \gtrless 0$. We additionally perform a threshold sensitivity study in Table~\ref{tab:tau-f1}. For each dataset and generator, we sweep $\tau$ over the full score range on the test set, compute precision/recall, and identify an optimal threshold $\tau^*$ that maximizes F1. We then compare F1 at our fixed choice $\tau=0$. Across all generators and datasets, F1 at $\tau=0$ is typically about 95-97\% of the oracle F1. This shows that in practice, our parameter-free sign-based rule already operates very close to the best threshold.

\begin{table}[htbp]
\centering
{\scriptsize
\begin{tabular}{l c c c c}
\toprule
Setting                     & AUROC & $\tau^\ast$              & F1@$\,\tau^\ast$ & F1@$\tau=0$ \\
\midrule
XSum@GPT-J-6B               & 89.12 & $2.92\times 10^{-5}$     & 83.56            & 80.36       \\
WritingPrompts@Llama-2-13B  & 99.75 & $-9.29\times 10^{-6}$    & 98.66            & 98.66       \\
SQuAD@Llama-3-8B            & 93.56 & $-4.49\times 10^{-5}$    & 87.58            & 82.69       \\
WritingPrompts@GPT-5-chat   & 80.63 & $-1.34\times 10^{-5}$    & 76.13            & 75.07       \\
\bottomrule
\end{tabular}}
\caption{AUROC and F1 scores at the optimal threshold $\tau^\ast$ and at $\tau=0$ across different settings.}
\label{tab:tau-f1}
\end{table}

In Lastde, the authors propose a fixed threshold of 2 for Lastde++ regardless of the source model, motivated by plotting score distributions and empirical performance across their experiments. In Table~\ref{tab:metric-lastde-surpmark}, we therefore compare F1 of two methods at their respective threshold. Across three of the four settings, SurpMark achieves higher AUROC, and in all four settings it attains a higher F1. On SQuAD@Llama-3-8B, Lastde++ has slightly higher AUROC, but at their fixed thresholds SurpMark still achieves higher F1, indicating SurpMark's sign-based decision rule is better calibrated and less sensitive to threshold choice.

\begin{table}[htbp]
\centering
{\tiny
\begin{tabular}{l l c c c c}
\toprule
Metric                         & Method        & XSum@GPT-J-6B & WritingPrompts@Llama-2-13B & SQuAD@Llama-3-8B & WritingPrompts@GPT-5-chat \\
\midrule
AUROC                          & Lastde++      & 85.38         & 99.14                       & 94.72            & 30.64                     \\
                               & SurpMark $k=6$& 88.35         & 99.47                       & 93.76            & 82.25                     \\
F1 at respective& Lastde++    & 63.44         & 95.56                       & 80.93            & 0.00                      \\
fixed threshold                               & SurpMark $k=6$& 80.36         & 98.66                       & 82.69            & 75.07                     \\
\bottomrule
\end{tabular}}
\caption{Comparison of AUROC and F1 at fixed thresholds for Lastde++ and SurpMark ($k=6$) across different settings.}
\label{tab:metric-lastde-surpmark}
\end{table}
%%%%%%%%%%%%%%%%%%%%%%%%%%%%%%%%%%%%%%%%%%%%%%%%%%%%%%%%%%%%%%%%%%%%%%%%%%%%%%%
%%%%%%%%%%%%%%%%%%%%%%%%%%%%%%%%%%%%%%%%%%%%%%%%%%%%%%%%%%%%%%%%%%%%%%%%%%%%%%%
\subsubsection{\quad Character-Word Perturbation OOD}
\label{app:cw-perturbation-ood}
We consider character- and word-level perturbations as a controlled form of surface-level distribution shift, including spelling noise, character edits, and word substitutions. In this setting, the reference sets remain clean and source-pure, while the test sets are constructed by applying character-level, word-level perturbations to both human-written and LLM-generated texts. This creates an OOD shift by modifying the surface form of the test texts through perturbations.

\begin{table}[htbp]
\centering
\scriptsize
\caption{Results on character-word perturbation on the DetectRL benchmark. The reference sets remain clean, while the test sets are constructed by applying character-level or word-level perturbations. The LLM-generated samples are drawn from a mixture of outputs from multiple LMs, including Llama-2-70B, Claude-Instant, and GPT-3.5-Turbo.}
\label{table:char_word_perturbation}
\begin{tabular}{llcc}
\toprule
Perturbation & Method & AUROC & TPR@FPR=5\% \\
\midrule
\multirow{3}{*}{Character-level}
& Fast-DetectGPT & 97.23 & 89.33 \\
& Lastde++ & 83.86 & 40.67 \\
& SurpMark & \textbf{99.80} & \textbf{99.33} \\
\midrule
\multirow{3}{*}{Word-level}
& Fast-DetectGPT & 98.82 & 93.33 \\
& Lastde++ & 90.75 & 68.00 \\
& SurpMark & \textbf{99.59} & \textbf{99.33} \\
\bottomrule
\end{tabular}
\end{table}

\subsubsection{Detection Using Larger Proxy LMs}
We further evaluate the effect of proxy model choice on GPT-5-chat generations. 
As shown in Table~\ref{tab:gpt5_proxy_models}, using larger proxy models generally improves the performance of the baselines, especially when moving from Deepseek-R1-Distill-Qwen-14B to Qwen-3-32B. 
However, SurpMark consistently outperforms Fast-DetectGPT and Lastde++ across all proxy models and both datasets. 
Notably, SurpMark remains strong even with smaller proxies. 
These results suggest that SurpMark is less sensitive to proxy model choice and can provide robust detection signals without relying on the largest available proxy model.

\begin{table*}[htbp]
\centering
\caption{AUROC on XSum and WritingPrompts generated by GPT-5-chat, using different proxy models: Qwen-3-14B, Deepseek-R1-Distill-Qwen-14B, and Qwen-3-32B.}
\label{tab:gpt5_proxy_models}
{\scriptsize
\begin{tabular}{lcc|cc|cc}
\toprule
\multirow{2}{*}{Method}
& \multicolumn{2}{c}{Qwen-3-14B}
& \multicolumn{2}{c}{Deepseek-R1-Distill-Qwen-14B}
& \multicolumn{2}{c}{Qwen-3-32B} \\
\cmidrule(lr){2-3}
\cmidrule(lr){4-5}
\cmidrule(lr){6-7}
& XSum@GPT-5-chat & WP@GPT-5-chat
& XSum@GPT-5-chat & WP@GPT-5-chat
& XSum@GPT-5-chat & WP@GPT-5-chat \\
\midrule
Fast-DetectGPT
& 77.88 & 52.68
& 49.10 & 35.40
& 84.91 & 69.90 \\

Lastde++
& 75.26 & 46.98
& 46.25 & 29.71
& 80.09 & 61.80 \\

SurpMark
& \textbf{81.12} & \textbf{90.06}
& \textbf{80.81} & \textbf{84.12}
& \textbf{87.10} & \textbf{92.36} \\
\bottomrule
\end{tabular}}
\end{table*}

\subsubsection{White-box Detection Results}
We further evaluate the effect of proxy-model access under white-box and black-box settings. 
As shown in Table~\ref{tab:xsum_white_black_box}, all methods achieve very high AUROC in the white-box setting, where the true source model is available. 
This is consistent with prior observations in Lastde++, which suggest that detection becomes easier when the detector has access to the source model. 
More importantly, SurpMark remains the strongest method in the more practical black-box setting, where GPT2-Large is used as the proxy model instead of the true generator.
\begin{table*}[htbp]
\centering
\caption{AUROC on XSum under white-box and black-box settings for four source models. The white-box setting assumes access to the true source model, while the black-box setting uses GPT2-Large as the proxy model.}
\label{tab:xsum_white_black_box}
{\scriptsize
\begin{tabular}{lcc|cc|cc|cc}
\toprule
\multirow{2}{*}{Method}
& \multicolumn{2}{c}{Llama-3-8B}
& \multicolumn{2}{c}{Llama-3.2-3B}
& \multicolumn{2}{c}{GPT-Neo-2.7B}
& \multicolumn{2}{c}{OPT-2.7B} \\
\cmidrule(lr){2-3}
\cmidrule(lr){4-5}
\cmidrule(lr){6-7}
\cmidrule(lr){8-9}
& White-box & Black-box
& White-box & Black-box
& White-box & Black-box
& White-box & Black-box \\
\midrule
Fast-DetectGPT
& 99.14 & 96.78
& 90.54 & 61.86
& 98.78 & 81.84
& 97.94 & 90.55 \\

Lastde++
& 98.64 & 95.90
& 88.11 & 59.90
& 99.57 & 87.50
& 98.75 & 87.93 \\

SurpMark
& \textbf{99.39} & \textbf{97.77}
& \textbf{92.15} & \textbf{73.07}
& 99.41 & \textbf{92.92}
& \textbf{99.21} & \textbf{91.16} \\
\bottomrule
\end{tabular}}
\end{table*}

\subsubsection{Results on Kaggle DAIGT V2 Train Dataset}
We also evaluate SurpMark on the Kaggle DAIGT V2 Train Dataset, which represents a challenging mixed-generator setting with texts generated by more than 15 models. 
As shown in Table~\ref{tab:daigt_v2}, SurpMark consistently outperforms the baselines across multiple runs, achieving the highest AUROC and TPR@FPR=5\%. 
These results further demonstrate the robustness of SurpMark under highly heterogeneous generator sources.
\begin{table}[htbp]
\centering
\caption{Results on the Kaggle DAIGT V2 Train Dataset, a challenging mixed-generator setting containing texts generated by more than 15 models. We report AUROC and TPR@FPR=5\% with standard deviation over multiple runs.}
\label{tab:daigt_v2}
{\scriptsize
\begin{tabular}{lcc}
\toprule
Method & AUROC & TPR@FPR=5\%  \\
\midrule
Lastde++ & 88.37 \com 2.45 & 78.00 \com 4.06 \\
Fast-DetectGPT & 88.75 \com 2.49 & 81.33 \com 4.16 \\
SurpMark & \textbf{96.17} \com 1.56 & \textbf{85.78} \com 1.39 \\
\bottomrule
\end{tabular}}
\end{table}

\subsubsection{\quad Impact of Pretraining Exposure on Human Texts} \label{app:pretraining_contamination}

We further examine whether our detection signal could be explained by human-written texts appearing in the pretraining data of the proxy LM. 
We conduct an experiment where the human-written texts are newer than the proxy model. 
Specifically, we add results on Dolly, whose human-written responses were released in 2023, and use GPT2-Large, released in 2019, as the proxy model. 
This provides a recent-human / older-proxy setting where pretraining exposure of the human responses is unlikely.

As shown in Table~\ref{tab:dolly_proxy_auroc}, the AUROC obtained with GPT2-Large is comparable to that obtained with Qwen-2-0.5B-Instruct, a newer proxy model released in 2024. 
Interestingly, GPT2-Large even achieves slightly higher AUROC, suggesting that the detection performance is unlikely to be merely an artifact of human references appearing in the proxy model’s pretraining data.

To further analyze this effect, we measure the smoothness of token surprisal sequences using the mean absolute first-order difference, where lower values indicate smoother token dynamics. 
As shown in Table~\ref{tab:surprisal_smoothness}, machine-generated texts remain substantially smoother than human-written texts, consistent with prior observations in Lastde. 
Moreover, Dolly human-written texts are slightly smoother when scored by Qwen-2-0.5B-Instruct than by GPT2-Large, which is consistent with prior findings that pretraining exposure can make human-written text exhibit smoother likelihood trajectories. 
In our setting, however, this effect makes human text more LM-like and therefore makes detection harder rather than easier.

Overall, these results suggest that pretraining exposure alone is unlikely to explain the effectiveness of SurpMark. 
If anything, such exposure may increase the machine-likeness of some human-written texts and make our reported detection performance a conservative estimate.

\begin{table}[htbp]
\centering
\caption{AUROC on Dolly under different generator/proxy model pairs.}
\label{tab:dolly_proxy_auroc}{\small
\begin{tabular}{lc}
\toprule
Generator@Proxy & Dolly \\
\midrule
GPT-J-6B@GPT2-Large & 92.96 \\
GPT-J-6B@Qwen-2-0.5B-Instruct & 90.54 \\
GPT-Neo-2.7B@GPT2-Large & 95.06 \\
GPT-Neo-2.7B@Qwen-2-0.5B-Instruct & 92.71 \\
\bottomrule
\end{tabular}}
\end{table}
\begin{table*}[b]
\centering
\caption{Smoothness of mean token surprisal sequences on Dolly, measured by $\frac{1}{T-1}\sum_{t=2}^{T}|s_t-s_{t-1}|$, where $s_t$ is the negative log probability of the $t$-th token. Lower values indicate smoother token dynamics. HWT proxied by Qwen-2-0.5B-Instruct is slightly smoother than HWT proxied by GPT2-Large, while MGT proxied by Qwen-2-0.5B-Instruct is much smoother overall.}
\label{tab:surprisal_smoothness}{\scriptsize
\begin{tabular}{llccc}
\toprule
Metric & Split 
& HWT (Dolly) proxied by GPT2-Large
& HWT (Dolly) proxied by Qwen-2-0.5B-Instruct
& MGT proxied by Qwen-2-0.5B-Instruct \\
\midrule
\multirow{2}{*}{Smoothness ($\downarrow$ smoother)}
& Ref  & 2.79 & 2.74 & 2.12 \\
& Test & 2.86 & 2.80 & 2.06 \\
\bottomrule
\end{tabular}}
\end{table*}
\end{document}



% This document was modified from the file originally made available by
% Pat Langley and Andrea Danyluk for ICML-2K. This version was created
% by Iain Murray in 2018, and modified by Alexandre Bouchard in
% 2019 and 2021 and by Csaba Szepesvari, Gang Niu and Sivan Sabato in 2022.
% Modified again in 2023 and 2024 by Sivan Sabato and Jonathan Scarlett.
% Previous contributors include Dan Roy, Lise Getoor and Tobias
% Scheffer, which was slightly modified from the 2010 version by
% Thorsten Joachims & Johannes Fuernkranz, slightly modified from the
% 2009 version by Kiri Wagstaff and Sam Roweis's 2008 version, which is
% slightly modified from Prasad Tadepalli's 2007 version which is a
% lightly changed version of the previous year's version by Andrew
% Moore, which was in turn edited from those of Kristian Kersting and
% Codrina Lauth. Alex Smola contributed to the algorithmic style files.
